% Generated mechanically from the frozen groupoids.tex target.
% Do not edit this generated file; edit the bound locale source instead.

% OK, start here.
%

\setcounter{chapter}{38}
\chapter{群胚概形}



\phantomsection
\label{groupoids-section-phantom}


\section{引言}
\label{groupoids-section-introduction}

\noindent
本章讨论群胚概形的一般理论。例如，可参阅 Keel 和 Mori
的优美论文 \cite{K-M}。





\section{记号}
\label{groupoids-section-notation}

\noindent
设 $S$ 为概形。若 $U$、$T$ 是 $S$ 上的概形，则以
$U(T)$ 表示 $U$ {\it 在} $S$ 上的 $T$-值点集合。换言之，
$U(T) = \Mor_S(T, U)$。在下文的讨论中，我们尽量保留字母 $T$
来表示 $S$ 上的一个“测试概形”。设 $X$、$Y$ 是
$S$ 上的概形，并设 $f : X \to Y$ 是 $S$ 上的概形态射。
对任意 $S$ 上的概形 $T$，都得到一个诱导的集合映射
$$
f : X(T) \longrightarrow Y(T)
$$
如上所示，我们也以 $f$ 表示此映射。事实上，该构造关于概形
$T/S$ 具有函子性。米田引理（见《范畴》引理 \ref{categories-lemma-yoneda}）说明，$f$ 决定函子变换
$f : h_X \to h_Y$，且也由这一变换唯一决定。更一般地，
对于纤维积之间的映射，我们采用相同的记号。例如，若
$X$、$Y$、$Z$ 是 $S$ 上的概形，且
$m : X \times_S Y \to Z \times_S Z$
是 $S$ 上的概形态射，则可将 $m$ 视为对应于一族 $T$-值点之间的映射
$$
X(T) \times Y(T) \longrightarrow Z(T) \times Z(T).
$$
其余情形依此类推。

\medskip\noindent
我们继续沿用从下标 $0$ 开始标记投影映射的约定；因此有
$\text{pr}_0 : X \times_S Y \to X$ 以及
$\text{pr}_1 : X \times_S Y \to Y$。






\section{等价关系}
\label{groupoids-section-equivalence-relations}

\noindent
回顾集合 $A$ 上的一个{\it 关系} $R$，它就是满足
$R \subset A \times A$ 的子集。通常用 $a R b$ 表示
$(a, b) \in R$。若
$a R b, b R c \Rightarrow a R c$，则称该关系具有{\it 传递性}；
若对所有 $a \in A$ 都有 $a R a$，则称其具有{\it 自反性}；
若 $a R b \Rightarrow b R a$，则称其具有{\it 对称性}。
同时具有传递性、自反性和对称性的关系称为{\it 等价关系}。

\medskip\noindent
在概形的情形中，我们略微放宽关系的概念，只要求
$R \to A \times A$ 是一个映射。定义如下。

\begin{definition}
\label{groupoids-definition-equivalence-relation}
设 $S$ 为概形，$U$ 为 $S$ 上的概形。
\begin{enumerate}
\item $U$ 在 $S$ 上的一个{\it 预关系}，是任意概形态射
$j : R \to U \times_S U$。在此情形下，令
$t = \text{pr}_0 \circ j$、$s = \text{pr}_1 \circ j$，从而
$j = (t, s)$。
\item $U$ 在 $S$ 上的一个{\it 关系}，是概形的单态射
$j : R \to U \times_S U$。
\item 一个{\it 预等价关系}是预关系
$j : R \to U \times_S U$，并且对所有 $T/S$，
$j : R(T) \to U(T) \times U(T)$ 的像都是等价关系。
\item 当且仅当对每个 $S$ 上的概形 $T$，$R$ 的 $T$-值点
在 $U$ 的 $T$-值点集合上给出一个等价关系时，我们称概形态射
$R \to U \times_S U$ 是
{\it $U$ 在 $S$ 上的一个等价关系}。
\end{enumerate}
\end{definition}

\noindent
换言之，等价关系就是一个预等价关系，并且 $j$ 是关系。

\begin{lemma}
\label{groupoids-lemma-restrict-relation}
设 $S$ 为概形，
$U$ 为 $S$ 上的概形，
$j : R \to U \times_S U$ 为预关系，
$g : U' \to U$ 为概形态射。最后令
$$
R' = (U' \times_S U')\times_{U \times_S U} R
\xrightarrow{j'}
U' \times_S U'
$$
则 $j'$ 是 $U'$ 在 $S$ 上的预关系。
若 $j$ 是关系，则 $j'$ 是关系；
若 $j$ 是预等价关系，则 $j'$ 是预等价关系；
若 $j$ 是等价关系，则 $j'$ 是等价关系。
\end{lemma}

\begin{proof}
证明略。
\end{proof}

\begin{definition}
\label{groupoids-definition-restrict-relation}
设 $S$ 为概形，$U$ 为 $S$ 上的概形，
$j : R \to U \times_S U$ 为预关系，
$g : U' \to U$ 为概形态射。
预关系 $j' : R' \to U' \times_S U'$ 称为预关系 $j$
在 $U'$ 上的{\it 限制}或{\it 拉回}。
在这种情形下，有时记 $R' = R|_{U'}$。
\end{definition}

\begin{lemma}
\label{groupoids-lemma-pre-equivalence-equivalence-relation-points}
设 $j : R \to U \times_S U$ 为预关系。
考虑由下述规则在概形 $U$ 的点上定义的关系：
$$
x \sim y
\Leftrightarrow
\exists\ r \in R :
t(r) = x,
s(r) = y.
$$
若 $j$ 是预等价关系，则这是一个等价关系。
\end{lemma}

\begin{proof}
假设 $x \sim y$ 且 $y \sim z$。
取 $r \in R$，使得 $t(r) = x$、$s(r) = y$；再取
$r' \in R$，使得 $t(r') = y$、$s(r') = z$。
取域 $K$，使其出现在如下交换图中：
$$
\xymatrix{
\kappa(r) \ar[r] & K \\
\kappa(y) \ar[u] \ar[r] & \kappa(r') \ar[u]
}
$$
用 $x_K, y_K, z_K : \Spec(K) \to U$
表示态射
$$
\begin{matrix}
\Spec(K) \to \Spec(\kappa(r))
\to
\Spec(\kappa(x)) \to U \\
\Spec(K) \to \Spec(\kappa(r))
\to
\Spec(\kappa(y)) \to U \\
\Spec(K) \to \Spec(\kappa(r'))
\to
\Spec(\kappa(z)) \to U
\end{matrix}
$$
由构造可知 $(x_K, y_K) \in j(R(K))$ 且
$(y_K, z_K) \in j(R(K))$。由于 $j$ 是预等价关系，
还可得到 $(x_K, z_K) \in j(R(K))$。
这显然蕴含 $x \sim z$。

\medskip\noindent
$\sim$ 的自反性和对称性的证明从略。
\end{proof}

\begin{lemma}
\label{groupoids-lemma-etale-equivalence-relation}
设 $j : R \to U \times_S U$ 为预关系。假设
\begin{enumerate}
\item $s, t$ 都是非分歧的；
\item 对 $S$ 上的任意代数闭域 $k$，映射
$R(k) \to U(k) \times U(k)$ 是等价关系；
\item 存在态射 $e : U \to R$、$i : R \to R$、
$c : R \times_{s, U, t} R \to R$，使得
$$
\xymatrix{
U \ar[r]_e \ar[d]_\Delta &
R \ar[d]_j &
R \ar[d]^j \ar[r]_i &
R \ar[d]^j &
R \times_{s, U, t} R \ar[d]^{j \times j} \ar[r]_c &
R \ar[d]^j \\
U \times_S U \ar[r] &
U \times_S U &
U \times_S U \ar[r]^{flip} &
U \times_S U &
U \times_S U \times_S U \ar[r]^{\text{pr}_{02}} &
U \times_S U
}
$$
均为交换图。
\end{enumerate}
则 $j$ 是等价关系。
\end{lemma}

\begin{proof}
由条件 (1) 以及
《概形态射》引理 \ref{morphisms-lemma-unramified-permanence}
可知 $j$ 是非分歧的。因此根据
《概形态射》引理 \ref{morphisms-lemma-diagonal-unramified-morphism}，
$\Delta_j : R \to R \times_{U \times_S U} R$ 是开浸入。
另一方面，条件 (2) 表明 $\Delta_j$ 在 $k$-值点上是双射，
因而 $\Delta_j$ 是同构，所以 $j$ 是单态射。再由上述交换图，
容易推出对 $S$ 上的所有概形 $T$，
$R(T) \subset U(T) \times U(T)$ 都是等价关系。
\end{proof}













\section{群概形}
\label{groupoids-section-group-schemes}

\noindent
回顾一个{\it 群}是二元组 $(G, m)$，其中 $G$ 是集合，
$m : G \times G \to G$ 是满足下列性质的集合映射：
\begin{enumerate}
\item （结合律）$m(g, m(g', g'')) = m(m(g, g'), g'')$
对所有 $g, g', g'' \in G$ 均成立；
\item （单位元）存在唯一的元素 $e \in G$
（称为 $G$ 的{\it 恒等元}、{\it 单位元}或 $1$），使得
$m(g, e) = m(e, g) = g$ 对所有 $g \in G$ 均成立；
\item （逆元）对所有 $g \in G$，存在 $i(g) \in G$，
使得 $m(g, i(g)) = m(i(g), g) = e$，其中 $e$ 是单位元。
\end{enumerate}
于是得到映射 $e : \{*\} \to G$ 和映射 $i : G \to G$，
并且四元组 $(G, m, e, i)$ 满足上列公理。

\medskip\noindent
一个{\it 群同态} $\psi : (G, m) \to (G', m')$ 是集合映射
$\psi : G \to G'$，且满足
$m'(\psi(g), \psi(g')) = \psi(m(g, g'))$。这自动保证
$\psi(e) = e'$ 以及 $i'(\psi(g)) = \psi(i(g))$。
（记号含义显然。）下文将使用这一点。

\begin{definition}
\label{groupoids-definition-group-scheme}
设 $S$ 为概形。
\begin{enumerate}
\item 一个{\it $S$ 上的群概形}是二元组 $(G, m)$，其中
$G$ 是 $S$ 上的概形，而 $m : G \times_S G \to G$
是 $S$ 上的概形态射，且满足如下性质：
对每个 $S$ 上的概形 $T$，二元组 $(G(T), m)$ 都是群。
\item 一个{\it $S$ 上群概形的态射
$\psi : (G, m) \to (G', m')$}
是 $S$ 上的概形态射 $\psi : G \to G'$，并且对每个 $T/S$，
诱导映射 $\psi : G(T) \to G'(T)$ 都是群同态。
\end{enumerate}
\end{definition}

\noindent
设 $(G, m)$ 是概形 $S$ 上的群概形。
由上述讨论（以及第 \ref{groupoids-section-notation} 节中的讨论），
可得到 $S$ 上的概形态射：
（单位元）$e : S \to G$ 以及（取逆）$i : G \to G$；
对每个 $T$，四元组 $(G(T), m, e, i)$ 都满足上列群公理。

\medskip\noindent
设 $(G, m)$、$(G', m')$ 是 $S$ 上的群概形，
$f : G \to G'$ 是 $S$ 上的概形态射。
由定义可知，$f$ 是 $S$ 上群概形的态射，当且仅当下图交换：
$$
\xymatrix{
G \times_S G \ar[r]_-{f \times f} \ar[d]_m &
G' \times_S G' \ar[d]^{m'} \\
G \ar[r]^f & G'
}
$$

\begin{lemma}
\label{groupoids-lemma-base-change-group-scheme}
设 $(G, m)$ 是 $S$ 上的群概形，
$S' \to S$ 是概形态射。
则拉回 $(G_{S'}, m_{S'})$ 是 $S'$ 上的群概形。
\end{lemma}

\begin{proof}
证明略。
\end{proof}

\begin{definition}
\label{groupoids-definition-closed-subgroup-scheme}
设 $S$ 为概形，$(G, m)$ 为 $S$ 上的群概形。
\begin{enumerate}
\item $G$ 的一个{\it 闭子群概形}是闭子概形
$H \subset G$，使得 $m|_{H \times_S H}$ 分解通过 $H$，
并在 $H$ 上诱导出 $S$ 上的群概形结构。
\item $G$ 的一个{\it 开子群概形}是开子概形
$G' \subset G$，使得 $m|_{G' \times_S G'}$ 分解通过 $G'$，
并在 $G'$ 上诱导出 $S$ 上的群概形结构。
\end{enumerate}
\end{definition}

\noindent
等价地，可以说 $H$ 是 $G$ 的闭子群概形，如果它是
$S$ 上的群概形，并带有一个 $S$ 上的群概形态射
$i : H \to G$，且该态射将 $H$ 识别为 $G$ 的闭子概形。

\begin{lemma}
\label{groupoids-lemma-closed-subgroup-scheme}
设 $S$ 为概形，$(G, m, e, i)$ 为 $S$ 上的群概形。
\begin{enumerate}
\item 闭子概形 $H \subset G$ 是闭子群概形，当且仅当
$e : S \to G$、$m|_{H \times_S H} : H \times_S H \to G$
及 $i|_H : H \to G$ 均分解通过 $H$。
\item 开子概形 $H \subset G$ 是开子群概形，当且仅当
$e : S \to G$、$m|_{H \times_S H} : H \times_S H \to G$
及 $i|_H : H \to G$ 均分解通过 $H$。
\end{enumerate}
\end{lemma}

\begin{proof}
考察 $T$-值点后，这恰好转化为刻画一个群的子集何时为子群的熟知条件。
\end{proof}

\begin{definition}
\label{groupoids-definition-smooth-group-scheme}
设 $S$ 为概形，$(G, m)$ 为 $S$ 上的群概形。
\begin{enumerate}
\item 若结构态射 $G \to S$ 是光滑的，则称 $G$ 为{\it 光滑群概形}。
\item 若结构态射 $G \to S$ 是平坦的，则称 $G$ 为{\it 平坦群概形}。
\item 若结构态射 $G \to S$ 是分离的，则称 $G$ 为{\it 分离群概形}。
\end{enumerate}
可按需增补。
\end{definition}






\section{群概形的例子}
\label{groupoids-section-examples-group-schemes}

\begin{example}[乘法群概形]
\label{groupoids-example-multiplicative-group}
考虑这样一个函子：它把任意概形 $T$ 映到结构层整体截面环的单位群
$\Gamma(T, \mathcal{O}_T^*)$。该函子由概形
$$
\mathbf{G}_m = \Spec(\mathbf{Z}[x, x^{-1}])
$$
表示。给出群结构的态射为
\begin{eqnarray*}
\mathbf{G}_m \times \mathbf{G}_m & \to & \mathbf{G}_m \\
\Spec(\mathbf{Z}[x, x^{-1}] \otimes_{\mathbf{Z}} \mathbf{Z}[x, x^{-1}])
& \to &
\Spec(\mathbf{Z}[x, x^{-1}]) \\
\mathbf{Z}[x, x^{-1}] \otimes_{\mathbf{Z}} \mathbf{Z}[x, x^{-1}]
& \leftarrow &
\mathbf{Z}[x, x^{-1}] \\
x \otimes x & \leftarrow & x
\end{eqnarray*}
由此可见，$\mathbf{G}_m$ 是 $\mathbf{Z}$ 上的群概形。
对任意概形 $S$，基变换 $\mathbf{G}_{m, S}$ 是 $S$ 上的群概形，
其点函子仍为
$$
T/S
\longmapsto
\mathbf{G}_{m, S}(T) = \mathbf{G}_m(T) = \Gamma(T, \mathcal{O}_T^*)
$$
。
\end{example}

\begin{example}[单位根]
\label{groupoids-example-roots-of-unity}
设 $n \in \mathbf{N}$。
考虑这样一个函子：它把任意概形 $T$ 映到
$\Gamma(T, \mathcal{O}_T^*)$ 中由 $n$ 次单位根组成的子群。
该函子由概形
$$
\mu_n = \Spec(\mathbf{Z}[x]/(x^n - 1)).
$$
表示。给出群结构的态射为
\begin{eqnarray*}
\mu_n \times \mu_n & \to & \mu_n \\
\Spec(
\mathbf{Z}[x]/(x^n - 1)
\otimes_{\mathbf{Z}}
\mathbf{Z}[x]/(x^n - 1))
& \to &
\Spec(\mathbf{Z}[x]/(x^n - 1)) \\
\mathbf{Z}[x]/(x^n - 1) \otimes_{\mathbf{Z}} \mathbf{Z}[x]/(x^n - 1)
& \leftarrow &
\mathbf{Z}[x]/(x^n - 1) \\
x \otimes x & \leftarrow & x
\end{eqnarray*}
由此可见，$\mu_n$ 是 $\mathbf{Z}$ 上的群概形。
对任意概形 $S$，基变换 $\mu_{n, S}$ 是 $S$ 上的群概形，
其点函子仍为
$$
T/S
\longmapsto
\mu_{n, S}(T) = \mu_n(T) = \{f \in \Gamma(T, \mathcal{O}_T^*) \mid f^n = 1\}
$$
。
\end{example}


\begin{example}[加法群概形]
\label{groupoids-example-additive-group}
考虑这样一个函子：它把任意概形 $T$ 映到结构层整体截面所成的群
$\Gamma(T, \mathcal{O}_T)$。该函子由概形
$$
\mathbf{G}_a = \Spec(\mathbf{Z}[x])
$$
表示。给出群结构的态射为
\begin{eqnarray*}
\mathbf{G}_a \times \mathbf{G}_a & \to & \mathbf{G}_a \\
\Spec(\mathbf{Z}[x] \otimes_{\mathbf{Z}} \mathbf{Z}[x])
& \to &
\Spec(\mathbf{Z}[x]) \\
\mathbf{Z}[x] \otimes_{\mathbf{Z}} \mathbf{Z}[x]
& \leftarrow &
\mathbf{Z}[x] \\
x \otimes 1 + 1 \otimes x & \leftarrow & x
\end{eqnarray*}
由此可见，$\mathbf{G}_a$ 是 $\mathbf{Z}$ 上的群概形。
对任意概形 $S$，基变换 $\mathbf{G}_{a, S}$ 是 $S$ 上的群概形，
其点函子仍为
$$
T/S
\longmapsto
\mathbf{G}_{a, S}(T) = \mathbf{G}_a(T) = \Gamma(T, \mathcal{O}_T)
$$
。
\end{example}

\begin{example}[一般线性群概形]
\label{groupoids-example-general-linear-group}
设 $n \geq 1$。
考虑这样一个函子：它把任意概形 $T$ 映到群
$$
\text{GL}_n(\Gamma(T, \mathcal{O}_T))
$$
即以结构层整体截面环为系数的可逆 $n \times n$ 矩阵所成的群。
该函子由概形
$$
\text{GL}_n = \Spec(\mathbf{Z}[\{x_{ij}\}_{1 \leq i, j \leq n}][1/d])
$$
表示，其中 $d = \det((x_{ij}))$，而 $(x_{ij})$ 是在
$(i, j)$ 位置取值 $x_{ij}$ 的 $n \times n$ 矩阵。
给出群结构的态射为
\begin{eqnarray*}
\text{GL}_n \times \text{GL}_n & \to & \text{GL}_n \\
\Spec(\mathbf{Z}[x_{ij}, 1/d] \otimes_{\mathbf{Z}}
\mathbf{Z}[x_{ij}, 1/d])
& \to &
\Spec(\mathbf{Z}[x_{ij}, 1/d]) \\
\mathbf{Z}[x_{ij}, 1/d] \otimes_{\mathbf{Z}} \mathbf{Z}[x_{ij}, 1/d]
& \leftarrow &
\mathbf{Z}[x_{ij}, 1/d] \\
\sum x_{ik} \otimes x_{kj} & \leftarrow & x_{ij}
\end{eqnarray*}
由此可见，$\text{GL}_n$ 是 $\mathbf{Z}$ 上的群概形。
对任意概形 $S$，基变换 $\text{GL}_{n, S}$ 是 $S$ 上的群概形，
其点函子仍为
$$
T/S
\longmapsto
\text{GL}_{n, S}(T) = \text{GL}_n(T) =\text{GL}_n(\Gamma(T, \mathcal{O}_T))
$$
。
\end{example}

\begin{example}
\label{groupoids-example-determinant}
行列式定义了 $\mathbf{Z}$ 上的群概形态射
$$
\det : \text{GL}_n \longrightarrow \mathbf{G}_m
$$
经基变换，对任意基概形 $S$，它给出其上的群概形态射
$\text{GL}_{n, S} \to \mathbf{G}_{m, S}$。
\end{example}

\begin{example}[常值群]
\label{groupoids-example-constant-group}
设 $G$ 为抽象群。考虑这样一个函子：它把任意概形 $T$
映到局部常值映射 $T \to G$ 所成的群（这里 $T$ 带扎里斯基拓扑，
$G$ 带离散拓扑）。该函子由概形
$$
G_{\Spec(\mathbf{Z})} =
\coprod\nolimits_{g \in G} \Spec(\mathbf{Z}).
$$
表示。给出群结构的态射为
$$
G_{\Spec(\mathbf{Z})}
\times_{\Spec(\mathbf{Z})}
G_{\Spec(\mathbf{Z})}
\longrightarrow
G_{\Spec(\mathbf{Z})}
$$
它把对应于二元组 $(g, g')$ 的分支映到对应于 $gg'$ 的分支。
对任意概形 $S$，基变换 $G_S$ 是 $S$ 上的群概形，
其点函子仍为
$$
T/S
\longmapsto
G_S(T) = \{f : T \to G \text{ 局部常值}\}
$$
。
\end{example}





\section{群概形的性质}
\label{groupoids-section-properties-group-schemes}

\noindent
本节汇集若干在任意基上都成立的群概形的简单性质。

\begin{lemma}
\label{groupoids-lemma-group-scheme-separated}
设 $S$ 为概形，$G$ 为 $S$ 上的群概形。
则 $G \to S$ 是分离的（相应地，\ 拟分离的），当且仅当单位元态射
$e : S \to G$ 是闭浸入（相应地，\ 拟紧的）。
\end{lemma}

\begin{proof}
回顾由《概形》引理 \ref{schemes-lemma-section-immersion}，
$e$ 是一个浸入；而当 $G \to S$ 分离（相应地，\ 拟分离）时，
它是闭浸入（相应地，\ 拟紧的）。
反过来，考虑下图：
$$
\xymatrix{
G \ar[r]_-{\Delta_{G/S}} \ar[d] &
G \times_S G \ar[d]^{(g, g') \mapsto m(i(g), g')} \\
S \ar[r]^e & G
}
$$
从代数几何的函子观点出发，验证此图为笛卡尔图是一个练习。
换言之，$\Delta_{G/S}$ 是 $e$ 的一个基变换。因此，若 $e$
是闭浸入（相应地，\ 拟紧的），则 $\Delta_{G/S}$ 也具有相同性质；
参见《概形》引理 \ref{schemes-lemma-base-change-immersion}
（相应地，\ 《概形》引理 \ref{schemes-lemma-quasi-compact-preserved-base-change}）。
\end{proof}

\begin{lemma}
\label{groupoids-lemma-flat-action-on-group-scheme}
设 $S$ 为概形，$G$ 为 $S$ 上的群概形。
设 $T$ 为 $S$ 上的概形，并设 $\psi : T \to G$ 是 $S$ 上的态射。
若 $T$ 在 $S$ 上平坦，则态射
$$
T \times_S G \longrightarrow G, \quad
(t, g) \longmapsto m(\psi(t), g)
$$
是平坦的。特别地，若 $G$ 在 $S$ 上平坦，则
$m : G \times_S G \to G$ 是平坦的。
\end{lemma}

\begin{proof}
考虑下图：
$$
\xymatrix{
T \times_S G \ar[rrr]_{(t, g) \mapsto (t, m(\psi(t), g))} & & &
T \times_S G \ar[r]_{\text{pr}} \ar[d] &
G \ar[d] \\
& & &
T \ar[r] &
S
}
$$
左上方的水平箭头是同构，方块为笛卡尔图。
因此结论由《概形态射》引理 \ref{morphisms-lemma-base-change-flat} 得出。
\end{proof}

\begin{lemma}
\label{groupoids-lemma-group-scheme-module-differentials}
\begin{reference}
\cite[Proposition 3.15]{BookAV}
\end{reference}
设 $(G, m, e, i)$ 为概形 $S$ 上的群概形，
$f : G \to S$ 为结构态射。则存在典范同构
$$
\Omega_{G/S} \cong f^*\mathcal{C}_{S/G} \cong f^*e^*\Omega_{G/S}
$$
其中 $\mathcal{C}_{S/G}$ 表示浸入 $e$ 的余法层。
特别地，若 $S$ 是一个域的谱，则 $\Omega_{G/S}$
是自由 $\mathcal{O}_G$-模。
\end{lemma}

\begin{proof}
由《概形态射》引理 \ref{morphisms-lemma-base-change-differentials} 有
$$
\Omega_{G \times_S G/G} = \text{pr}_0^*\Omega_{G/S}
$$
其中左端借助 $\text{pr}_1$ 将 $G \times_S G$ 视为 $G$ 上的概形。
令 $\tau : G \times_S G \to G \times_S G$ 为“剪切映射”，
它在点上由 $(g, h) \mapsto (m(g, h), h)$ 给出。
通过投影 $\text{pr}_1$ 将 $G \times_S G$ 视为 $G$ 上的概形时，
此映射是其自同构。结合以上两点，得到同构
$$
\tau^*\text{pr}_0^*\Omega_{G/S} \to \text{pr}_0^*\Omega_{G/S}
$$
由于 $\text{pr}_0 \circ \tau = m$，可将它改写为同构
$$
m^*\Omega_{G/S} \to \text{pr}_0^*\Omega_{G/S}
$$
再沿
$(e \circ f, \text{id}_G) : G \to G \times_S G$
拉回此同构，并利用 $m \circ (e \circ f, \text{id}_G) = \text{id}_G$
以及 $\text{pr}_0 \circ (e \circ f, \text{id}_G) = e \circ f$，
便得到所需的同构
$$
\Omega_{G/S} \to f^*e^*\Omega_{G/S}
$$
。由
《概形态射》引理 \ref{morphisms-lemma-differentials-relative-immersion-section}
还有 $\mathcal{C}_{S/G} \cong e^*\Omega_{G/S}$。
若 $S$ 是一个域的谱，则 $S$ 上的任意 $\mathcal{O}_S$-模都是自由的，
故最后一个断言成立。
\end{proof}

\begin{lemma}
\label{groupoids-lemma-group-scheme-addition-tangent-vectors}
设 $S$ 为概形，$G$ 为 $S$ 上的群概形，并设 $s \in S$。
则复合映射
$$
T_{G/S, e(s)} \oplus T_{G/S, e(s)} = T_{G \times_S G/S, (e(s), e(s))}
\rightarrow T_{G/S, e(s)}
$$
就是切向量的加法。这里的 $=$ 来自
《代数簇》引理 \ref{varieties-lemma-tangent-space-product}，
而右侧箭头由 $m : G \times_S G \to G$ 通过
《代数簇》引理 \ref{varieties-lemma-map-tangent-spaces} 诱导。
\end{lemma}

\begin{proof}
我们将使用 Varieties, Equation (\ref{varieties-equation-tangent-space-fibre})，
并在纤维中处理切向量。第一因子 $T_{G_s/s, e(s)}$ 中的元素
$\theta$，经由 $(1, e) : G_s \to G_s \times G_s$ 所诱导的映射
$T_{G_s/s, e(s)} \to T_{G_s \times G_s/s, (e(s), e(s))}$，
其像就是 $\theta$。由于 $m \circ (1, e) = 1$，
由函子性可知 $\theta$ 映到 $\theta$。
又因该映射是线性的，$(\theta_1, \theta_2)$ 映到
$\theta_1 + \theta_2$。
\end{proof}





\section{域上群概形的性质}
\label{groupoids-section-properties-group-schemes-field}

\noindent
本节汇集群概形在域上的若干性质。对于域上（局部）代数的群概形，
还可以得到更多结论；参见第 \ref{groupoids-section-algebraic-group-schemes} 节。

\begin{lemma}
\label{groupoids-lemma-group-scheme-over-field-open-multiplication}
若 $(G, m)$ 是域 $k$ 上的群概形，则乘法映射
$m : G \times_k G \to G$ 是开映射。
\end{lemma}

\begin{proof}
乘法映射同构于投影映射
$\text{pr}_0 : G \times_k G \to G$，因为交换图
$$
\xymatrix{
G \times_k G \ar[d]^m \ar[rrr]_{(g, g') \mapsto (m(g, g'), g')} & & &
G \times_k G \ar[d]^{(g, g') \mapsto g} \\
G \ar[rrr]^{\text{id}} & & & G
}
$$
的水平箭头都是同构。根据
《概形态射》引理 \ref{morphisms-lemma-scheme-over-field-universally-open}，
该投影是开映射。
\end{proof}

\begin{lemma}
\label{groupoids-lemma-group-scheme-over-field-translate-open}
设 $(G, m)$ 是域 $k$ 上的群概形，$U \subset G$ 是开集，
$T \to G$ 是概形态射。则复合映射
$T \times_k U \to G \times_k G \to G$ 的像是开集。
\end{lemma}

\begin{proof}
对任意域扩张 $K/k$，态射 $G_K \to G$ 是开映射
（《概形态射》引理 \ref{morphisms-lemma-scheme-over-field-universally-open}）。
$T \times_k U$ 的每个点 $\xi$ 都是某个 $K$ 下态射
$(t, u) : \Spec(K) \to T \times_k U$ 的像。此时
$T_K \times_K U_K = (T \times_k U)_K \to G_K$ 的像包含开平移
$t \cdot U_K$。综合这些事实可知，$T \times_k U \to G$ 的像
包含 $\xi$ 之像的一个开邻域。由于 $\xi$ 任意，结论成立。
\end{proof}

\begin{lemma}
\label{groupoids-lemma-group-scheme-over-field-separated}
设 $G$ 为域上的群概形。则 $G$ 是分离概形。
\end{lemma}

\begin{proof}
记 $S = \Spec(k)$，其中 $k$ 为域，并设 $G$ 为 $S$ 上的群概形。
由引理 \ref{groupoids-lemma-group-scheme-separated}，只需证明
$e : S \to G$ 是闭浸入。根据《概形态射》引理 \ref{morphisms-lemma-algebraic-residue-field-extension-closed-point-fibre}，
$e : S \to G$ 的像是 $G$ 的闭点。
显然，$\mathcal{O}_G \to e_*\mathcal{O}_S$ 是满射，因为
$e_*\mathcal{O}_S$ 是支集在 $G$ 的单位元处、取值为 $k$ 的摩天层。
由《概形》引理 \ref{schemes-lemma-characterize-closed-immersions}，
$e$ 是闭浸入。
\end{proof}

\begin{lemma}
\label{groupoids-lemma-group-scheme-field-geometrically-irreducible}
设 $G$ 是域 $k$ 上的群概形。则
\begin{enumerate}
\item $G$ 的每个局部环 $\mathcal{O}_{G, g}$ 都有唯一的极小素理想；
\item $G$ 中恰有一个经过 $e$ 的不可约分支 $Z$；且
\item $Z$ 在 $k$ 上几何不可约。
\end{enumerate}
\end{lemma}

\begin{proof}
对任意点 $g \in G$，存在域扩张 $K/k$ 以及映到 $g$ 的
$K$-值点 $g' \in G(K)$。把 $g'$ 看成群概形 $G_K$ 的
$K$-有理点，则
$\mathcal{O}_{G, g} \to \mathcal{O}_{G_K, g'}$
是忠实平坦的局部环同态（因为 $G_K \to G$ 平坦，且平坦的局部环同态
是忠实平坦的；见《交换代数》引理 \ref{algebra-lemma-local-flat-ff}）。
$\mathcal{O}_{G_K, g'}$ 的结论蕴含 $\mathcal{O}_{G, g}$ 的结论；
参见《交换代数》引理 \ref{algebra-lemma-injective-minimal-primes-in-image}。
因此，为证明 (1)，只需对 $G$ 的 $k$-有理点 $g$ 证明。
此时由 $g$ 作平移定义了自同构 $G \to G$，它把 $e$ 映到 $g$，
故 $\mathcal{O}_{G, g} \cong \mathcal{O}_{G, e}$。
这样便看出 (2) 蕴含 (1)，因为经过 $e$ 的不可约分支与
$\mathcal{O}_{G, e}$ 的极小素理想一一对应。

\medskip\noindent
为证明 (2) 和 (3)，只需在 $k$ 代数闭时证明 (2)。
在这种情形下，设 $Z_1$、$Z_2$ 是 $G$ 的两个经过 $e$ 的不可约分支。
由于 $k$ 代数闭，闭子概形
$Z_1 \times_k Z_2 \subset G \times_k G$ 也不可约；参见
《代数簇》引理 \ref{varieties-lemma-bijection-irreducible-components}。
因此 $m(Z_1 \times_k Z_2)$ 包含在 $G$ 的某个不可约分支中。
另一方面，由于 $m|_{e \times G} = \text{id}_G$ 且
$m|_{G \times e} = \text{id}_G$，它同时包含 $Z_1$ 和 $Z_2$。
故如所需，$Z_1 = Z_2$。
\end{proof}

\begin{remark}
\label{groupoids-remark-warning-group-scheme-geometrically-irreducible}
警告：引理 \ref{groupoids-lemma-group-scheme-field-geometrically-irreducible}
并不意味着 $G/k$ 的每个不可约分支都是几何不可约的。
例如，$\mathbf{Q}$ 上的群概形
$\mu_{3, \mathbf{Q}} = \Spec(\mathbf{Q}[x]/(x^3 - 1))$
有两个不可约分支，它们对应于分解
$x^3 - 1 = (x - 1)(x^2 + x + 1)$。
第一个因子对应经过单位元的不可约分支；第二个不可约分支在
$\Spec(\mathbf{Q})$ 上并非几何不可约。
\end{remark}

\begin{lemma}
\label{groupoids-lemma-reduced-subgroup-scheme-perfect}
设 $G$ 是完美域 $k$ 上的群概形。则 $G$ 的约化 $G_{red}$
是 $G$ 的闭子群概形。
\end{lemma}

\begin{proof}
证明略。提示：利用《代数簇》引理 \ref{varieties-lemma-perfect-reduced} 和
\ref{varieties-lemma-geometrically-reduced-any-base-change}
可知 $G_{red} \times_k G_{red}$ 是约化的。
\end{proof}

\begin{lemma}
\label{groupoids-lemma-open-subgroup-closed-over-field}
设 $k$ 为域，$\psi : G' \to G$ 是 $k$ 上的群概形态射。
若 $\psi(G')$ 在 $G$ 中是开的，则 $\psi(G')$ 在 $G$ 中是闭的。
\end{lemma}

\begin{proof}
令 $U = \psi(G') \subset G$，并令
$Z = G \setminus \psi(G') = G \setminus U$ 带有诱导的约化闭子概形结构。
由引理 \ref{groupoids-lemma-group-scheme-over-field-translate-open}，映射
$$
Z \times_k G' \longrightarrow
Z \times_k U \longrightarrow G
$$
的像是开的（第一个箭头是满射）。另一方面，由于 $\psi$
是群概形同态，$Z \times_k G' \to G$ 的像包含在 $Z$ 中
（因为对 $G'$ 的所有点 $g'$，由 $\psi(g')$ 作平移都保持
$U$；略去一个小细节）。因此 $Z \subset G$ 是开子集
（尽管未必是开子概形）。所以 $U = \psi(G')$ 是闭的。
\end{proof}

\begin{lemma}
\label{groupoids-lemma-immersion-group-schemes-closed-over-field}
设 $i : G' \to G$ 是域 $k$ 上群概形的浸入。
则 $i$ 是闭浸入；也就是说，$i(G')$ 是 $G$ 的闭子群概形。
\end{lemma}

\begin{proof}
为证明 $i$ 是闭浸入，只需证明 $i(G')$ 是 $G$ 的闭子集。
设 $k \subset k'$ 是 $k$ 的一个完美扩张。若
$i(G'_{k'}) \subset G_{k'}$ 是闭的，则由
《概形态射》引理 \ref{morphisms-lemma-fpqc-quotient-topology}，
$i(G') \subset G$ 也是闭的（因为 $G_{k'} \to G$ 平坦、拟紧且满射）。
因此可以并且确实假设 $k$ 完美。下文将直接使用这样一个事实：
$k$ 上约化概形的乘积仍是约化的。由
引理 \ref{groupoids-lemma-reduced-subgroup-scheme-perfect}，
可将 $G'$ 和 $G$ 换成各自的约化。
令 $\overline{G'} \subset G$ 为 $i(G')$ 的闭包，并赋予它约化闭子概形结构。
由《代数簇》引理 \ref{varieties-lemma-closure-of-product}，
$\overline{G'} \times_k \overline{G'}$ 是
$G' \times_k G' \to G \times_k G$ 之像的闭包。因此，由 $m$ 的连续性，
$$
m\Big(\overline{G'} \times_k \overline{G'}\Big)
\subset \overline{G'}
$$
成立。于是 $\overline{G'} \subset G$ 是一个约化闭子群概形。
由引理 \ref{groupoids-lemma-open-subgroup-closed-over-field}，
$i(G') \subset \overline{G'}$ 也是闭的，故如所需，
$i(G') = \overline{G'}$。
\end{proof}

\begin{lemma}
\label{groupoids-lemma-irreducible-group-scheme-over-field-qc}
设 $G$ 是域 $k$ 上的群概形。若 $G$ 不可约，则 $G$ 拟紧。
\end{lemma}

\begin{proof}
设 $K/k$ 为域扩张。若 $G_K$ 拟紧，则因 $G_K \to G$ 满射，
$G$ 也拟紧。由
引理 \ref{groupoids-lemma-group-scheme-field-geometrically-irreducible}，
$G_K$ 不可约。因此，只需把 $k$ 换成某个扩张后证明本引理。
取 $K$ 为基数非常大的代数闭域扩张。由
《代数簇》引理 \ref{varieties-lemma-make-Jacobson}，
$G_K$ 是 Jacobson 概形，并且其所有闭点的剩余域都等于 $K$。
换言之，可假设 $G$ 是 Jacobson 概形，且其所有闭点的剩余域都是 $k$。

\medskip\noindent
设 $U \subset G$ 是非空仿射开集，$g \in G(k)$。于是
$gU \cap U \not = \emptyset$，所以 $g$ 位于态射
$$
U \times_{\Spec(k)} U \longrightarrow G, \quad
(u_1, u_2) \longmapsto u_1u_2^{-1}
$$
的像中。根据引理 \ref{groupoids-lemma-group-scheme-over-field-open-multiplication}，该像是开的。
又因为 $G$ 是 Jacobson 概形且闭点都是 $k$-有理的，所以该像等于整个
$G$。由于 $U$ 仿射，$U \times_{\Spec(k)} U$ 也仿射。
因此 $G$ 是一个拟紧概形的像，从而拟紧。
\end{proof}

\begin{lemma}
\label{groupoids-lemma-connected-group-scheme-over-field-irreducible}
设 $G$ 是域 $k$ 上的群概形。若 $G$ 连通，则 $G$ 不可约。
\end{lemma}

\begin{proof}
由《代数簇》引理 \ref{varieties-lemma-geometrically-connected-if-connected-and-point}，
$G$ 几何连通。若对某个域扩张 $K/k$ 能证明 $G_K$ 不可约，
则本引理成立。因此可应用
《代数簇》引理 \ref{varieties-lemma-make-Jacobson}，
将问题约化到如下情形：$k$ 代数闭，$G$ 是 Jacobson 概形，
且所有闭点都是 $k$-有理的。

\medskip\noindent
设 $Z \subset G$ 是 $G$ 中唯一经过单位元的不可约分支；见
引理 \ref{groupoids-lemma-group-scheme-field-geometrically-irreducible}。
赋予 $Z$ 诱导的约化闭子概形结构，则 $Z \times_k Z$
约化且不可约（《代数簇》引理 \ref{varieties-lemma-geometrically-reduced-any-base-change} 和
\ref{varieties-lemma-bijection-irreducible-components}）。
因此 $m|_{Z \times_k Z} : Z \times_k Z \to G$ 分解通过 $Z$，
从而 $Z$ 成为 $G$ 的闭子群概形。

\medskip\noindent
为导出矛盾，假设还存在另一个不可约分支 $Z' \subset G$。
由引理 \ref{groupoids-lemma-group-scheme-field-geometrically-irreducible}，
$Z \cap Z' = \emptyset$。由
引理 \ref{groupoids-lemma-irreducible-group-scheme-over-field-qc}，$Z$ 拟紧。
于是可以选取拟紧开集 $U \subset G$，使得
$Z \subset U$ 且 $U \cap Z' = \emptyset$。
由引理 \ref{groupoids-lemma-group-scheme-over-field-translate-open}，
$Z \times_k U \to G$ 的像 $W$ 在 $G$ 中是开的。
另一方面，$W$ 是拟紧空间的像，因而拟紧。我们断言 $W$ 是闭的。

\medskip\noindent
证明该断言。由于 $W$ 拟紧而 $G$ 分离
（引理 \ref{groupoids-lemma-group-scheme-over-field-separated}），根据
《概形》引理 \ref{schemes-lemma-quasi-compact-permanence}，
包含映射 $W \to G$ 是拟紧的。因此，$W$ 闭包中的点都是
$W$ 中某些点的特化（《概形态射》引理 \ref{morphisms-lemma-reach-points-scheme-theoretic-image}）。
所以只需证明：$G$ 中任何与 $W$ 相交的不可约分支
$Z'' \subset G$ 都包含在 $W$ 中。
由于 $G$ 是 Jacobson 概形且闭点均为有理点，
$Z'' \cap W$ 有一个有理点 $g \in Z''(k) \cap W(k)$，
从而 $Z'' = Zg$。但由构造 $W = m(Z \times_k W)$，所以
$Z'' \cap W \not = \emptyset$ 蕴含 $Z'' \subset W$。

\medskip\noindent
由此断言，$W \subset G$ 是 $G$ 的既开又闭子集。
此外 $W \cap Z' = \emptyset$，否则由上一段的论证，
对某个 $g \in W(k)$ 会有 $Z' = Zg$。由于 $Z$ 是子群，
甚至可以取 $g \in U(k)$，这与 $Z' \cap U = \emptyset$ 矛盾。
因此 $W \subset G$ 是真既开又闭子集，与 $G$ 连通的假设矛盾。
\end{proof}

\begin{proposition}
\label{groupoids-proposition-connected-component}
设 $G$ 为域 $k$ 上的群概形。存在一个典范闭子群概形
$G^0 \subset G$，它具有下列性质：
\begin{enumerate}
\item $G^0 \to G$ 是平坦闭浸入；
\item $G^0 \subset G$ 是单位元的连通分支；
\item $G^0$ 几何不可约；且
\item $G^0$ 拟紧。
\end{enumerate}
\end{proposition}

\begin{proof}
令 $G^0$ 为单位元的连通分支，并赋予它典范概形结构
（《概形态射》定义 \ref{morphisms-definition-scheme-structure-connected-component}）。
为证明 $G^0$ 是闭子群概形，使用
引理 \ref{groupoids-lemma-closed-subgroup-scheme} 的判据。
态射 $e : \Spec(k) \to G$ 分解通过 $G^0$，因为我们所取的
$G^0$ 正是 $G$ 中包含 $e$ 的连通分支。
由于 $i : G \to G$ 是固定 $e$ 的自同构，$i$ 把 $G^0$ 映入自身。
根据《代数簇》引理 \ref{varieties-lemma-geometrically-connected-criterion}，
概形 $G^0$ 在 $k$ 上几何连通。因此 $G^0 \times_k G^0$ 连通
（《代数簇》引理 \ref{varieties-lemma-bijection-connected-components}）。
故在集合论意义下 $m(G^0 \times_k G^0) \subset G^0$。
由《概形态射》引理 \ref{morphisms-lemma-characterize-flat-closed-immersions}，
$m|_{G^0 \times_k G^0} : G^0 \times_k G^0 \to G$
分解通过 $G^0$。所以 $G^0$ 是 $G$ 的闭子群概形。
由引理 \ref{groupoids-lemma-connected-group-scheme-over-field-irreducible}，
$G^0$ 不可约；由
引理 \ref{groupoids-lemma-group-scheme-field-geometrically-irreducible}，
$G^0$ 几何不可约；再由
引理 \ref{groupoids-lemma-irreducible-group-scheme-over-field-qc}，
$G^0$ 拟紧。
\end{proof}

\begin{lemma}
\label{groupoids-lemma-profinite-product-over-field}
设 $k$ 为域，$T = \Spec(A)$，其中 $A$ 是一族代数的有向余极限，
而每个代数都是若干份 $k$ 的有限乘积。对任意 $k$ 上的概形 $X$，
作为拓扑空间有 $|T \times_k X| = |T| \times |X|$。
\end{lemma}

\begin{proof}
取仿射开覆盖后，可约化到 $X$ 仿射的情形。记
$X = \Spec(B)$。写成 $A = \colim A_i$，其中
$A_i = \prod_{t \in T_i} k$ 且 $T_i$ 有限。
于是 $T_i = |\Spec(A_i)|$ 带离散拓扑，而转移态射
$A_i \to A_{i'}$ 由集合映射 $T_{i'} \to T_i$ 给出。
因此，作为拓扑空间有 $|T| = \lim T_i$；见
《概形的极限》引理 \ref{limits-lemma-topology-limit}。类似地，
\begin{align*}
|T \times_k X| & =
|\Spec(A \otimes_k B)| \\
& =
|\Spec(\colim A_i \otimes_k B)| \\
& =
\lim |\Spec(A_i \otimes_k B)| \\
& =
\lim |\Spec(\prod\nolimits_{t \in T_i} B)| \\
& =
\lim T_i \times |X| \\
& =
(\lim T_i) \times |X| \\
& =
|T| \times |X|
\end{align*}
这里使用了上面的引理以及极限与极限可交换这一事实。
\end{proof}

\noindent
下一个引理说明，事实上可以把一个“点的代数性预有限族”
置于某个仿射开集中。强烈建议读者先阅读
引理 \ref{groupoids-lemma-points-in-affine}。

\begin{lemma}
\label{groupoids-lemma-compact-set-in-affine}
设 $k$ 为代数闭域，$G$ 为 $k$ 上的群概形。
假设 $G$ 是 Jacobson 概形，且所有闭点均为 $k$-有理点。
设 $T = \Spec(A)$，其中 $A$ 是一族代数的有向余极限，
而每个代数都是若干份 $k$ 的有限乘积。对任意态射 $f : T \to G$，
存在包含 $f(T)$ 的仿射开集 $U \subset G$。
\end{lemma}

\begin{proof}
令 $G^0 \subset G$ 为命题 \ref{groupoids-proposition-connected-component} 中得到的闭子群概形。
前两段将问题约化到 $G = G^0$ 的情形。

\medskip\noindent
注意，$T$ 是有限拓扑空间的有向逆极限
（《概形的极限》引理 \ref{limits-lemma-topology-limit}），
因而作为拓扑空间是预有限的
（《拓扑学》定义 \ref{topology-definition-profinite}）。
令 $W \subset G$ 为包含 $T \to G$ 之像的拟紧开集。
把 $W$ 替换为 $G^0 \times W \to G \times G \to G$ 的像后，
可假设 $W$ 在 $G^0$ 的左平移作用下不变；见
引理 \ref{groupoids-lemma-group-scheme-over-field-translate-open}。
考虑复合映射
$$
\psi = \pi \circ f : T \xrightarrow{f} W \xrightarrow{\pi} \pi_0(W)
$$
空间 $\pi_0(W)$ 是预有限的
（《拓扑学》引理 \ref{topology-lemma-spectral-pi0} 以及
《概形的性质》引理 \ref{properties-lemma-quasi-compact-quasi-separated-spectral}）。
对 $\xi \in \pi_0(W)$，令 $F_\xi \subset T$ 为
$T \to \pi_0(W)$ 在该点上的纤维。
假设对每个 $\xi \in \pi_0(W)$，都已找到仿射开集
$U_\xi \subset W$，使得 $f(F_\xi) \subset U_\xi$。
由《拓扑学》引理 \ref{topology-lemma-closed-map}，$\psi$ 是闭映射，
所以 $\psi(T \setminus f^{-1}(U_\xi))$ 是 $\pi_0(W)$
中不含 $\xi$ 的闭子集。其补集 $V_\xi$ 因而是 $\xi$ 的开邻域，
且 $\psi^{-1}(V_\xi) \subseteq f^{-1}(U_\xi)$。
由 $V_\xi$ 组成的 $\pi_0(W)$ 的开覆盖，可以细化为两两不交的
既开又闭子集所成的有限覆盖
$\pi_0(W) = V_1 \amalg \ldots \amalg V_n$
（《拓扑学》引理 \ref{topology-lemma-profinite-refine-open-covering}）。
对每个 $i = 1, \ldots, n$，取 $\xi_i \in \pi_0(W)$，
使得 $V_i \subset V_{\xi_i}$。由于 $V_i$ 闭而 $U_{\xi_i}$ 仿射，
$U_{\xi_i} \cap \pi^{-1}(V_i)$ 也仿射。
把 $U_{\xi_i}$ 替换为 $U_{\xi_i} \cap \pi^{-1}(V_i)$ 后，有
$\psi^{-1}(V_i) = f^{-1}(U_{\xi_i})$。因为这些 $V_i$ 覆盖
$\pi_0(W)$，所以这些 $f^{-1}(U_{\xi_i})$ 覆盖 $T$，从而
$f(T) \subseteq U = \bigcup_i U_{\xi_i}$。此时 $U$ 是有限多个
仿射开集的不交并，因此如所需为仿射的。

\medskip\noindent
设 $Z$ 为一个与 $f(T)$ 相交的 $G$ 的连通分支。则 $Z$
有一个 $k$-有理点 $z$（因为概形 $T$ 的所有剩余域都同构于 $k$）。
因此 $Z = G^0 z$。由 $W$ 的选取，$Z \subset W$。
上一段的论证把问题约化为：在 $W$ 中寻找 $f(T) \cap Z$
的仿射开邻域。经一个有理点平移后，可假设 $Z = G^0$
（细节略）。注意，概形论逆像
$T' = f^{-1}(G^0) \subset T$ 是同类型的闭子概形。
把 $T$ 替换为 $T'$ 后，可假设 $f(T) \subset G^0$。
选取 $e \in G$ 的一个仿射开邻域 $U \subset G$，
从而特别有 $U \cap G^0$ 非空。我们将证明，存在
$g \in G^0(k)$，使得 $f(T) \subset g^{-1}U$。
由于 $W$ 在 $G^0$ 的左平移下不变，$g^{-1}U \subset W$，
这将完成证明。

\medskip\noindent
前两段的论证允许我们转到 $G^0$，并将问题约化如下：
假设 $G$ 不可约，$U \subset G$ 是 $e$ 的仿射开邻域。
证明对某个 $g \in G(k)$ 有 $f(T) \subset g^{-1}U$。
考虑态射
$$
U \times_k T \longrightarrow G \times_k T,\quad
(t, u) \longrightarrow (uf(t)^{-1}, t)
$$
这是开浸入，因为它向 $G \times_k T \to G \times_k T$
的延拓是同构。由对 $T$ 的假设及
引理 \ref{groupoids-lemma-profinite-product-over-field}，有
$|U \times_k T| = |U| \times |T|$，对 $G \times_k T$ 也类似。
因此，上述开浸入的像是有限多个盒
$\bigcup_{i = 1, \ldots, n} U_i \times V_i$ 的并，其中
$V_i \subset T$，而 $U_i \subset G$ 是拟紧开集。
这意味着可能出现的开集 $Uf(t)^{-1}$（$t \in T$）只有有限多个，
记为 $Uf(t_1)^{-1}, \ldots, Uf(t_r)^{-1}$。
由于 $G$ 不可约，交集
$$
Uf(t_1)^{-1} \cap \ldots \cap Uf(t_r)^{-1}
$$
非空。又由于 $G$ 是 Jacobson 概形且闭点均为 $k$-有理点，
可在此交集中选取一个 $k$-值点 $g \in G(k)$。
于是对所有 $t \in T$ 都有 $g \in Uf(t)^{-1}$，
即如所需，$f(t) \in g^{-1}U$。
\end{proof}

\begin{remark}
\label{groupoids-remark-easy}
若 $G$ 是域上的群概形，是否总存在一个拟紧的既开又闭子群概形？
由命题 \ref{groupoids-proposition-connected-component}，
只有当 $G$ 在几何意义下有无穷多个连通分支时，这个问题才有意义。
\end{remark}

\begin{lemma}
\label{groupoids-lemma-group-scheme-field-countable-affine}
设 $G$ 为域上的群概形。存在既开又闭的子概形
$G' \subset G$，它是仿射概形的可数并。
\end{lemma}

\begin{proof}
设 $e \in U(k)$ 为单位元的拟紧开邻域。
把 $U$ 替换为 $U \cap i(U)$ 后，可假设 $U$ 在取逆映射下不变。
由于 $G$ 分离，这仍是拟紧集。令
$$
G' = \bigcup\nolimits_{n \geq 1} m_n(U \times_k \ldots \times_k U)
$$
其中 $m_n : G \times_k \ldots \times_k G \to G$ 是 $n$ 元乘法映射
$(g_1, \ldots, g_n) \mapsto m(m(\ldots (m(g_1, g_2), g_3), \ldots ), g_n)$。
这些映射都是开的（见
引理 \ref{groupoids-lemma-group-scheme-over-field-open-multiplication}），
故 $G'$ 是开子群概形。由
引理 \ref{groupoids-lemma-open-subgroup-closed-over-field}，它也是闭子群概形。
\end{proof}






\section{代数群概形的性质}
\label{groupoids-section-algebraic-group-schemes}

\noindent
回顾域 $k$ 上的概形称为（局部）代数的，是指它在
$\Spec(k)$ 上（局部）为有限型；见
《代数簇》定义 \ref{varieties-definition-algebraic-scheme}。
本节标题中的“代数”正是这个含义。

\begin{lemma}
\label{groupoids-lemma-group-scheme-finite-type-field}
设 $k$ 为域，$G$ 为 $k$ 上的局部代数群概形。
则 $G$ 等维，并且对所有 $g \in G$ 都有
$\dim(G) = \dim_g(G)$。对任意闭点 $g \in G$，有
$\dim(G) = \dim(\mathcal{O}_{G, g})$。
\end{lemma}

\begin{proof}
先证明，对任意一对点 $g, g' \in G$，都有
$\dim_g(G) = \dim_{g'}(G)$。由
《概形态射》引理 \ref{morphisms-lemma-dimension-fibre-after-base-change}，
可以任意扩张基域。因此可假设 $g$ 和 $g'$ 都在 $k$ 上有定义。
于是存在 $G$ 的一个自同构，把 $g$ 映到 $g'$，故二者相等。
根据《概形态射》引理 \ref{morphisms-lemma-dimension-fibre-at-a-point}，有
$\dim_g(G) = \dim(\mathcal{O}_{G, g}) +
\text{trdeg}_k(\kappa(g))$。
另一方面，$G$（或 $G$ 的任意开子集）的维数是 $G$ 的局部环维数的上确界；
见《概形的性质》引理 \ref{properties-lemma-codimension-local-ring}。
显然，该值在闭点 $g$ 处达到最大；此时由 Hilbert 零点定理，
$\text{trdeg}_k(\kappa(g)) = 0$（见
《概形态射》第 \ref{morphisms-section-points-finite-type} 节）。
引理得证。
\end{proof}

\noindent
下述结果有时称为 Cartier 定理。

\begin{lemma}
\label{groupoids-lemma-group-scheme-characteristic-zero-smooth}
设 $k$ 为特征 $0$ 的域，$G$ 为 $k$ 上的局部代数群概形。
则结构态射 $G \to \Spec(k)$ 是光滑的；也就是说，
$G$ 是光滑群概形。
\end{lemma}

\begin{proof}
由引理 \ref{groupoids-lemma-group-scheme-module-differentials}，
$G$ 在 $k$ 上的微分模是自由的。于是光滑性由
《代数簇》引理 \ref{varieties-lemma-char-zero-differentials-free-smooth}
推出。
\end{proof}

\begin{remark}
\label{groupoids-remark-when-reduced}
特征 $0$ 的域上的任意群概形都是约化的；参见
\cite[I, Theorem 1.1 and I, Corollary 3.9, and II, Theorem 2.4]{Perrin-thesis}
以及 \cite[Proposition 4.2.8]{Perrin}。
这一问题由 \cite[page 80]{Oort} 提出。
我们已在引理 \ref{groupoids-lemma-group-scheme-characteristic-zero-smooth}
中看到，当群概形局部为有限型时，该结论成立。
\end{remark}

\begin{lemma}
\label{groupoids-lemma-reduced-group-scheme-prefect-field-characteristic-p-smooth}
设 $k$ 为特征 $p > 0$ 的完美域（特征零情形见
引理 \ref{groupoids-lemma-group-scheme-characteristic-zero-smooth}）。
设 $G$ 为 $k$ 上的局部代数群概形。若 $G$ 约化，则结构态射
$G \to \Spec(k)$ 是光滑的；也就是说，$G$ 是光滑群概形。
\end{lemma}

\begin{proof}
由引理 \ref{groupoids-lemma-group-scheme-module-differentials}，
层 $\Omega_{G/k}$ 是自由的。因此结论由
《代数簇》引理 \ref{varieties-lemma-char-p-differentials-free-smooth}
推出。
\end{proof}

\begin{remark}
\label{groupoids-remark-reduced-smooth-not-true-general}
设 $k$ 为特征 $p > 0$ 的域，$\alpha \in k$ 不是 $p$ 次幂。
闭子群概形
$$
G = V(x^p + \alpha y^p) \subset \mathbf{G}_{a, k}^2
$$
约化且不可约，但不光滑，甚至不正规。
\end{remark}

\noindent
下一个引理是引理 \ref{groupoids-lemma-compact-set-in-affine} 的特殊情形，
其证明稍为简单。

\begin{lemma}
\label{groupoids-lemma-points-in-affine}
设 $k$ 为代数闭域，$G$ 为 $k$ 上的局部代数群概形，
$g_1, \ldots, g_n \in G(k)$ 为 $k$-有理点。
则存在包含 $g_1, \ldots, g_n$ 的仿射开集 $U \subset G$。
\end{lemma}

\begin{proof}
首先对 $n$ 归纳，说明可假设所有 $g_i$ 都位于 $G$ 的同一连通分支上。
否则，可以找到分解 $G = W_1 \amalg W_2$，其中 $W_i$ 在 $G$ 中开；
适当重编号后，对某个 $0 < r < n$ 有
$g_1, \ldots, g_r \in W_1$ 以及
$g_{r + 1}, \ldots, g_n \in W_2$。
由归纳假设，可找到 $G$ 的仿射开集 $U_1$ 和 $U_2$，使得
$g_1, \ldots, g_r \in U_1$ 且
$g_{r + 1}, \ldots, g_n \in U_2$。于是
$$
g_1, \ldots, g_n \in (U_1 \cap W_1) \cup (U_2 \cap W_2)
$$
给出所求解。因此可假设 $g_1, \ldots, g_n$ 都位于 $G$
的同一连通分支。由 $g_1^{-1}$ 作平移后，可假设
$g_1, \ldots, g_n \in G^0$，其中 $G^0 \subset G$ 如
命题 \ref{groupoids-proposition-connected-component} 所定义。
选取 $e$ 的仿射开邻域 $U$；特别地，$U \cap G^0$ 非空。
由于 $G^0$ 不可约，
$$
G^0 \cap (Ug_1^{-1} \cap \ldots \cap Ug_n^{-1})
$$
非空。由于 $G \to \Spec(k)$ 局部为有限型，
$G^0 \to \Spec(k)$ 也局部为有限型，故每个非空开集都有
$k$-有理点。因此可取 $g \in G^0(k)$，使得对所有 $i$ 都有
$g \in Ug_i^{-1}$。于是对所有 $i$ 都有 $g_i \in g^{-1}U$，
而 $g^{-1}U$ 就是所求仿射开集。
\end{proof}

\begin{lemma}
\label{groupoids-lemma-algebraic-quasi-projective}
设 $k$ 为域，$G$ 为 $k$ 上的代数群概形。则 $G$ 在 $k$ 上拟射影。
\end{lemma}

\begin{proof}
由《代数簇》引理 \ref{varieties-lemma-ample-after-field-extension}，
可假设 $k$ 代数闭。令 $G^0 \subset G$ 为命题 \ref{groupoids-proposition-connected-component} 中的 $G$ 的连通分支。
$G$ 的其他每个连通分支都有 $k$-有理点，因而作为概形同构于 $G^0$。
由于 $G$ 拟紧且 Noether，这样的连通分支只有有限多个。
于是问题约化到下一段所讨论的情形。

\medskip\noindent
设 $G$ 是代数闭域 $k$ 上的连通代数群概形。
若 $k$ 的特征为零，则由
引理 \ref{groupoids-lemma-group-scheme-characteristic-zero-smooth}，
$G$ 在 $k$ 上光滑。若 $k$ 的特征为 $p > 0$，令
$H = G_{red}$ 为 $G$ 的约化。根据
《除子》命题 \ref{divisors-proposition-push-down-ample}，
只需证明 $H$ 有充足可逆层。（对 $k$ 上的代数概形，
有充足可逆层等价于在 $k$ 上拟射影；例如见相当一般的
《态射进阶》引理 \ref{more-morphisms-lemma-quasi-projective}。）
由引理 \ref{groupoids-lemma-reduced-subgroup-scheme-perfect}，
$H$ 是 $k$ 上的群概形；由
引理 \ref{groupoids-lemma-reduced-group-scheme-prefect-field-characteristic-p-smooth}，
$H$ 在 $k$ 上光滑。这将问题约化到下一段所讨论的情形。

\medskip\noindent
设 $G$ 是代数闭域 $k$ 上拟紧、不可约且光滑的群概形。
注意，$G$ 的局部环都是正则的，因而是唯一分解整环
（《代数簇》引理 \ref{varieties-lemma-smooth-regular} 以及
《代数进阶》引理 \ref{more-algebra-lemma-regular-local-UFD}）。
$G$ 的一个非空仿射开集的补集，是某个有效 Cartier 除子 $D$ 的支集；
这由《除子》引理 \ref{divisors-lemma-complement-open-affine-effective-cartier-divisor} 得到。
（由引理 \ref{groupoids-lemma-group-scheme-over-field-separated}，$G$ 分离。）
因此存在有效 Cartier 除子 $D \subset G$，使得
$G \setminus D$ 仿射。下文还将用到：对任意 $n \geq 1$ 以及
$g_1, \ldots, g_n \in G(k)$，补集
$G \setminus \bigcup D g_i$ 仿射。事实上，它是分离概形 $G$ 中
诸仿射开集 $G \setminus Dg_i \cong G \setminus D$ 的交。

\medskip\noindent
可以选取下图的第一行
$$
\xymatrix{
G & U \ar[l]_j \ar[r]^\pi & \mathbf{A}^d_k \\
& W \ar[r]^{\pi'} \ar[u] & V \ar[u]
}
$$
使得 $U \not = \emptyset$，$j : U \to G$ 是开浸入，而
$\pi$ 是平展的（\'etale）；参见
《概形态射》引理 \ref{morphisms-lemma-smooth-etale-over-affine-space}。
存在非空仿射开集 $V \subset \mathbf{A}^d_k$，使得取
$W = \pi^{-1}(V)$ 时，态射 $\pi' = \pi|_W : W \to V$
是有限平展的（finite \'etale）。特别地，$\pi'$ 有限局部自由，
设其次数为 $n$。考虑有效 Cartier 除子
$$
\mathcal{D} = \{(g, w) \mid m(g, j(w)) \in D\} \subset G \times W
$$
（这是把 $D \subset G$ 沿平坦态射 $m : G \times G \to G$
拉回后，再限制到 $G \times W$ 所得的除子。）
考虑闭子集\footnote{利用《除子》第 \ref{divisors-section-norms} 节的材料，可以把有效 Cartier 除子
$\mathcal{D}$ 沿有限局部自由态射 $1 \times \pi'$ 的范数取作
有效 Cartier 除子 $E$，从而省去部分论证。}
$T = (1 \times \pi')(\mathcal{D}) \subset G \times V$。
由于 $\pi'$ 有限局部自由，$T$ 的每个不可约分支在
$G \times V$ 中的余维都是 $1$。又因 $G \times V$ 在 $k$ 上光滑，
这些分支都是有效 Cartier 除子（《除子》引理 \ref{divisors-lemma-weil-divisor-is-cartier-UFD} 以及上引诸引理），
所以 $T$ 是 $G \times V$ 中某个有效 Cartier 除子 $E$ 的支集。
若 $v \in V(k)$，则
$(\pi')^{-1}(v) = \{w_1, \ldots, w_n\} \subset W(k)$，并且在集合论意义下，
$$
E_v = \bigcup\nolimits_{i = 1, \ldots, n} D j(w_i)^{-1}
$$
在 $G$ 中成立。特别地，$G \setminus E_v$ 是仿射开集（见上文）。
此外，对任意 $g \in G(k)$，存在 $v \in V$，使得
$g \not \in E_v$。事实上，满足 $g \not \in Dj(w)^{-1}$ 的
$w \in W$ 所成的集合 $W'$ 是非空开集；只需选取 $v$，
使得 $W' \to V$ 在 $v$ 上的纤维有 $n$ 个元素。

\medskip\noindent
考虑 $G \times V$ 上的可逆层
$\mathcal{M} = \mathcal{O}_{G \times V}(E)$。
由《代数簇》引理 \ref{varieties-lemma-rational-equivalence-for-Pic}，
限制 $\mathcal{M}_v = \mathcal{O}_G(E_v)$ 的同构类
$\mathcal{L}$ 与 $v \in V(k)$ 无关。另一方面，对每个
$g \in G(k)$，都可找到 $v$，使得 $g \not \in E_v$ 且
$G \setminus E_v$ 仿射。因此，$\mathcal{O}_G(E_v)$ 的典范截面
（《除子》定义 \ref{divisors-definition-invertible-sheaf-effective-Cartier-divisor}）
对应于 $\mathcal{L}$ 的一个截面 $s_v$；它在 $g$ 处不消失，
并且 $G_{s_v}$ 仿射。按照定义，这说明 $\mathcal{L}$ 充足
（《概形的性质》定义 \ref{properties-definition-ample}）。
\end{proof}

\begin{lemma}
\label{groupoids-lemma-algebraic-center}
设 $k$ 为域，$G$ 为 $k$ 上的局部代数群概形。
则 $G$ 的中心是 $G$ 的闭子群概形。
\end{lemma}

\begin{proof}
以 $\text{Aut}(G)$ 表示 $k$ 上概形范畴中的反变函子：
它把 $S/k$ 映到基变换 $G_S$ 作为 $S$ 上群概形的自同构集合。
有自然变换
$$
G \longrightarrow \text{Aut}(G),\quad
g \longmapsto \text{inn}_g
$$
它把 $G$ 的一个 $S$-值点 $g$ 映到由 $g$ 决定的 $G$ 的内自同构。
按照定义，$G$ 的中心 $C$ 是该变换的核；也就是说，它是把
$S$ 映到如下 $g \in G(S)$ 所成集合的函子：其对应的内自同构
是恒等的。本引理的断言是，该函子可由 $G$ 的一个闭子群概形表示。

\medskip\noindent
选取整数 $n \geq 1$。令 $G_n \subset G$ 为 $G$ 的单位元 $e$
的第 $n$ 阶无穷小邻域。对每个概形 $S/k$，基变换 $G_{n, S}$
都是 $e_S : S \to G_S$ 的第 $n$ 阶无穷小邻域。
因此有自然变换
$\text{Aut}(G) \to \text{Aut}(G_n)$，右端是 $G_n$
作为概形的自同构函子（一般而言，$G_n$ 不是群概形）。
注意，$G_n$ 是某个 Artin 局部环 $A_n$ 的谱；其剩余域为 $k$，
并且作为 $k$-向量空间是有限维的
（《代数簇》引理 \ref{varieties-lemma-algebraic-scheme-dim-0}）。
由于 $G_n$ 的每个自同构都特别诱导一个可逆线性映射
$A_n \to A_n$，得到函子变换
$$
G \to \text{Aut}(G) \to \text{Aut}(G_n) \to \text{GL}(A_n)
$$
最后一个群值函子是可表示的；见
例 \ref{groupoids-example-general-linear-group}。并且最后一个箭头显然为单射。
因此，对每个 $n$ 都得到闭子群概形
$$
H_n = \Ker(G \to \text{Aut}(G_n)) = \Ker(G \to \text{GL}(A_n)).
$$
作为第一步近似，令 $H = \bigcap_{n \geq 1} H_n$
（概形论交）。这是一个包含中心 $C$ 的闭子群概形。

\medskip\noindent
设 $h$ 是 $H$ 的一个 $S$-值点，其中 $S$ 局部 Noether。
于是自同构 $\text{inn}_h$ 在所有闭子概形 $G_{n, S}$ 上诱导恒等映射。
考虑核
$K = \Ker(\text{inn}_h : G_S \to G_S)$。
它是 $S$ 上 $G_S$ 的闭子群概形，并包含所有 $n \geq 1$ 的
闭子概形 $G_{n, S}$。这蕴含 $K$ 包含
$e(S) \subset G_S$ 的一个开邻域；见
《交换代数》注 \ref{algebra-remark-intersection-powers-ideal}。
令 $G^0 \subset G$ 如命题 \ref{groupoids-proposition-connected-component} 所定义。
由于 $G^0$ 几何不可约，可断定 $K$ 包含 $G^0_S$
（对任意非空开集 $U \subset G^0_{k'}$ 和任意域扩张 $k'/k$，
都有 $U \cdot U^{-1} = G^0_{k'}$；见
引理 \ref{groupoids-lemma-irreducible-group-scheme-over-field-qc} 的证明）。
取 $S = H$ 应用上述结论，得到 $G^0$ 和 $H$ 是 $G$
中其点两两交换的子群概形：对任意概形 $S$ 以及任意 $S$-值点
$g \in G^0(S)$、$h \in H(S)$，在 $G(S)$ 中都有 $gh = hg$。

\medskip\noindent
假设 $k$ 代数闭。在 $G$ 的每个不可约分支 $G_i$ 中选取一个
$k$-值点 $g_i$。注意，在此情形下，$G$ 的连通分支就是 $G$ 的不可约分支，
而它们都是 $G^0$ 经各 $g_i$ 平移所得。我们断言
$$
C = H \cap \bigcap\nolimits_i \Ker(\text{inn}_{g_i} : G \to G)
\quad (\text{概形论交})
$$
事实上，$C$ 包含于右端。反过来，右端的每个 $S$-值点 $h$
都与 $G^0$ 以及各 $g_i$ 交换，因而与
$G = \bigcup G^0g_i$ 中的一切元素交换。

\medskip\noindent
一般基域 $k$ 的情形由代数闭包 $\overline{k}$ 上的结果经下降得到。
具体地，令 $A \subset G_{\overline{k}}$ 为表示
$G_{\overline{k}}$ 之中心的闭子群概形。由中心构造的函子性，
作为 $G_{\overline{k} \otimes_k \overline{k}}$ 的闭子概形有
$$
A \times_{\Spec(k)} \Spec(\overline{k}) =
\Spec(\overline{k}) \times_{\Spec(k)} A
$$
所以由《下降》引理 \ref{descent-lemma-closed-immersion}
（以及《下降》引理 \ref{descent-lemma-descending-property-closed-immersion}），
$A$ 下降为闭子群概形 $Z \subset G$。
于是 $Z$ 表示 $C$（略去一个小论证），证明完毕。
\end{proof}







\section{阿贝尔簇}
\label{groupoids-section-abelian-varieties}

\noindent
Mumford 关于阿贝尔簇的著作是这一主题的极佳参考文献；
见 \cite{AVar}，建议读者参阅。存在许多等价定义，下面给出其中一个。

\begin{definition}
\label{groupoids-definition-abelian-variety}
设 $k$ 为域。一个{\it 阿贝尔簇}是 $k$ 上的群概形，同时也是
$k$ 上固有且几何整的簇\footnote{等价定义见
注 \ref{groupoids-remark-16-equivalent}。}。
\end{definition}

\noindent
下面先证明关于这一概念的几个引理，再在
命题 \ref{groupoids-proposition-review-abelian-varieties}
中汇总所有结论。

\begin{lemma}
\label{groupoids-lemma-abelian-variety-projective}
设 $k$ 为域，$A$ 为 $k$ 上的阿贝尔簇。则 $A$ 是射影的。
\end{lemma}

\begin{proof}
这由引理 \ref{groupoids-lemma-algebraic-quasi-projective} 以及
《态射进阶》引理 \ref{more-morphisms-lemma-projective} 得出。
\end{proof}

\begin{lemma}
\label{groupoids-lemma-abelian-variety-change-field}
设 $k$ 为域，$A$ 为 $k$ 上的阿贝尔簇。
对任意域扩张 $K/k$，基变换 $A_K$ 是 $K$ 上的阿贝尔簇。
\end{lemma}

\begin{proof}
证明略。注意，这正是我们坚持要求 $A$ 几何整的原因；
若无此条件，本引理以及下文许多结论都会错误。
\end{proof}

\begin{lemma}
\label{groupoids-lemma-abelian-variety-smooth}
设 $k$ 为域，$A$ 为 $k$ 上的阿贝尔簇。则 $A$ 在 $k$ 上光滑。
\end{lemma}

\begin{proof}
若 $k$ 完美，则该结论由
引理 \ref{groupoids-lemma-group-scheme-characteristic-zero-smooth}（特征零）
以及引理 \ref{groupoids-lemma-reduced-group-scheme-prefect-field-characteristic-p-smooth}
（正特征）得出。一般情形可利用光滑性的下降
（《下降》引理 \ref{descent-lemma-descending-property-smooth}）
以及引理 \ref{groupoids-lemma-abelian-variety-change-field} 转到完美闭包，
从而约化到上述情形。
\end{proof}

\begin{lemma}
\label{groupoids-lemma-abelian-variety-abelian}
阿贝尔簇是交换群概形；也就是说，群律可交换。
\end{lemma}

\begin{proof}
设 $k$ 为域，$A$ 为 $k$ 上的阿贝尔簇。
由引理 \ref{groupoids-lemma-abelian-variety-change-field}，可把 $k$
替换为其代数闭包。考虑态射
$$
h : A \times_k A \longrightarrow A \times_k A,\quad
(x, y) \longmapsto (x, xyx^{-1}y^{-1})
$$
两边都取第一投影时，这是 $A$ 上的态射。设 $e \in A(k)$ 为单位元。
于是 $h|_{e \times A}$ 是取常值 $(e, e)$ 的映射。根据
《态射进阶》引理 \ref{more-morphisms-lemma-flat-proper-family-cannot-collapse-fibre}，
存在 $e$ 的开邻域 $U \subset A$，使得 $h|_{U \times A}$
分解通过某个在 $U$ 上有限的 $Z \subset U \times A$。
这意味着，对 $x \in U(k)$，态射
$A \to A$，$y \mapsto xyx^{-1}y^{-1}$ 只取有限多个值。
当然，这意味着它恒取值 $e$。所以
$(x, y) \mapsto xyx^{-1}y^{-1}$ 在 $U \times A$ 上恒取值 $e$，
从而 $A$ 上的群律可交换。
\end{proof}

\begin{lemma}
\label{groupoids-lemma-apply-cube}
设 $k$ 为域，$A$ 为 $k$ 上的阿贝尔簇，并设 $\mathcal{L}$
为可逆 $\mathcal{O}_A$-模。则在 $A \times_k A \times_k A$ 上
存在可逆模的同构
$$
m_{1, 2, 3}^*\mathcal{L} \otimes
m_1^*\mathcal{L} \otimes
m_2^*\mathcal{L} \otimes
m_3^*\mathcal{L} \cong
m_{1, 2}^*\mathcal{L} \otimes
m_{1, 3}^*\mathcal{L} \otimes
m_{2, 3}^*\mathcal{L}
$$
其中 $m_{i_1, \ldots, i_t} : A \times_k A \times_k A \to A$
是态射 $(x_1, x_2, x_3) \mapsto \sum x_{i_j}$。
\end{lemma}

\begin{proof}
把立方定理（《态射进阶》定理 \ref{more-morphisms-theorem-of-the-cube}）应用于差
$$
\mathcal{M} =
m_{1, 2, 3}^*\mathcal{L} \otimes
m_1^*\mathcal{L} \otimes
m_2^*\mathcal{L} \otimes
m_3^*\mathcal{L} \otimes
m_{1, 2}^*\mathcal{L}^{\otimes -1} \otimes
m_{1, 3}^*\mathcal{L}^{\otimes -1} \otimes
m_{2, 3}^*\mathcal{L}^{\otimes -1}
$$
即可。其理由是，$\mathcal{M}$ 在
$A \times A \times e = A \times A$ 上的限制等于
$$
n_{1, 2}^*\mathcal{L} \otimes
n_1^*\mathcal{L} \otimes
n_2^*\mathcal{L} \otimes
n_{1, 2}^*\mathcal{L}^{\otimes -1} \otimes
n_1^*\mathcal{L}^{\otimes -1} \otimes
n_2^*\mathcal{L}^{\otimes -1} \cong \mathcal{O}_{A \times_k A}
$$
其中 $n_{i_1, \ldots, i_t} : A \times_k A \to A$
为态射 $(x_1, x_2) \mapsto \sum x_{i_j}$。
对 $A \times e \times A$ 和 $e \times A \times A$ 亦然。
\end{proof}

\begin{lemma}
\label{groupoids-lemma-pullbacks-by-n}
设 $k$ 为域，$A$ 为 $k$ 上的阿贝尔簇，并设 $\mathcal{L}$
为可逆 $\mathcal{O}_A$-模。则
$$
[n]^*\mathcal{L} \cong
\mathcal{L}^{\otimes n(n + 1)/2} \otimes
([-1]^*\mathcal{L})^{\otimes n(n - 1)/2}
$$
其中 $[n] : A \to A$ 把 $x$ 映到有 $n$ 个加项的
$x + x + \ldots + x$，而 $[-1] : A \to A$ 是 $A$ 的取逆映射。
\end{lemma}

\begin{proof}
考虑态射 $A \to A \times_k A \times_k A$，
$x \mapsto (x, x, -x)$，其中 $-x = [-1](x)$。
拉回引理 \ref{groupoids-lemma-apply-cube} 中的关系，得到
$$
\mathcal{L} \otimes
\mathcal{L} \otimes
\mathcal{L} \otimes
[-1]^*\mathcal{L} \cong
[2]^*\mathcal{L}
$$
这证明了 $n = 2$ 的情形。归纳假设结论对
$1, 2, \ldots, n$ 成立。再考虑态射
$A \to A \times_k A \times_k A$，$x \mapsto (x, x, [n - 1]x)$。
拉回引理 \ref{groupoids-lemma-apply-cube} 中的关系，得到
$$
[n + 1]^*\mathcal{L} \otimes
\mathcal{L} \otimes
\mathcal{L} \otimes
[n - 1]^*\mathcal{L} \cong
[2]^*\mathcal{L} \otimes
[n]^*\mathcal{L} \otimes
[n]^*\mathcal{L}
$$
再作初等算术运算即得结论。
\end{proof}

\begin{lemma}
\label{groupoids-lemma-degree-multiplication-by-d}
设 $k$ 为域，$A$ 为 $k$ 上的阿贝尔簇，并设
$[d] : A \to A$ 为乘以 $d$ 的映射。则 $[d]$ 有限局部自由，
次数为 $d^{2\dim(A)}$。
\end{lemma}

\begin{proof}
由引理 \ref{groupoids-lemma-abelian-variety-projective}
（以及《态射进阶》引理 \ref{more-morphisms-lemma-projective}），$A$ 有一个充足可逆模
$\mathcal{L}$。因为 $[-1] : A \to A$ 是自同构，
$[-1]^*\mathcal{L}$ 也是充足可逆 $\mathcal{O}_A$-模。
因此 $\mathcal{N} = \mathcal{L} \otimes [-1]^*\mathcal{L}$ 充足；
见《概形的性质》引理 \ref{properties-lemma-ample-tensor-globally-generated}。
由于 $\mathcal{N} \cong [-1]^*\mathcal{N}$，由
引理 \ref{groupoids-lemma-pullbacks-by-n} 有
$[d]^*\mathcal{N} \cong \mathcal{N}^{\otimes d^2}$。

\medskip\noindent
为导出矛盾，设 $C$ 是包含在 $[d]$ 的某一纤维中的固有曲线。
此时 $\mathcal{N}^{\otimes d^2}|_C \cong \mathcal{O}_C$
是次数为 $0$ 的充足可逆 $\mathcal{O}_C$-模，这例如与
《代数簇》引理 \ref{varieties-lemma-ample-curve} 矛盾。
（也可使用《代数簇》引理 \ref{varieties-lemma-ample-positive}。）
因此 $[d]$ 的每个纤维维数都是 $0$，故例如由
《概形上同调》引理 \ref{coherent-lemma-characterize-finite}，
$[d]$ 有限。此外，由引理 \ref{groupoids-lemma-abelian-variety-smooth}，
$A$ 在 $k$ 上光滑，因而根据
《交换代数》引理 \ref{algebra-lemma-CM-over-regular-flat}，
$[d] : A \to A$ 平坦（这里还使用了域上光滑的概形正则，以及
正则环是 Cohen--Macaulay 环；见
《代数簇》引理 \ref{varieties-lemma-smooth-regular} 和
《交换代数》引理 \ref{algebra-lemma-regular-ring-CM}）。
所以由《概形态射》引理 \ref{morphisms-lemma-finite-flat}，
$[d]$ 有限局部自由。

\medskip\noindent
最后计算次数。由《代数簇》引理 \ref{varieties-lemma-degree-finite-morphism-in-terms-degrees}，
$$
\deg_{\mathcal{N}^{\otimes d^2}}(A) = \deg([d]) \deg_\mathcal{N}(A)
$$
由于 $A$ 关于 $\mathcal{N}^{\otimes d^2}$（相应地，
$\mathcal{N}$）的次数，是下列多项式中 $n^{\dim(A)}$ 的系数：
$$
n \longmapsto \chi(A, \mathcal{N}^{\otimes nd^2}),\quad
\text{分别}\quad n \longmapsto \chi(A, \mathcal{N}^{\otimes n})
$$
所以 $\deg([d]) = d^{2 \dim(A)}$。
\end{proof}

\begin{lemma}
\label{groupoids-lemma-abelian-variety-multiplication-by-d-etale}
\begin{slogan}
阿贝尔簇上乘以一个整数的态射为平展态射，当且仅当该整数在基域中可逆。
\end{slogan}
设 $k$ 为域，$A$ 为 $k$ 上的非零阿贝尔簇。
则 $[d] : A \to A$ 为平展（\'etale）态射，当且仅当 $d$ 在 $k$ 中可逆。
\end{lemma}

\begin{proof}
注意 $[d](x + y) = [d](x) + [d](y)$。由于由一点作平移是 $A$
的自同构，$[d] : A \to A$ 为平展（\'etale）的点集要么为空，
要么等于 $A$（略去一些细节）。因此，只需检验 $[d]$
在单位元 $e \in A(k)$ 处是否平展（\'etale）。
已知 $[d]$ 有限局部自由
（引理 \ref{groupoids-lemma-degree-multiplication-by-d}），故它在 $e$ 处
平展（\'etale），等价于证明
$\text{d}[d] : T_{A/k, e} \to T_{A/k, e}$ 是单射；见
《代数簇》引理 \ref{varieties-lemma-injective-tangent-spaces-unramified}
以及《概形态射》引理 \ref{morphisms-lemma-flat-unramified-etale}。
由引理 \ref{groupoids-lemma-group-scheme-addition-tangent-vectors}，
$\text{d}[d]$ 是 $T_{A/k, e}$ 上乘以 $d$ 的映射。
\end{proof}

\begin{lemma}
\label{groupoids-lemma-abelian-variety-multiplication-by-p}
设 $k$ 为特征 $p > 0$ 的域，$A$ 为 $k$ 上维数为 $g$ 的阿贝尔簇。
则 $[p] : A \to A$ 在 $0$ 上的纤维至多有 $p^g$ 个不同的点。
\end{lemma}

\begin{proof}
可将并且确实将 $k$ 替换为其代数闭包。由
引理 \ref{groupoids-lemma-group-scheme-addition-tangent-vectors}，
$[p]$ 的导数作为映射 $T_{A/k, e} \to T_{A/k, e}$
是乘以 $p$，因而为零（比较引理 \ref{groupoids-lemma-abelian-variety-multiplication-by-d-etale} 的证明）。
由于 $[p]$ 与平移可交换，$[p]$ 的导数处处为零；也就是说，诱导映射
$[p]^*\Omega_{A/k} \to \Omega_{A/k}$ 为零。
考察泛点可知，相应的函数域映射
$[p]^* : k(A) \to k(A)$ 在 $\Omega_{k(A)/k}$ 上诱导零映射。
设 $t_1, \ldots, t_g$ 为 $k(A)$ 在 $k$ 上的一个 p-基
（《代数进阶》定义 \ref{more-algebra-definition-p-basis}
以及引理 \ref{more-algebra-lemma-p-basis}）。
由《交换代数》引理 \ref{algebra-lemma-derivative-zero-pth-power}，
$[p]^*(t_i)$ 有一个 $p$ 次根。因此，
$k(A)[x_1, \ldots, x_g]/(x_1^p - t_1, \ldots, x_g^p - t_g)$
是 $[p]^* : k(A) \to k(A)$ 的一个子扩张。
于是可找到仿射开集 $U \subset A$，使得
$t_i \in \mathcal{O}_A(U)$ 且
$x_i \in \mathcal{O}_A([p]^{-1}(U))$。得到 $[p]$ 在 $U$ 上的分解
$$
[p]^{-1}(U)
\xrightarrow{\pi_1}
\Spec(\mathcal{O}(U)[x_1, \ldots, x_g]/(x_1^p - t_1, \ldots, x_g^p - t_g))
\xrightarrow{\pi_2}
U
$$
缩小 $U$ 后，可假设 $\pi_1$ 有限局部自由
（例如使用一般平坦性；实际上在本情形中它已经有限局部自由）。
由引理 \ref{groupoids-lemma-degree-multiplication-by-d}，$[p]$ 的次数为
$p^{2g}$。由于 $\pi_2$ 的次数为 $p^g$，$\pi_1$ 的次数也为 $p^g$。
态射 $\pi_2$ 是万有同胚，因而其纤维都是单点集。
所以 $[p]^{-1}(U) \to U$ 的集合论纤维就是 $\pi_1$ 的纤维，
至多有 $p^g$ 个元素。由于 $[p]$ 是 $k$ 上的群概形同态，
对每个 $a \in A(k)$，$[p] : A(k) \to A(k)$ 的纤维基数都相同，
证明完毕。
\end{proof}

\begin{proposition}
\label{groupoids-proposition-review-abelian-varieties}
\begin{reference}
\cite{AVar} 对此作了精彩阐释。
\end{reference}
设 $A$ 为域 $k$ 上的阿贝尔簇。则
\begin{enumerate}
\item $A$ 在 $k$ 上射影；
\item $A$ 是交换群概形；
\item 对所有 $n \geq 1$，态射 $[n] : A \to A$ 都是满射；
\item 若 $k$ 代数闭，则 $A(k)$ 是可除阿贝尔群；
\item $A[n] = \Ker([n] : A \to A)$ 是 $k$ 上次数为
$n^{2\dim A}$ 的有限群概形；
\item $A[n]$ 在 $k$ 上平展（\'etale），当且仅当 $n \in k^*$；
\item 若 $n \in k^*$ 且 $k$ 代数闭，则
$A(k)[n] \cong (\mathbf{Z}/n\mathbf{Z})^{\oplus 2\dim(A)}$；
\item 若 $k$ 是特征 $p > 0$ 的代数闭域，则存在整数
$0 \leq f \leq \dim(A)$，使得对所有 $m \geq 1$ 都有
$A(k)[p^m] \cong (\mathbf{Z}/p^m\mathbf{Z})^{\oplus f}$。
\end{enumerate}
\end{proposition}

\begin{proof}
(1) 由引理 \ref{groupoids-lemma-abelian-variety-projective} 得出。
(2) 由引理 \ref{groupoids-lemma-abelian-variety-abelian} 得出。
(3) 由引理 \ref{groupoids-lemma-degree-multiplication-by-d} 得出。
若 $k$ 代数闭，则 $k$ 上簇的满态射在 $k$-有理点上诱导满射，
故由 (3) 得 (4)。
(5) 由引理 \ref{groupoids-lemma-degree-multiplication-by-d} 以及如下事实得出：
次数为 $N$ 的有限局部自由态射经基变换后，仍是次数为 $N$
的有限局部自由态射。(6) 由
引理 \ref{groupoids-lemma-abelian-variety-multiplication-by-d-etale} 得出。
具体地，若 $n$ 在 $k$ 中可逆，则 $[n]$ 平展（\'etale），
因而 $A[n]$ 在 $k$ 上平展（\'etale）。
另一方面，若 $n$ 在 $k$ 中不可逆，则 $[n]$ 在 $e$ 处不平展
（\'etale），从而 $A[n]$ 在 $k$ 上的 $e$ 处不平展（\'etale）
（使用《概形态射》引理 \ref{morphisms-lemma-flat-unramified-etale} 和
\ref{morphisms-lemma-set-points-where-fibres-unramified}）。

\medskip\noindent
假设 $k$ 代数闭，令 $g = \dim(A)$。证明 (7)。
设 $\ell$ 是在 $k$ 中可逆的素数。于是
$$
A[\ell](k) = A(k)[\ell]
$$
是被 $\ell$ 零化、阶为 $\ell^{2g}$ 的有限阿贝尔群。
由有限阿贝尔群结构定理，它同构于
$(\mathbf{Z}/\ell\mathbf{Z})^{2g}$。接着考虑短正合列
$$
0 \to A(k)[\ell] \to A(k)[\ell^2] \xrightarrow{\ell} A(k)[\ell] \to 0
$$
用与上面相同的论证可得
$A(k)[\ell^2] \cong (\mathbf{Z}/\ell^2\mathbf{Z})^{2g}$。
对指数归纳，得到
$A(k)[\ell^m] \cong (\mathbf{Z}/\ell^m\mathbf{Z})^{2g}$。
对与 $k$ 的特征互素的合数 $n$，取各主分量，即得到
$A(k)$ 中 $n$-挠子群的正确形式。(8) 的证明完全相同；
使用引理 \ref{groupoids-lemma-abelian-variety-multiplication-by-p} 所给出的
$A(k)[p] \cong (\mathbf{Z}/p\mathbf{Z})^{\oplus f}$，
其中某个 $0 \leq f \leq g$。
\end{proof}

\begin{remark}
\label{groupoids-remark-16-equivalent}
设 $k$ 为域。阿贝尔簇有 $2 \times 4 \times 2 = 16$ 个等价定义。
考虑以下三组性质：
\begin{itemize}
\item 射影、固有；
\item 几何不可约、不可约、几何连通、连通；
\item 光滑、几何约化。
\end{itemize}
从每组各选一个性质，并设 $A$ 为 $k$ 上的群概形，
且这些所选性质都在 $k$ 上成立，
则 $A$ 是阿贝尔簇。若选取“固有、几何不可约、几何约化”，
便恢复定义 \ref{groupoids-definition-abelian-variety}
（使用《代数簇》引理 \ref{varieties-lemma-geometrically-integral}）。
最弱的选择是“固有、连通、几何约化”；例如参见
《概形态射》引理 \ref{morphisms-lemma-locally-projective-proper} 和
《代数簇》引理 \ref{varieties-lemma-smooth-geometrically-normal}。
因此，设 $A$ 是 $k$ 上固有、连通且几何约化的群概形。
由引理 \ref{groupoids-lemma-connected-group-scheme-over-field-irreducible} 和
\ref{groupoids-lemma-group-scheme-field-geometrically-irreducible}，
$A$ 几何不可约，因而是阿贝尔簇。
最后，若 $A/k$ 是阿贝尔簇，则由命题 \ref{groupoids-proposition-review-abelian-varieties}，它在 $k$ 上射影且光滑，
因而满足最强的选择“射影、几何不可约、光滑”。
\end{remark}







\section{群概形的作用}
\label{groupoids-section-action-group-scheme}

\noindent
设 $(G, m)$ 为群，$V$ 为集合。回顾 $G$ 在 $V$ 上的一个
{\it （左）作用}由映射 $a : G \times V \to V$ 给出，并满足
\begin{enumerate}
\item （结合律）对所有 $g, g' \in G$ 和 $v \in V$，
$a(m(g, g'), v) = a(g, a(g', v))$；
\item （单位元）对所有 $v \in V$，$a(e, v) = v$。
\end{enumerate}
也称 $V$ 是一个{\it $G$-集}（这通常意味着从记号中省略 $a$，
是一种记号滥用）。{\it $G$-集的映射}
$\psi : V \to V'$ 是满足如下条件的任意集合映射：
对所有 $v \in V$ 都有 $\psi(a(g, v)) = a(g, \psi(v))$。

\begin{definition}
\label{groupoids-definition-action-group-scheme}
设 $S$ 为概形，$(G, m)$ 为 $S$ 上的群概形。
\begin{enumerate}
\item {\it $G$ 在概形 $X/S$ 上的作用}是 $S$ 上的态射
$a : G \times_S X \to X$，使得对每个 $T/S$，映射
$a : G(T) \times X(T) \to X(T)$ 都在 $X(T)$ 上定义
$G(T)$-集结构。
\item 设 $X$、$Y$ 是 $S$ 上各自带有 $G$-作用的概形。
一个{\it 等变}态射，更准确地说，一个{\it $G$-等变态射}
$\psi : X \to Y$，是 $S$ 上的概形态射，并使得对每个 $T/S$，
$\psi : X(T) \to Y(T)$ 都是 $G(T)$-集的态射。
\end{enumerate}
\end{definition}

\noindent
在情形 (1) 中，这表示下列诸图交换：
\begin{equation}
\label{groupoids-equation-action}
\vcenter{
\xymatrix{
G \times_S G \times_S X \ar[r]_-{1_G \times a} \ar[d]_{m \times 1_X} &
G \times_S X \ar[d]^a \\
G \times_S X \ar[r]^a & X
}
}
\quad\quad
\vcenter{
\xymatrix{
G \times_S X \ar[r]_-a & X \\
X\ar[u]^{e \times 1_X} \ar[ru]_{1_X}
}
}
\end{equation}
在情形 (2) 中，这仅表示下图交换：
$$
\xymatrix{
G \times_S X \ar[r]_-{\text{id} \times \psi} \ar[d]_a &
G \times_S Y \ar[d]^a \\
X \ar[r]^\psi & Y
}
$$

\begin{definition}
\label{groupoids-definition-free-action}
设 $S$、$G \to S$、$X \to S$ 如
定义 \ref{groupoids-definition-action-group-scheme} 所述。
设 $a : G \times_S X \to X$ 是 $G$ 在 $X/S$ 上的作用。
若对每个 $S$ 上的概形 $T$，作用
$a : G(T) \times X(T) \to X(T)$ 都是群 $G(T)$
在集合 $X(T)$ 上的自由作用，则称该作用是{\it 自由的}。
\end{definition}

\begin{lemma}
\label{groupoids-lemma-free-action}
情形如定义 \ref{groupoids-definition-free-action}。
作用 $a$ 自由，当且仅当
$$
G \times_S X \to X \times_S X, \quad (g, x) \mapsto (a(g, x), x)
$$
是单态射。
\end{lemma}

\begin{proof}
由定义立即可得。
\end{proof}






\section{主齐性空间}
\label{groupoids-section-principal-homogeneous}

\noindent
在《位点上同调》定义 \ref{sites-cohomology-definition-torsor} 中，我们定义了一个位点上
群层的扭子。设 $\tau \in \{Zariski, \etale, smooth, syntomic, fppf\}$
为一种拓扑，而 $(G, m)$ 是 $S$ 上的群概形。
由于 $\tau$ 强于典范拓扑（见《下降》引理 \ref{descent-lemma-fpqc-universal-effective-epimorphisms}），
$\underline{G}$（见《位点与层》定义 \ref{sites-definition-representable-sheaf}）是
$(\Sch/S)_\tau$ 上的群层。因此，我们已经知道
$(\Sch/S)_\tau$ 上 $\underline{G}$-扭子的含义。
若该层可表示，便出现一种特殊情形。下述定义直接说明：
表示它的概形何时为 $G$-扭子。

\begin{definition}
\label{groupoids-definition-pseudo-torsor}
设 $S$ 为概形，$(G, m)$ 为 $S$ 上的群概形，$X$ 为 $S$ 上的概形，
并设 $a : G \times_S X \to X$ 是 $G$ 在 $X$ 上的作用。
\begin{enumerate}
\item 若诱导的 $S$ 上概形态射
$G \times_S X \to X \times_S X$，
$(g, x) \mapsto (a(g, x), x)$ 是同构，则称 $X$ 是
{\it 伪 $G$-扭子}，或称 $X$ {\it 在 $G$ 下形式主齐}。
\item 若存在 $S$ 上的 $G$-等变同构 $G \to X$，其中 $G$
通过左乘作用于 $G$，则称伪 $G$-扭子 $X$ 是{\it 平凡的}。
\end{enumerate}
\end{definition}

\noindent
显然，若 $S' \to S$ 是概形态射，则 $S$ 上伪 $G$-扭子的拉回
$X_{S'}$ 是 $S'$ 上的伪 $G_{S'}$-扭子。

\begin{lemma}
\label{groupoids-lemma-characterize-trivial-pseudo-torsors}
情形如定义 \ref{groupoids-definition-pseudo-torsor}。
\begin{enumerate}
\item 概形 $X$ 是伪 $G$-扭子，当且仅当对 $S$ 上的每个概形 $T$，
集合 $X(T)$ 要么为空，要么群 $G(T)$ 在 $X(T)$ 上的作用是单传递的。
\item 伪 $G$-扭子 $X$ 平凡，当且仅当态射 $X \to S$ 有截面。
\end{enumerate}
\end{lemma}

\begin{proof}
证明略。
\end{proof}

\begin{definition}
\label{groupoids-definition-principal-homogeneous-space}
设 $S$ 为概形，$(G, m)$ 为 $S$ 上的群概形，$X$ 为
$S$ 上的伪 $G$-扭子。
\begin{enumerate}
\item 若存在 fpqc 覆盖\footnote{这表示默认类型的扭子是在某个
fpqc 覆盖上平凡的伪扭子。这是
\cite[Expos\'e IV, 6.5]{SGA3} 中的定义。由于我们最常在 fppf
拓扑中工作，这一定义对我们略有不便。}
$\{S_i \to S\}_{i \in I}$，使得每个
$X_{S_i} \to S_i$ 都有截面（即为平凡伪 $G_{S_i}$-扭子），
则称 $X$ 是{\it 主齐性空间}或{\it $G$-扭子}。
\item 设 $\tau \in \{Zariski, \etale, smooth, syntomic, fppf\}$。
若存在 $\tau$-覆盖 $\{S_i \to S\}_{i \in I}$，使得每个
$X_{S_i} \to S_i$ 都有截面，则称 $X$ 是
{\it $\tau$ 拓扑中的 $G$-扭子}、{\it $\tau$ $G$-扭子}，
或简称{\it $\tau$ 扭子}。
\item 若 $X$ 是 $G$-扭子，并且它是 \'etale 拓扑的扭子，
则称其{\it 拟等平凡}。
\item 若 $X$ 是 $G$-扭子，并且它是 Zariski 拓扑的扭子，
则称其{\it 局部平凡}。
\end{enumerate}
\end{definition}

\noindent
我们有时说“设 $X$ 是 $S$ 上的 $G$-扭子”，意指 $X$ 是
$S$ 上带有 $G$-作用的概形，并且该作用使它成为 $S$ 上的主齐性空间。
下面说明，当两种说法都适用时，这与先前引入的记号一致。

\begin{lemma}
\label{groupoids-lemma-torsor}
设 $S$ 为概形，$(G, m)$ 为 $S$ 上的群概形，$X$ 为 $S$ 上的概形，
并设 $a : G \times_S X \to X$ 是 $G$ 在 $X$ 上的作用。
设 $\tau \in \{Zariski, \etale, smooth, syntomic, fppf\}$。
则 $X$ 是 $\tau$-拓扑中的 $G$-扭子，当且仅当
$\underline{X}$ 是 $(\Sch/S)_\tau$ 上的 $\underline{G}$-扭子。
\end{lemma}

\begin{proof}
证明略。
\end{proof}

\begin{remark}
\label{groupoids-remark-fun-with-torsors}
设 $(G, m)$ 为概形 $S$ 上的群概形。在这种情形下，自然会问：
\begin{enumerate}
\item 若 $X \to S$ 是伪 $G$-扭子且 $X \to S$ 满射，
则 $X$ 是否必为 $G$-扭子？
\item $(\Sch/S)_{fppf}$ 上的每个 $\underline{G}$-扭子是否可表示？
换言之，每个 $\underline{G}$-扭子是否都来自某个 fppf $G$-扭子？
\item 每个 $G$-扭子是否都是
fppf（相应地，\ smooth；相应地，\ \'etale；相应地，\ Zariski）扭子？
\end{enumerate}
一般而言，这些问题的答案都是否定的。要得到肯定答案，需要对
$G \to S$ 施加附加条件。例如：若 $S$ 是一个域的谱，则 (1)
的答案为肯定，因为此时 $\{X \to S\}$ 是使 $X$ 平凡化的 fpqc 覆盖。
若 $G \to S$ 仿射，则 (2) 的答案为肯定
（由《下降》引理 \ref{descent-lemma-affine} 得出）。
若 $G = \text{GL}_{n, S}$，则 (3) 的答案为肯定；
事实上，任意 $\text{GL}_{n, S}$-扭子都局部平凡
（由《下降》引理 \ref{descent-lemma-finite-locally-free-descends} 得出）。
\end{remark}



\section{等变拟凝聚层}
\label{groupoids-section-equivariant}

\noindent
我们把“函数”看作与“空间”对偶。因此，空间的态射所诱导的函数映射
方向相反。此外，也把模层的截面看作“函数”。这自然导向下述定义中
所选的箭头方向。

\begin{definition}
\label{groupoids-definition-equivariant-module}
设 $S$ 为概形，$(G, m)$ 为 $S$ 上的群概形，并设
$a : G \times_S X \to X$ 是群概形 $G$ 在 $X/S$ 上的作用。
一个{\it $G$-等变拟凝聚 $\mathcal{O}_X$-模}，或简称
{\it 等变拟凝聚 $\mathcal{O}_X$-模}，是二元组
$(\mathcal{F}, \alpha)$，其中 $\mathcal{F}$ 是拟凝聚
$\mathcal{O}_X$-模，而 $\alpha$ 是
$\mathcal{O}_{G \times_S X}$-模同态
$$
\alpha : a^*\mathcal{F} \longrightarrow \text{pr}_1^*\mathcal{F}
$$
这里 $\text{pr}_1 : G \times_S X \to X$ 是投影，并且满足：
\begin{enumerate}
\item 下图在 $\mathcal{O}_{G \times_S G \times_S X}$-模范畴中交换：
$$
\xymatrix{
(1_G \times a)^*\text{pr}_1^*\mathcal{F} \ar[r]_-{\text{pr}_{12}^*\alpha} &
\text{pr}_2^*\mathcal{F} \\
(1_G \times a)^*a^*\mathcal{F} \ar[u]^{(1_G \times a)^*\alpha} \ar@{=}[r] &
(m \times 1_X)^*a^*\mathcal{F} \ar[u]_{(m \times 1_X)^*\alpha}
}
$$
\item 拉回
$$
(e \times 1_X)^*\alpha : \mathcal{F} \longrightarrow \mathcal{F}
$$
是恒等映射。
\end{enumerate}
有关说明，可与 Equation (\ref{groupoids-equation-action}) 中的相应图比较。
\end{definition}

\noindent
注意，第一个图的交换性保证 $(e \times 1_X)^*\alpha$
是 $\mathcal{F}$ 上的幂等算子；因此条件 (2) 仅仅要求它是同构。

\begin{lemma}
\label{groupoids-lemma-pullback-equivariant}
设 $S$ 为概形，$G$ 为 $S$ 上的群概形。
设 $f : Y \to X$ 是两个带 $G$-作用的 $S$-概形之间的
$G$-等变态射。规则
$(\mathcal{F}, \alpha) \mapsto (f^*\mathcal{F}, (1_G \times f)^*\alpha)$
定义了从 $G$-等变拟凝聚 $\mathcal{O}_X$-模范畴到
$G$-等变拟凝聚 $\mathcal{O}_Y$-模范畴的函子。
\end{lemma}

\begin{proof}
证明略。
\end{proof}

\noindent
下面给出一个例子。

\begin{example}
\label{groupoids-example-Gm-on-affine}
设 $A$ 为 $\mathbf{Z}$-分次环；也就是说，$A$ 带有直和分解
$A = \bigoplus_{n \in \mathbf{Z}} A_n$，且
$A_n \cdot A_m \subset A_{n + m}$。
令 $X = \Spec(A)$。环同态
$\mu : A \to A \otimes \mathbf{Z}[x, x^{-1}]$，
$f \mapsto f \otimes x^{\deg(f)}$，给出一个 $\mathbf{G}_m$-作用
$$
a : \mathbf{G}_m \times X \longrightarrow X
$$
。确切地说，只需验证
$$
\xymatrix{
A \ar[r]_\mu \ar[d]_\mu &
A \otimes \mathbf{Z}[x, x^{-1}] \ar[d]^{\mu \otimes 1} \\
A \otimes \mathbf{Z}[x, x^{-1}] \ar[r]^-{1 \otimes m} &
A \otimes \mathbf{Z}[x, x^{-1}] \otimes \mathbf{Z}[x, x^{-1}]
}
$$
其中 $m(x) = x \otimes x$；见例 \ref{groupoids-example-multiplicative-group}。在齐次元上求值便立即可见。
设 $M$ 是分次 $A$-模。使用定义 \ref{groupoids-definition-equivariant-module} 中的 $\alpha$，即可得到
$\mathbf{G}_m$-等变拟凝聚 $\mathcal{O}_X$-模
$\mathcal{F} = \widetilde{M}$；该映射对应于
$A \otimes \mathbf{Z}[x, x^{-1}]$-模同态
$$
M \otimes_{A, \mu} (A \otimes \mathbf{Z}[x, x^{-1}])
\longrightarrow
M \otimes_{A, \text{id}_A \otimes 1} (A \otimes \mathbf{Z}[x, x^{-1}])
$$
它对齐次的 $m \in M$，把 $m \otimes 1 \otimes 1$
映到 $m \otimes 1 \otimes x^{\deg(m)}$。
\end{example}

\begin{lemma}
\label{groupoids-lemma-complete-reducibility-Gm}
设 $a : \mathbf{G}_m \times X \to X$ 是仿射概形上的作用。
则 $X$ 是某个 $\mathbf{Z}$-分次环的谱，且该作用如
例 \ref{groupoids-example-Gm-on-affine} 所述。
\end{lemma}

\begin{proof}
设 $f \in A = \Gamma(X, \mathcal{O}_X)$。可以唯一地写成有限和
$$
a^\sharp(f) = \sum\nolimits_{n \in \mathbf{Z}} f_n \otimes x^n
\quad\text{于}\quad
A \otimes \mathbf{Z}[x, x^{-1}] =
\Gamma(\mathbf{G}_m \times X, \mathcal{O}_{\mathbf{G}_m \times X})
$$
其中 $f_n$ 属于 $A$。于是得到映射 $A \to A$，$f \mapsto f_n$。
由于 $a$ 是作用，在 $x = 1$ 处求值可得 $f = \sum f_n$。
再由 $a$ 是作用，有
$$
\sum (f_n)_m \otimes x^m \otimes x^n = \sum f_n x^n \otimes x^n
$$
（比较例 \ref{groupoids-example-Gm-on-affine} 中的计算）。
所以当 $n \not = m$ 时 $(f_n)_m = 0$，且 $(f_n)_n = f_n$。
令
$$
A_n = \{f \in A \mid f_n = f\}
$$
则 $A = \sum A_n$。另一方面，在上述情形中，$f = 0$
蕴含 $f_n = 0$，所以该和必为直和。
\end{proof}

\begin{lemma}
\label{groupoids-lemma-Gm-equivariant-module}
设 $A$ 为分次环，$X = \Spec(A)$ 带有例 \ref{groupoids-example-Gm-on-affine} 中的作用
$a : \mathbf{G}_m \times X \to X$。
设 $\mathcal{F}$ 为 $\mathbf{G}_m$-等变拟凝聚
$\mathcal{O}_X$-模。则 $M = \Gamma(X, \mathcal{F})$
具有典范分次，使其成为分次 $A$-模，并使同构
$\widetilde{M} \to \mathcal{F}$
（《概形》引理 \ref{schemes-lemma-quasi-coherent-affine}）
成为 $\mathbf{G}_m$-等变模的同构；其中 $\widetilde{M}$
上的 $\mathbf{G}_m$-等变结构来自
例 \ref{groupoids-example-Gm-on-affine}。
\end{lemma}

\begin{proof}
可以对模 $M$ 重复引理 \ref{groupoids-lemma-complete-reducibility-Gm}
的论证。也可以考虑概形
$(X', \mathcal{O}_{X'}) = (X, \mathcal{O}_X \oplus \mathcal{F})$，
其中把 $\mathcal{F}$ 看成平方为零的理想。
利用 $X$ 上以及 $\mathcal{F}$ 上的作用，可定义自然作用
$a' : \mathbf{G}_m \times X' \to X'$。然后对 $X'$ 应用
引理 \ref{groupoids-lemma-complete-reducibility-Gm} 即得结论。
（这一论证的优点是，它立即说明 $A$ 和 $M$ 上的分次相容；
也就是说，$M$ 是分次 $A$-模。）
细节略。
\end{proof}



\section{群胚}
\label{groupoids-section-groupoids}

\noindent
回顾群胚是每个态射均为同构的范畴；见
《范畴》定义 \ref{categories-definition-groupoid}。
因此，群胚有对象集 $\text{Ob}$、箭头集 $\text{Arrows}$、
{\it 源}与{\it 靶}映射
$s, t : \text{Arrows} \to \text{Ob}$，以及{\it 合成律}
$c : \text{Arrows} \times_{s, \text{Ob}, t} \text{Arrows}
\to \text{Arrows}$。这些映射恰好满足下列公理：
\begin{enumerate}
\item （结合律）作为映射
$\text{Arrows} \times_{s, \text{Ob}, t}
\text{Arrows} \times_{s, \text{Ob}, t}
\text{Arrows} \to \text{Arrows}$，有
$c \circ (1, c) = c \circ (c, 1)$；
\item （恒等）存在映射 $e : \text{Ob} \to \text{Arrows}$，使得
\begin{enumerate}
\item 作为映射 $\text{Ob} \to \text{Ob}$，有
$s \circ e = t \circ e = \text{id}$；
\item 作为映射 $\text{Arrows} \to \text{Arrows}$，有
$c \circ (1, e \circ s) = c \circ (e \circ t, 1) = 1$；
\end{enumerate}
\item （取逆）存在映射 $i : \text{Arrows} \to \text{Arrows}$，使得
\begin{enumerate}
\item 作为映射 $\text{Arrows} \to \text{Ob}$，有
$s \circ i = t$、$t \circ i = s$；且
\item 作为映射 $\text{Arrows} \to \text{Arrows}$，有
$c \circ (1, i) = e \circ t$ 以及
$c \circ (i, 1) = e \circ s$。
\end{enumerate}
\end{enumerate}
若这些条件成立，则映射 $e$ 和 $i$ 唯一确定，且 $i$ 是双射。
注意，若 $(\text{Ob}', \text{Arrows}', s', t', c')$
是另一个群胚范畴，则函子
$f : (\text{Ob}, \text{Arrows}, s, t, c) \to
(\text{Ob}', \text{Arrows}', s', t', c')$
由一对集合映射 $f : \text{Ob} \to \text{Ob}'$ 与
$f : \text{Arrows} \to \text{Arrows}'$ 给出，它们满足
$s' \circ f = f \circ s$、$t' \circ f = f \circ t$ 以及
$c' \circ (f, f) = f \circ c$。与恒等和取逆的相容性自动成立。
下文将使用这一点。（警告：对一般范畴，必须另行要求与恒等的相容性。）

\begin{definition}
\label{groupoids-definition-groupoid}
设 $S$ 为概形。
\begin{enumerate}
\item 一个{\it $S$ 上的群胚概形}，或简称
{\it $S$ 上的群胚}，是五元组 $(U, R, s, t, c)$，其中
$U$ 和 $R$ 是 $S$ 上的概形，而
$s, t : R \to U$ 以及 $c : R \times_{s, U, t} R \to R$
是 $S$ 上的概形态射，并满足如下性质：对任意 $S$ 上的概形
$T$，五元组
$$
(U(T), R(T), s, t, c)
$$
都是上述意义下的群胚范畴。
\item 一个{\it $S$ 上群胚概形的态射
$f : (U, R, s, t, c) \to (U', R', s', t', c')$}
由概形态射 $f : U \to U'$ 与 $f : R \to R'$ 给出，
并满足如下性质：对任意 $S$ 上的概形 $T$，映射 $f$
定义了从群胚范畴 $(U(T), R(T), s, t, c)$ 到群胚范畴
$(U'(T), R'(T), s', t', c')$ 的函子。
\end{enumerate}
\end{definition}

\noindent
设 $(U, R, s, t, c)$ 为 $S$ 上的群胚。
由定义前的讨论以及米田引理，存在唯一的 $S$ 上概形态射
$e : U \to R$ 和 $i : R \to R$，使得对每个 $S$ 上的概形 $T$，
诱导映射 $e : U(T) \to R(T)$ 是恒等，而
$i : R(T) \to R(T)$ 是群胚范畴中的取逆。
七元组 $(U, R, s, t, c, e, i)$ 满足对应于上述公理
(1)、(2)(a)、(2)(b)、(3)(a)、(3)(b) 的各交换图。
反之，具有这一性质的七元组，其五元组 $(U, R, s, t, c)$
是群胚概形。注意，$i$ 是同构，而 $e$ 同时是 $s$ 和 $t$ 的截面。
此外，对一个给定的 $S$ 上群胚概形，记
$$
j = (t, s) : R \longrightarrow U \times_S U
$$
这与上文第 \ref{groupoids-section-equivalence-relations} 节中的约定相容。
我们有时说“设 $(U, R, s, t, c, e, i)$ 为 $S$ 上的群胚”，
以强调恒等与取逆的存在。

\begin{lemma}
\label{groupoids-lemma-groupoid-pre-equivalence}
给定 $S$ 上的群胚概形 $(U, R, s, t, c)$，态射
$j : R \to U \times_S U$ 是预等价关系。
\end{lemma}

\begin{proof}
证明略。这是一个很好的定义练习。
\end{proof}

\begin{lemma}
\label{groupoids-lemma-equivalence-groupoid}
给定 $S$ 上的等价关系 $j : R \to U \times_S U$，
存在唯一的方式把它扩张为 $S$ 上的群胚 $(U, R, s, t, c)$。
\end{lemma}

\begin{proof}
证明略。这是一个很好的定义练习。
\end{proof}

\begin{lemma}
\label{groupoids-lemma-diagram}
设 $S$ 为概形，$(U, R, s, t, c)$ 为 $S$ 上的群胚。
在交换图
$$
\xymatrix{
& U & \\
R \ar[d]_s \ar[ru]^t &
R \times_{s, U, t} R
\ar[l]^-{\text{pr}_0} \ar[d]^{\text{pr}_1} \ar[r]_-c &
R \ar[d]^s \ar[lu]_t \\
U & R \ar[l]_t \ar[r]^s & U
}
$$
中，下面两个方块都是纤维积方块。此外，上面的三角形
（实际上也是方块）同样是笛卡尔的。
\end{lemma}

\begin{proof}
证明略。这是关于定义以及代数几何函子观点的练习。
\end{proof}

\begin{lemma}
\label{groupoids-lemma-diagram-pull}
设 $S$ 为概形，$(U, R, s, t, c, e, i)$ 为 $S$ 上的群胚。图
\begin{equation}
\label{groupoids-equation-pull}
\xymatrix{
R \times_{t, U, t} R
\ar@<1ex>[r]^-{\text{pr}_1} \ar@<-1ex>[r]_-{\text{pr}_0}
\ar[d]_{(\text{pr}_0, c \circ (i, 1))} &
R \ar[r]^t \ar[d]^{\text{id}_R} &
U \ar[d]^{\text{id}_U} \\
R \times_{s, U, t} R
\ar@<1ex>[r]^-c \ar@<-1ex>[r]_-{\text{pr}_0} \ar[d]_{\text{pr}_1} &
R \ar[r]^t \ar[d]^s &
U \\
R \ar@<1ex>[r]^s \ar@<-1ex>[r]_t &
U
}
\end{equation}
交换。给出的竖直映射使上面两行彼此同构。
左下方的两个方块都是笛卡尔的。
\end{lemma}

\begin{proof}
该图的交换性由群胚公理得出。用群胚的语言说，左上方的竖直箭头
把具有相同靶的一对态射 $(\alpha, \beta)$ 映到一对态射
$(\alpha, \alpha^{-1} \circ \beta)$。在任意群胚中，这定义了
$\text{Arrows} \times_{t, \text{Ob}, t} \text{Arrows}$
与
$\text{Arrows} \times_{s, \text{Ob}, t} \text{Arrows}$
之间的双射，故引理的第二个断言成立。最后一个断言由
引理 \ref{groupoids-lemma-diagram} 得出。
\end{proof}

\begin{lemma}
\label{groupoids-lemma-base-change-groupoid}
设 $(U, R, s, t, c)$ 为概形 $S$ 上的群胚，$S' \to S$ 为态射。
令 $U' = S' \times_S U$、$R' = S' \times_S R$，并赋予它们
态射 $s, t, c$ 的基变换 $s'$、$t'$、$c'$。
则 $(U', R', s', t', c')$ 是 $S'$ 上的群胚概形，而诸投影决定
$S$ 上群胚概形的态射
$(U', R', s', t', c') \to (U, R, s, t, c)$。
\end{lemma}

\begin{proof}
证明略。提示：
$R' \times_{s', U', t'} R' = S' \times_S (R \times_{s, U, t} R)$。
\end{proof}









\section{群胚上的拟凝聚层}
\label{groupoids-section-groupoids-quasi-coherent}

\noindent
关于箭头方向的约定，见第 \ref{groupoids-section-equivariant} 节的引言。

\begin{definition}
\label{groupoids-definition-groupoid-module}
设 $S$ 为概形，$(U, R, s, t, c)$ 为 $S$ 上的群胚概形。
所谓 $(U, R, s, t, c)$ 上的一个{\it 拟凝聚模}，是指一对
$(\mathcal{F}, \alpha)$，其中 $\mathcal{F}$ 是拟凝聚
$\mathcal{O}_U$-模，而 $\alpha$ 是 $\mathcal{O}_R$-模态射
$$
\alpha : t^*\mathcal{F} \longrightarrow s^*\mathcal{F}
$$
满足：
\begin{enumerate}
\item 图
$$
\xymatrix{
& \text{pr}_1^*t^*\mathcal{F} \ar[r]_-{\text{pr}_1^*\alpha} &
\text{pr}_1^*s^*\mathcal{F} \ar@{=}[rd] & \\
\text{pr}_0^*s^*\mathcal{F} \ar@{=}[ru] & & & c^*s^*\mathcal{F} \\
& \text{pr}_0^*t^*\mathcal{F} \ar[lu]^{\text{pr}_0^*\alpha} \ar@{=}[r] &
c^*t^*\mathcal{F} \ar[ru]_{c^*\alpha}
}
$$
在 $\mathcal{O}_{R \times_{s, U, t} R}$-模范畴中交换；
\item 拉回
$$
e^*\alpha : \mathcal{F} \longrightarrow \mathcal{F}
$$
是恒等态射。
\end{enumerate}
可与引理 \ref{groupoids-lemma-diagram} 中的交换图比较。
\end{definition}

\noindent
第一个图的交换性迫使算子 $e^*\alpha$ 幂等。因此，第二个条件也可表述为
$e^*\alpha$ 是同构。事实上，该条件蕴含 $\alpha$ 是同构。

\begin{lemma}
\label{groupoids-lemma-isomorphism}
设 $S$ 为概形，$(U, R, s, t, c)$ 为 $S$ 上的群胚概形。
若 $(\mathcal{F}, \alpha)$ 是 $(U, R, s, t, c)$ 上的拟凝聚模，
则 $\alpha$ 是同构。
\end{lemma}

\begin{proof}
沿态射 $(i, 1) : R \to R \times_{s, U, t} R$ 拉回
定义 \ref{groupoids-definition-groupoid-module} 中的交换图。
由此得到 $i^*\alpha \circ \alpha = s^*e^*\alpha$。
再沿态射 $(1, i)$ 拉回，得到关系
$\alpha \circ i^*\alpha = t^*e^*\alpha$。由第二个假设，
这些态射都是恒等态射。因此 $i^*\alpha$ 是 $\alpha$ 的逆。
\end{proof}

\begin{lemma}
\label{groupoids-lemma-pullback}
设 $S$ 为概形，并考虑 $S$ 上群胚概形的态射
$f : (U, R, s, t, c) \to (U', R', s', t', c')$。
则由
$$
(\mathcal{F}, \alpha) \mapsto (f^*\mathcal{F}, f^*\alpha)
$$
给出的拉回 $f^*$，定义了从 $(U', R', s', t', c')$ 上的拟凝聚层范畴
到 $(U, R, s, t, c)$ 上的拟凝聚层范畴的函子。
\end{lemma}

\begin{proof}
证明略。
\end{proof}

\begin{lemma}
\label{groupoids-lemma-pushforward}
设 $S$ 为概形，并考虑 $S$ 上群胚概形的态射
$f : (U, R, s, t, c) \to (U', R', s', t', c')$。
假设：
\begin{enumerate}
\item $f : U \to U'$ 拟紧且拟分离；
\item 方块
$$
\xymatrix{
R \ar[d]_t \ar[r]_f & R' \ar[d]^{t'} \\
U \ar[r]^f & U'
}
$$
是笛卡尔的；
\item $s'$ 和 $t'$ 平坦。
\end{enumerate}
则由
$$
(\mathcal{F}, \alpha) \mapsto (f_*\mathcal{F}, f_*\alpha)
$$
给出的推前 $f_*$，定义了从 $(U, R, s, t, c)$ 上的拟凝聚层范畴
到 $(U', R', s', t', c')$ 上的拟凝聚层范畴的函子；它是
引理 \ref{groupoids-lemma-pullback} 所定义的拉回的右伴随。
\end{lemma}

\begin{proof}
由于 $U \to U'$ 拟紧且拟分离，$f_*$ 把拟凝聚层变为拟凝聚层
（《概形》引理 \ref{schemes-lemma-push-forward-quasi-coherent}）。
又因方块
$$
\vcenter{
\xymatrix{
R \ar[d]_t \ar[r]_f & R' \ar[d]^{t'} \\
U \ar[r]^f & U'
}
}
\quad\text{且}\quad
\vcenter{
\xymatrix{
R \ar[d]_s \ar[r]_f & R' \ar[d]^{s'} \\
U \ar[r]^f & U'
}
}
$$
均为笛卡尔的，故
$(t')^*f_*\mathcal{F} = f_*t^*\mathcal{F}$ 以及
$(s')^*f_*\mathcal{F} = f_*s^*\mathcal{F}$，见
《概形上同调》引理 \ref{coherent-lemma-flat-base-change-cohomology}。
所以可把 $f_*\alpha$ 视为态射
$(t')^*f_*\mathcal{F} \to (s')^*f_*\mathcal{F}$。
类似论证表明 $f_*\alpha$ 满足上循环条件。
该函子与拉回函子伴随，因为环空间上模的拉回与推前彼此伴随。
其余细节略去。
\end{proof}

\begin{lemma}
\label{groupoids-lemma-colimits}
设 $S$ 为概形，$(U, R, s, t, c)$ 为 $S$ 上的群胚概形。
则 $(U, R, s, t, c)$ 上的拟凝聚模范畴有余极限。
\end{lemma}

\begin{proof}
设 $i \mapsto (\mathcal{F}_i, \alpha_i)$ 是以 $\mathcal{I}$ 为指标范畴的图。
可以形成余极限
$\mathcal{F} = \colim \mathcal{F}_i$；
它是 $U$ 上的拟凝聚层，见
《概形》第 \ref{schemes-section-quasi-coherent} 节。
由于余极限与拉回交换，有
$s^*\mathcal{F} = \colim s^*\mathcal{F}_i$，类似地还有
$t^*\mathcal{F} = \colim t^*\mathcal{F}_i$。因此可以令
$\alpha = \colim \alpha_i$。我们略去证明：在
$(U, R, s, t, c)$ 上的拟凝聚模范畴中，$(\mathcal{F}, \alpha)$
确为该图的余极限。
\end{proof}

\begin{lemma}
\label{groupoids-lemma-abelian}
设 $S$ 为概形，$(U, R, s, t, c)$ 为 $S$ 上的群胚概形。
若 $s$、$t$ 平坦，则 $(U, R, s, t, c)$ 上的拟凝聚模范畴
是阿贝尔范畴。
\end{lemma}

\begin{proof}
设
$\varphi : (\mathcal{F}, \alpha) \to (\mathcal{G}, \beta)$
是 $(U, R, s, t, c)$ 上拟凝聚模的同态。由于 $s$ 平坦，
$$
0 \to s^*\Ker(\varphi)
\to s^*\mathcal{F} \to s^*\mathcal{G} \to s^*\Coker(\varphi) \to 0
$$
正合；沿 $t$ 拉回亦然。因此 $\alpha$ 和 $\beta$ 诱导同构
$\kappa : t^*\Ker(\varphi) \to s^*\Ker(\varphi)$ 以及
$\lambda : t^*\Coker(\varphi) \to s^*\Coker(\varphi)$，
且二者满足上循环条件。直接验证可知，
$(\Ker(\varphi), \kappa)$ 和
$(\Coker(\varphi), \lambda)$ 分别是
$(U, R, s, t, c)$ 上拟凝聚模范畴中的核与余核。
此外，条件 $\Coim(\varphi) = \Im(\varphi)$ 成立，
因为它在 $U$ 上成立。
\end{proof}













\section{拟凝聚模的余极限}
\label{groupoids-section-colimits}

\noindent
本节证明若干技术性结果：在适当假设下，群胚上的每个拟凝聚模
都是“较小”拟凝聚模的滤过余极限。

\begin{lemma}
\label{groupoids-lemma-construct-quasi-coherent}
设 $(U, R, s, t, c)$ 为 $S$ 上的群胚概形。
假设 $s, t$ 平坦、拟紧且拟分离。对于 $U$ 上任意拟凝聚模
$\mathcal{G}$，存在典范同构
$\alpha : t^*s_*t^*\mathcal{G} \to s^*s_*t^*\mathcal{G}$，
使 $(s_*t^*\mathcal{G}, \alpha)$ 成为 $(U, R, s, t, c)$ 上的拟凝聚模。
这一构造定义函子
$$
\QCoh(\mathcal{O}_U) \longrightarrow \QCoh(U, R, s, t, c)
$$
且该函子是遗忘函子
$(\mathcal{F}, \beta) \mapsto \mathcal{F}$ 的右伴随。
\end{lemma}

\begin{proof}
拟凝聚模沿拟紧拟分离态射的推前仍是拟凝聚的，见
《概形》引理 \ref{schemes-lemma-push-forward-quasi-coherent}。
因此 $s_*t^*\mathcal{G}$ 拟凝聚。采用引理 \ref{groupoids-lemma-diagram}
中的记号，有
$$
t^*s_*t^*\mathcal{G} =
\text{pr}_{1, *}\text{pr}_0^*t^*\mathcal{G} =
\text{pr}_{1, *}c^*t^*\mathcal{G} =
s^*s_*t^*\mathcal{G}
$$
中间的等式源于态射
$R \times_{s, U, t} R \to U$ 满足
$t \circ c = t \circ \text{pr}_0$；首尾两个等式则源于这些步骤中
基变换与推前交换，见《概形上同调》引理 \ref{coherent-lemma-flat-base-change-cohomology}。

\medskip\noindent
为了验证 $\alpha$ 满足定义 \ref{groupoids-definition-groupoid-module}
中的上循环条件以及伴随性质，我们换一种方式描述构造
$\mathcal{G} \mapsto (s_*t^*\mathcal{G}, \alpha)$。
考虑由 $R$ 上等价关系 $R \times_{t, U, t} R$ 所关联的群胚概形
$(R, R \times_{t, U, t} R, \text{pr}_0, \text{pr}_1, \text{pr}_{02})$，
见引理 \ref{groupoids-lemma-equivalence-groupoid}。存在群胚概形态射
$$
f :
(R, R \times_{t, U, t} R, \text{pr}_1, \text{pr}_0, \text{pr}_{02})
\longrightarrow
(U, R, s, t, c)
$$
它在对象上由 $s : R \to U$ 给出，在箭头上由
$R \times_{t, U, t} R \to R$，
$(r_0, r_1) \mapsto r_0^{-1} \circ r_1$ 给出；
所需诸图交换性的验证从略。由于
$t, s : R \to U$ 拟紧、拟分离且平坦，并且由
引理 \ref{groupoids-lemma-diagram-pull} 有笛卡尔方块
$$
\xymatrix{
R \times_{t, U, t} R \ar[d]_{\text{pr}_0}
\ar[rr]_-{(r_0, r_1) \mapsto r_0^{-1} \circ r_1} & & R \ar[d]^t \\
R \ar[rr]^s & & U
}
$$
故引理 \ref{groupoids-lemma-pushforward} 适用于 $f$。
于是拟凝聚模沿 $f$ 的推前与拉回互为伴随函子。
为完成证明，我们把这些函子与上面描述的函子识别起来。注意
$$
t^* :
\QCoh(\mathcal{O}_U)
\longrightarrow
\QCoh(R, R \times_{t, U, t} R, \text{pr}_1, \text{pr}_0, \text{pr}_{02})
$$
是范畴等价；这是拟凝聚层的下降理论，因为
$\{t : R \to U\}$ 是 fpqc 覆盖，见
《下降》命题 \ref{descent-proposition-fpqc-descent-quasi-coherent}。

\medskip\noindent
沿 $f$ 推前并预复合等价 $t^*$，会把
$\mathcal{G}$ 送到 $(s_*t^*\mathcal{G}, \alpha)$。
我们略去验证以此方式得到的同构 $\alpha$ 与上面构造的同构相同。

\medskip\noindent
沿 $f$ 拉回并后复合等价 $t^*$ 的逆，会把
$(\mathcal{F}, \beta)$ 送到模 $s^*\mathcal{F}$
相对于 $\{t : R \to U\}$ 的下降；这里该模配备
$R \times_{t, U, t} R$ 上的下降数据
$\gamma$，它是 $\beta$ 沿
$(r_0, r_1) \mapsto r_0^{-1} \circ r_1$ 的拉回。
考虑同构 $\beta : t^*\mathcal{F} \to s^*\mathcal{F}$。
$t^*\mathcal{F}$ 相对于 $\{t : R \to U\}$ 的典范下降数据
（《下降》定义 \ref{descent-definition-descent-datum-effective-quasi-coherent}）
经 $\beta$ 转化为态射
$$
\text{pr}_0^*s^*\mathcal{F}
\xrightarrow{\text{pr}_0^*\beta^{-1}}
\text{pr}_0^*t^*\mathcal{F}
\xrightarrow{can}
\text{pr}_1^*t^*\mathcal{F}
\xrightarrow{\text{pr}_1^*\beta}
\text{pr}_1^*s^*\mathcal{F}
$$
由于 $\beta$ 满足上循环条件，此态射等于 $\beta$ 沿
$(r_0, r_1) \mapsto r_0^{-1} \circ r_1$ 的拉回。
为此，只需取定义 \ref{groupoids-definition-groupoid-module}
中实际的上循环关系，并沿态射
$(\text{pr}_0, c \circ (i, 1)) : R \times_{t, U, t} R \to R \times_{s, U, t} R$
将其拉回；该态射也出现在
引理 \ref{groupoids-lemma-diagram-pull} 的交换图中。因此
$(s^*\mathcal{F}, \gamma)$ 同构于 $(t^*\mathcal{F}, can)$。
综上，沿 $f$ 拉回再后复合等价 $t^*$ 的逆所得函子，
同构于遗忘函子
$(\mathcal{F}, \beta) \mapsto \mathcal{F}$。
\end{proof}

\begin{remark}
\label{groupoids-remark-adjunction-map}
在引理 \ref{groupoids-lemma-construct-quasi-coherent} 的情形下，记遗忘函子为
$$
F : \QCoh(U, R, s, t, c) \to \QCoh(\mathcal{O}_U),\quad
(\mathcal{F}, \beta) \mapsto \mathcal{F}
$$
并记该引理中构造的右伴随为
$$
G : \QCoh(\mathcal{O}_U) \to \QCoh(U, R, s, t, c),\quad
\mathcal{G} \mapsto (s_*t^*\mathcal{G}, \alpha)
$$
则伴随的单位 $\eta : \text{id} \to G \circ F$
在 $(\mathcal{F}, \beta)$ 上取值为态射
$$
\mathcal{F} \to s_*s^*\mathcal{F} \xrightarrow{\beta^{-1}} s_*t^*\mathcal{F}
$$
验证从略。
\end{remark}

\begin{lemma}
\label{groupoids-lemma-push-pull}
设 $f : Y \to X$ 为概形态射，$\mathcal{F}$ 为拟凝聚
$\mathcal{O}_X$-模，$\mathcal{G}$ 为拟凝聚
$\mathcal{O}_Y$-模，并设
$\varphi : \mathcal{G} \to f^*\mathcal{F}$ 为模态射。假设：
\begin{enumerate}
\item $\varphi$ 是单射；
\item $f$ 拟紧、拟分离、平坦且满射；
\item $X$, $Y$ 局部 Noether；
\item $\mathcal{G}$ 是凝聚 $\mathcal{O}_Y$-模。
\end{enumerate}
则由拉回
$$
\xymatrix{
\mathcal{F} \ar[r] & f_*f^*\mathcal{F} \\
\mathcal{F} \cap f_*\mathcal{G} \ar[u] \ar[r] &
f_*\mathcal{G} \ar[u]
}
$$
定义的 $\mathcal{F} \cap f_*\mathcal{G}$ 是凝聚
$\mathcal{O}_X$-模。
\end{lemma}

\begin{proof}
我们将直接使用《概形上同调》引理 \ref{coherent-lemma-coherent-Noetherian} 对凝聚模的刻画，
以及凝聚模构成 $\QCoh(\mathcal{O}_X)$ 的 Serre 子范畴这一事实；见
《概形上同调》引理 \ref{coherent-lemma-coherent-Noetherian-quasi-coherent-sub-quotient}。
若 $f$ 有截面 $\sigma$，则
$\mathcal{F} \cap f_*\mathcal{G}$ 包含于态射
$\sigma^*\mathcal{G} \to \sigma^*f^*\mathcal{F} = \mathcal{F}$
的像中，因而是凝聚的。一般地，为证明
$\mathcal{F} \cap f_*\mathcal{G}$ 凝聚，只需证明
$f^*(\mathcal{F} \cap f_*\mathcal{G})$ 凝聚（见
《下降》引理 \ref{descent-lemma-finite-type-descends}）。
由于 $f$ 平坦，它等于
$f^*\mathcal{F} \cap f^*f_*\mathcal{G}$。
又因 $f$ 平坦、拟紧且拟分离，故
$f^*f_*\mathcal{G} = p_*q^*\mathcal{G}$，其中
$p, q : Y \times_X Y \to Y$ 是投影；见
《概形上同调》引理 \ref{coherent-lemma-flat-base-change-cohomology}。
由于 $p$ 有截面，结论成立。
\end{proof}

\noindent
设 $S$ 为概形，$(U, R, s, t, c)$ 为 $S$ 上的概形群胚。
假设 $U$ 局部 Noether。在下面的引理中，如果
$(U, R, s, t, c)$ 上拟凝聚层 $(\mathcal{F}, \alpha)$ 的
$\mathcal{O}_U$-模 $\mathcal{F}$ 是凝聚的，我们就称它是
{\it 凝聚的}。

\begin{lemma}
\label{groupoids-lemma-colimit-coherent}
设 $(U, R, s, t, c)$ 为 $S$ 上的群胚概形。假设：
\begin{enumerate}
\item $U$, $R$ 是 Noether 概形；
\item $s, t$ 平坦、拟紧且拟分离。
\end{enumerate}
则 $(U, R, s, t, c)$ 上的每个拟凝聚模
$(\mathcal{F}, \beta)$ 都是凝聚模的滤过余极限。
\end{lemma}

\begin{proof}
我们将不再特别说明而使用《概形上同调》引理 \ref{coherent-lemma-coherent-Noetherian} 对局部 Noether 概形上凝聚模的刻画。
由《概形上同调》引理 \ref{coherent-lemma-directed-colimit-coherent}，可将
$\mathcal{F} = \colim \mathcal{H}_i$ 写成凝聚子模
$\mathcal{H}_i \subset \mathcal{F}$ 的滤过余极限。
给定 $U$ 上拟凝聚层 $\mathcal{H}$，以
$(s_*t^*\mathcal{H}, \alpha)$ 表示
引理 \ref{groupoids-lemma-construct-quasi-coherent} 给出的
$(U, R, s, t, c)$ 上拟凝聚层。考虑
$\QCoh(U, R, s, t, c)$ 中的伴随态射
$(\mathcal{F}, \beta) \to (s_*t^*\mathcal{F}, \alpha)$，
见注 \ref{groupoids-remark-adjunction-map}。令
$$
(\mathcal{F}_i, \beta_i) =
(\mathcal{F}, \beta)
\times_{(s_*t^*\mathcal{F}, \alpha)}
(s_*t^*\mathcal{H}_i, \alpha)
$$
这是在 $\QCoh(U, R, s, t, c)$ 中取的纤维积。根据
引理 \ref{groupoids-lemma-abelian} 的证明，限制到 $U$ 是
$\QCoh(U, R, s, t, c)$ 上的正合函子，故得到拉回图
$$
\xymatrix{
\mathcal{F} \ar[r] & s_*t^*\mathcal{F} \\
\mathcal{F}_i \ar[r] \ar[u] & s_*t^*\mathcal{H}_i \ar[u]
}
$$
换言之，$\mathcal{F}_i = \mathcal{F} \cap s_*t^*\mathcal{H}_i$。
由注 \ref{groupoids-remark-adjunction-map} 中对伴随态射的描述，
该图同构于
$$
\xymatrix{
\mathcal{F} \ar[r] & s_*s^*\mathcal{F} \\
\mathcal{F}_i \ar[r] \ar[u] & s_*t^*\mathcal{H}_i \ar[u]
}
$$
其中右侧竖直箭头由对态射
$$
t^*\mathcal{H}_i \to t^*\mathcal{F} \xrightarrow{\beta} s^*\mathcal{F}
$$
施以 $s_*$ 得到。由于 $t$ 平坦，该箭头是单射。因此由
引理 \ref{groupoids-lemma-push-pull}，$\mathcal{F}_i$ 凝聚。
最后，因为 $s$ 拟紧且拟分离，$s_*$ 与余极限交换
（见《概形上同调》引理 \ref{coherent-lemma-colimit-cohomology}）。故
$s_*t^*\mathcal{F} = \colim s_*t^*\mathcal{H}_i$，进而如所需有
$(\mathcal{F}, \beta) = \colim (\mathcal{F}_i, \beta_i)$。
\end{proof}

\noindent
下面这个颇为有趣的引理在处理域上的群胚时很有用。事实上，
它正是证明代数群的任意表示都是有限维表示之余极限的标准论证。

\begin{lemma}
\label{groupoids-lemma-colimit-finite-type}
设 $(U, R, s, t, c)$ 为 $S$ 上的群胚概形。假设：
\begin{enumerate}
\item $U$, $R$ 仿射；
\item 存在 $e_i \in \mathcal{O}_R(R)$，使得每个元素
$g \in \mathcal{O}_R(R)$ 都能唯一写成
$\sum s^*(f_i)e_i$，其中 $f_i \in \mathcal{O}_U(U)$。
\end{enumerate}
则 $(U, R, s, t, c)$ 上每个拟凝聚模
$(\mathcal{F}, \alpha)$ 都是有限型拟凝聚模的滤过余极限。
\end{lemma}

\begin{proof}
该假设意指 $\mathcal{O}_R(R)$ 经由 $s$ 成为以 $e_i$ 为基的自由
$\mathcal{O}_U(U)$-模。因此对任意拟凝聚
$\mathcal{O}_U$-模 $\mathcal{G}$，有
$s^*\mathcal{G}(R) = \bigoplus_i \mathcal{G}(U)e_i$。
我们用 $s(-)$ 表示截面沿 $s$ 的拉回，其他态射亦采用类似记号。
设 $(\mathcal{F}, \alpha)$ 是 $(U, R, s, t, c)$ 上的拟凝聚模，
并取 $\sigma \in \mathcal{F}(U)$。由上可唯一写成
$$
\alpha(t(\sigma)) = \sum s(\sigma_i) e_i
$$
其中 $\sigma_i \in \mathcal{F}(U)$（当然，除有限多个外均为零）。
还可把 $R \times_{s, U, t}R$ 上的函数写成
$$
c(e_i) = \sum \text{pr}_1(f_{ij}) \text{pr}_0(e_j)
$$
定义 \ref{groupoids-definition-groupoid-module} 中图的交换性意味着
$$
\sum \text{pr}_1(\alpha(t(\sigma_i))) \text{pr}_0(e_i)
=
\sum \text{pr}_1(s(\sigma_i)f_{ij}) \text{pr}_0(e_j)
$$
（计算略）。取出 $\text{pr}_0(e_l)$ 的系数，得到
$\alpha(t(\sigma_l)) = \sum s(\sigma_i)f_{il}$。
于是，由诸元素 $\sigma_i$ 生成的子模
$\mathcal{G} \subset \mathcal{F}$ 定义一个由 $\alpha$ 保持的
有限型拟凝聚模，因而它是 $\QCoh(U, R, s, t, c)$ 中
$\mathcal{F}$ 的子对象。这个子模包含 $\sigma$
（将第一个关系沿 $e$ 拉回即可看出），故结论成立。
\end{proof}

\noindent
我们建议读者跳过本节的其余部分。设 $S$ 为概形，
$(U, R, s, t, c)$ 为 $S$ 上的概形群胚，$\kappa$ 为一个基数。
下文称 $(U, R, s, t, c)$ 上的拟凝聚层
$(\mathcal{F}, \alpha)$ 是 $\kappa$-生成的，是指
$\mathcal{F}$ 是 $\kappa$-生成的 $\mathcal{O}_U$-模；见
《概形的性质》定义 \ref{properties-definition-kappa-generated}。

\begin{lemma}
\label{groupoids-lemma-set-of-iso-classes}
设 $(U, R, s, t, c)$ 为 $S$ 上的群胚概形，$\kappa$ 为基数。
存在集合 $T$ 以及 $(U, R, s, t, c)$ 上一族
$\kappa$-生成拟凝聚模
$(\mathcal{F}_t, \alpha_t)_{t \in T}$，使得
$(U, R, s, t, c)$ 上每个 $\kappa$-生成拟凝聚模
都同构于某个 $(\mathcal{F}_t, \alpha_t)$。
\end{lemma}

\begin{proof}
对 $U$ 上每个拟凝聚模 $\mathcal{F}$，使
$(\mathcal{F}, \alpha)$ 成为 $(U, R, s, t, c)$ 上拟凝聚模的态射
$\alpha : t^*\mathcal{F} \to s^*\mathcal{F}$ 构成一个
（可能为空的）集合。由《概形的性质》引理 \ref{properties-lemma-set-of-iso-classes}，
$\kappa$-生成拟凝聚 $\mathcal{O}_U$-模的同构类构成一个集合。
\end{proof}

\begin{lemma}
\label{groupoids-lemma-colimit-kappa}
设 $(U, R, s, t, c)$ 为 $S$ 上的群胚概形，并假设 $s, t$ 平坦。
则存在基数 $\kappa$，使得 $(U, R, s, t, c)$ 上每个拟凝聚模
$(\mathcal{F}, \alpha)$ 都是其 $\kappa$-生成拟凝聚子模的有向余极限。
\end{lemma}

\begin{proof}
在引理的陈述及本证明中，拟凝聚模 $(\mathcal{F}, \alpha)$ 的
{\it 子模}，是指拟凝聚子模
$\mathcal{G} \subset \mathcal{F}$，满足作为
$s^*\mathcal{F}$ 的子层有
$\alpha(t^*\mathcal{G}) = s^*\mathcal{G}$。
这是有意义的，因为 $s, t$ 平坦，所以拉回 $s^*$ 与 $t^*$
正合，也就是说它们保持子层。本证明将重复
《概形的性质》引理 \ref{properties-lemma-colimit-kappa} 的证明。
我们建议读者先阅读该证明。

\medskip\noindent
选择仿射开覆盖 $U = \bigcup_{i \in I} U_i$。
对每一对 $i, j$，选择仿射开覆盖
$$
U_i \cap U_j = \bigcup\nolimits_{k \in I_{ij}} U_{ijk}
\quad\text{且}\quad
s^{-1}(U_i) \cap t^{-1}(U_j) = \bigcup\nolimits_{k \in J_{ij}} W_{ijk}.
$$
记 $U_i = \Spec(A_i)$、$U_{ijk} = \Spec(A_{ijk})$、
$W_{ijk} = \Spec(B_{ijk})$。令 $\kappa$ 为任意无限基数，
它 $\geq$ 任一集合 $I$, $I_{ij}$, $J_{ij}$ 的基数。

\medskip\noindent
设 $(\mathcal{F}, \alpha)$ 为 $(U, R, s, t, c)$ 上的拟凝聚模。
令 $M_i = \mathcal{F}(U_i)$、$M_{ijk} = \mathcal{F}(U_{ijk})$。
注意
$$
M_i \otimes_{A_i} A_{ijk} = M_{ijk} = M_j \otimes_{A_j} A_{ijk}
$$
并且 $\alpha$ 给出同构
$$
\alpha|_{W_{ijk}} :
M_i \otimes_{A_i, t} B_{ijk}
\longrightarrow
M_j \otimes_{A_j, s} B_{ijk}
$$
见《概形》引理 \ref{schemes-lemma-widetilde-pullback}。
利用选择公理，我们选择映射
$$
(i, j, k, m) \mapsto S(i, j, k, m)
$$
使它对每个 $i, j \in I$、$k \in I_{ij}$ 或 $k \in J_{ij}$
以及 $m \in M_i$，指定有限子集
$S(i, j, k, m) \subset M_j$，从而有
$$
m \otimes 1 = \sum\nolimits_{m' \in S(i, j, k, m)} m' \otimes a_{m'}
\quad\text{或}\quad
\alpha(m \otimes 1) = \sum\nolimits_{m' \in S(i, j, k, m)} m' \otimes b_{m'}
$$
在 $M_{ijk}$ 中成立，其中
$a_{m'} \in A_{ijk}$ 或 $b_{m'} \in B_{ijk}$。
此外，约定当 $k \in I_{ij}$ 时，对所有
$i, j = i, k, m$ 都有 $S(i, i, k, m) = \{m\}$。
固定这样的一族 $S(i, j, k, m)$。

\medskip\noindent
给定一族子集 $\mathcal{S} = (S_i)_{i \in I}$，
其中 $S_i \subset M_i$ 的基数至多为 $\kappa$，定义
$\mathcal{S}' = (S'_i)$，其中
$$
S'_j = \bigcup\nolimits_{(i, j, k, m)\text{ 使得 }m \in S_i}
S(i, j, k, m)
$$
注意 $S_i \subset S'_i$。由于 $S'_i$ 是在一个基数至多为
$\kappa$ 的集合上对有限集取并，其基数至多为 $\kappa$。
令 $\mathcal{S}^{(0)} = \mathcal{S}$、
$\mathcal{S}^{(1)} = \mathcal{S}'$，并归纳定义
$\mathcal{S}^{(n + 1)} = (\mathcal{S}^{(n)})'$。再令
$\mathcal{S}^{(\infty)} = \bigcup_{n \geq 0} \mathcal{S}^{(n)}$。
写作 $\mathcal{S}^{(\infty)} = (S^{(\infty)}_i)$，则对任意元素
$m \in S^{(\infty)}_i$，$m$ 在 $M_{ijk}$ 中的像可写成有限和
$\sum m' \otimes a_{m'}$，其中 $m' \in S_j^{(\infty)}$。
因此，令
$$
N_i = A_i\text{-子模，属于 }M_i\text{ 由下列对象生成： }S^{(\infty)}_i
$$
便有
$$
N_i \otimes_{A_i} A_{ijk} = N_j \otimes_{A_j} A_{ijk}
\quad\text{且}\quad
\alpha(N_i \otimes_{A_i, t} B_{ijk}) = N_j \otimes_{A_j, s} B_{ijk}
$$
这些等式分别作为 $M_{ijk}$ 或
$M_j \otimes_{A_j, s} B_{ijk}$ 的子模等式成立。
故存在拟凝聚子模
$\mathcal{G} \subset \mathcal{F}$，满足
$\mathcal{G}(U_i) = N_i$，且作为 $s^*\mathcal{F}$ 的子模有
$\alpha(t^*\mathcal{G}) = s^*\mathcal{G}$。
换言之，$(\mathcal{G}, \alpha|_{t^*\mathcal{G}})$ 是
$(\mathcal{F}, \alpha)$ 的子模。此外，依构造
$\mathcal{G}$ 是 $\kappa$-生成的。

\medskip\noindent
设 $\{(\mathcal{G}_t, \alpha_t)\}_{t \in T}$ 为
$(\mathcal{F}, \alpha)$ 的全体 $\kappa$-生成拟凝聚子模之集合。
若 $t, t' \in T$，则 $\mathcal{G}_t + \mathcal{G}_{t'}$
也是 $\kappa$-生成拟凝聚子模，因为它是态射
$\mathcal{G}_t \oplus \mathcal{G}_{t'} \to \mathcal{F}$ 的像。
因此，该系统按包含关系排序后是有向的。上述论证表明，
$\mathcal{F}$ 在 $U_i$ 上的每个截面都属于某个
$\mathcal{G}_t$（因为可从使该截面属于 $S_i$ 的
$\mathcal{S}$ 开始）。故
$\colim_t \mathcal{G}_t \to \mathcal{F}$ 既单又满，如所需。
\end{proof}





\section{群胚与群概形}
\label{groupoids-section-groupoids-group-schemes}

\noindent
由群 $G$ 在集合 $V$ 上的作用 $a$ 构造群胚有许多方式。
我们采用如下方式：把元素 $g \in G$ 看作一条以 $v$ 为源、
以 $a(g, v)$ 为靶的箭头。这就导出概形之群作用的下述构造。

\begin{lemma}
\label{groupoids-lemma-groupoid-from-action}
设 $S$ 为概形，$Y$ 为 $S$ 上的概形，$(G, m)$ 为 $Y$ 上的群概形，
其单位元为 $e_G$、逆为 $i_G$。设 $X/Y$ 为 $Y$ 上的概形，
$a : G \times_Y X \to X$ 是 $G$ 在 $X/Y$ 上的作用。
则可按如下方式得到 $S$ 上的群胚概形
$(U, R, s, t, c, e, i)$：
\begin{enumerate}
\item 令 $U = X$，$R = G \times_Y X$。
\item 令 $s : R \to U$ 为 $(g, x) \mapsto x$。
\item 令 $t : R \to U$ 为 $(g, x) \mapsto a(g, x)$。
\item 令 $c : R \times_{s, U, t} R \to R$ 为
$((g, x), (g', x')) \mapsto (m(g, g'), x')$。
\item 令 $e : U \to R$ 为 $x \mapsto (e_G(x), x)$。
\item 令 $i : R \to R$ 为 $(g, x) \mapsto (i_G(g), a(g, x))$。
\end{enumerate}
\end{lemma}

\begin{proof}
证明略。提示：只需在集合层面验证这一构造即可。
为此，使用引理上方的描述，把 $g$ 看成从 $v$ 到
$a(g, v)$ 的箭头。
\end{proof}

\begin{lemma}
\label{groupoids-lemma-action-groupoid-modules}
设 $S$ 为概形，$Y$ 为 $S$ 上的概形，$(G, m)$ 为 $Y$ 上的群概形。
设 $X$ 为 $Y$ 上的概形，$a : G \times_Y X \to X$
是 $G$ 在 $X$ 上且相对于 $Y$ 的作用。令 $(U, R, s, t, c)$ 为
引理 \ref{groupoids-lemma-groupoid-from-action} 构造的群胚概形。
规则 $(\mathcal{F}, \alpha) \mapsto (\mathcal{F}, \alpha)$
定义了 $G$-等变 $\mathcal{O}_X$-模范畴与
$(U, R, s, t, c)$ 上拟凝聚模范畴之间的范畴等价。
\end{lemma}

\begin{proof}
该断言是有意义的，因为作为态射
$R = G \times_Y X \to X$，有
$t = a$ 以及 $s = \text{pr}_1$；见
定义 \ref{groupoids-definition-equivariant-module} 和
\ref{groupoids-definition-groupoid-module}。
采用引理 \ref{groupoids-lemma-groupoid-from-action} 中的转译，
两个定义对交换性的要求完全一致。
\end{proof}





\section{稳定子群概形}
\label{groupoids-section-stabilizer}

\noindent
给定群胚概形，可如下得到一个群概形。

\begin{lemma}
\label{groupoids-lemma-groupoid-stabilizer}
设 $S$ 为概形，$(U, R, s, t, c)$ 为 $S$ 上的群胚。
由笛卡尔方块
$$
\xymatrix{
G \ar[r] \ar[d] & R \ar[d]^{j = (t, s)} \\
U \ar[r]^-{\Delta} & U \times_S U
}
$$
定义的概形 $G$ 是 $U$ 上的群概形，其乘法
$m$ 由复合律 $c$ 诱导。
\end{lemma}

\begin{proof}
这是因为在群胚范畴中，任一对象的自态射集构成一个群。
\end{proof}

\noindent
由于 $\Delta$ 是浸入，$G = j^{-1}(\Delta_{U/S})$
是 $R$ 的局部闭子概形。从这个角度看，结构态射
$j^{-1}(\Delta_{U/S}) \to U$ 可由 $s$ 或 $t$ 诱导
（二者相同），而 $m$ 由 $c$ 诱导。

\begin{definition}
\label{groupoids-definition-stabilizer-groupoid}
设 $S$ 为概形，$(U, R, s, t, c)$ 为 $S$ 上的群胚。
群概形 $j^{-1}(\Delta_{U/S})\to U$ 称为
群胚概形 $(U, R, s, t, c)$ 的{\it 稳定子}。
\end{definition}

\noindent
文献中常把稳定子群概形记作 $S$
（大概是因为 stabilizer 一词以字母 s 开头）；
这里不能这样做，因为我们已经用 $S$ 表示基概形。

\begin{lemma}
\label{groupoids-lemma-groupoid-action-stabilizer}
设 $S$ 为概形，$(U, R, s, t, c)$ 为 $S$ 上的群胚，
$G/U$ 为其稳定子。把 $R$ 经态射 $t : R \to U$
视为 $U$ 上的概形，记作 $R_t/U$。
存在由复合律 $c$ 诱导的典范左作用
$$
a : G \times_U R_t \longrightarrow R_t
$$
\end{lemma}

\begin{proof}
就 $T/S$ 上的点而言，定义 $a(g, r) = c(g, r)$。
\end{proof}

\begin{lemma}
\label{groupoids-lemma-groupoid-action-stabilizer-pseudo-torsor}
设 $S$ 为概形，$(U, R, s, t, c)$ 为 $S$ 上的群胚概形，
$G$ 为 $R$ 的稳定子群概形。令
$$
G_0 = G \times_{U, \text{pr}_0} (U \times_S U) = G \times_S U
$$
并把它视为 $U \times_S U$ 上的群概形。引理 \ref{groupoids-lemma-groupoid-action-stabilizer} 中 $G$ 在 $R$ 上的作用
诱导 $G_0$ 在 $R$ 上相对于 $U \times_S U$ 的作用，使 $R$
成为 $U \times_S U$ 上的伪 $G_0$-扭子。
\end{lemma}

\begin{proof}
这是因为在群胚范畴 $\mathcal{C}$ 中，集合
$\Mor_\mathcal{C}(x, y)$ 是群
$\Mor_\mathcal{C}(y, y)$ 下的主齐性集。
\end{proof}

\begin{lemma}
\label{groupoids-lemma-fibres-j}
设 $S$ 为概形，$(U, R, s, t, c)$ 为 $S$ 上的群胚概形。
设 $p \in U \times_S U$ 为一点，并以 $R_p$ 表示
$j = (t, s) : R \to U \times_S U$ 的概形论纤维。
若 $R_p \not = \emptyset$，则作用
$$
G_{0, \kappa(p)} \times_{\kappa(p)} R_p \longrightarrow R_p
$$
（见引理 \ref{groupoids-lemma-groupoid-action-stabilizer-pseudo-torsor}）
使 $R_p$ 成为 $\kappa(p)$ 上的 $G_{\kappa(p)}$-扭子。
\end{lemma}

\begin{proof}
由陈述中引用的引理，该作用是伪扭子。若 $R_p$ 不是空概形，
则 $\{R_p \to p\}$ 是使该伪扭子平凡化的 fpqc 覆盖。
\end{proof}







\section{群胚的限制}
\label{groupoids-section-restrict-groupoid}

\noindent
考虑一个（通常意义下的）群胚
$\mathcal{C} = (\text{Ob}, \text{Arrows}, s, t, c)$，
并假设给定集合映射 $g : \text{Ob}' \to \text{Ob}$。
可以构造群胚
$\mathcal{C}' = (\text{Ob}', \text{Arrows}', s', t', c')$：
把 $\text{Ob}'$ 中元素 $x', y'$ 之间的态射理解为
$\mathcal{C}$ 中 $g(x'), g(y')$ 之间的态射。换言之，令
$$
\text{Arrows}' =
\text{Ob}'
\times_{g, \text{Ob}, t}
\text{Arrows}
\times_{s, \text{Ob}, g}
\text{Ob}'.
$$
并以显然的方式选择 $s'$、$t'$ 和 $c'$。存在典范函子
$\mathcal{C}' \to \mathcal{C}$，它全忠实，但未必本质满。
配备函子 $\mathcal{C}' \to \mathcal{C}$ 的群胚
$\mathcal{C}'$ 称为群胚 $\mathcal{C}$ 在 $\text{Ob}'$ 上的
{\it 限制}。

\begin{lemma}
\label{groupoids-lemma-restrict-groupoid}
设 $S$ 为概形，$(U, R, s, t, c)$ 为 $S$ 上的群胚概形，
$g : U' \to U$ 为概形态射。考虑下图
$$
\xymatrix{
R' \ar[d] \ar[r] \ar@/_3pc/[dd]_{t'} \ar@/^1pc/[rr]^{s'}&
R \times_{s, U} U' \ar[r] \ar[d] &
U' \ar[d]^g \\
U' \times_{U, t} R \ar[d] \ar[r] &
R \ar[r]^s \ar[d]_t &
U \\
U' \ar[r]^g &
U
}
$$
其中所有方块都是纤维积方块。则存在典范复合律
$c' : R' \times_{s', U', t'} R' \to R'$，使
$(U', R', s', t', c')$ 成为 $S$ 上的群胚概形，并使
$U' \to U$, $R' \to R$ 定义 $S$ 上群胚概形的态射
$(U', R', s', t', c') \to (U, R, s, t, c)$。
此外，对 $S$ 上任意概形 $T$，群胚函子
$$
(U'(T), R'(T), s', t', c') \to (U(T), R(T), s, t, c)
$$
是 $(U(T), R(T), s, t, c)$ 经映射
$U'(T) \to U(T)$ 的限制（如上所述）。
\end{lemma}

\begin{proof}
证明略。
\end{proof}

\begin{definition}
\label{groupoids-definition-restrict-groupoid}
设 $S$ 为概形，$(U, R, s, t, c)$ 为 $S$ 上的群胚概形，
$g : U' \to U$ 为概形态射。引理 \ref{groupoids-lemma-restrict-groupoid} 中构造的群胚态射
$(U', R', s', t', c') \to (U, R, s, t, c)$
称为 $(U, R, s, t, c)$ 在 $U'$ 上的{\it 限制}。
在这种情形下，有时使用记号 $R' = R|_{U'}$。
\end{definition}

\begin{lemma}
\label{groupoids-lemma-restrict-groupoid-relation}
定义 \ref{groupoids-definition-restrict-groupoid} 与
\ref{groupoids-definition-restrict-relation} 所定义的群胚限制和
（预）等价关系限制，经引理 \ref{groupoids-lemma-groupoid-pre-equivalence} 与
\ref{groupoids-lemma-equivalence-groupoid} 的构造彼此一致。
\end{lemma}

\begin{proof}
这里所断言的是，引理 \ref{groupoids-lemma-restrict-groupoid} 中的
$R'$ 也等于
$$
R' = (U' \times_S U')\times_{U \times_S U} R
\longrightarrow
U' \times_S U'
$$
事实上，这也许是陈述该引理更清楚的方式。
\end{proof}

\begin{lemma}
\label{groupoids-lemma-restrict-stabilizer}
设 $S$ 为概形，$(U, R, s, t, c)$ 为 $S$ 上的群胚概形，
$g : U' \to U$ 为概形态射。令
$(U', R', s', t', c')$ 为 $(U, R, s, t, c)$ 经 $g$ 的限制。
设 $G$ 为 $(U, R, s, t, c)$ 的稳定子，$G'$ 为
$(U', R', s', t', c')$ 的稳定子。则 $G'$ 是 $G$ 沿 $g$ 的基变换；
换言之，存在典范识别 $G' = U' \times_{g, U} G$。
\end{lemma}

\begin{proof}
证明略。
\end{proof}






\section{不变子概形}
\label{groupoids-section-invariant}

\noindent
本节简要讨论不变子概形的概念。

\begin{definition}
\label{groupoids-definition-invariant-open}
设 $(U, R, s, t, c)$ 为基概形 $S$ 上的群胚概形。
\begin{enumerate}
\item 若 $t(s^{-1}(W)) \subset W$，则称子集
$W \subset U$ 在{\it 集合论意义下是 $R$-不变的}。
\item 若 $t(s^{-1}(W)) \subset W$，则称开子集
$W \subset U$ 是{\it $R$-不变的}。
\item 若 $t^{-1}(Z) = s^{-1}(Z)$，则称闭子概形
$Z \subset U$ 是{\it $R$-不变的}。这里采用概形论逆像；见
《概形》定义 \ref{schemes-definition-inverse-image-closed-subscheme}。
\item 若作为 $R$ 上的概形有
$T \times_{U, t} R = R \times_{s, U} T$，则称概形单态射
$T \to U$ 是{\it $R$-不变的}。
\end{enumerate}
\end{definition}

\noindent
对于子集与开子概形 $W \subset U$，$R$-不变性也等价于要求
作为 $R$ 的子集有 $s^{-1}(W) = t^{-1}(W)$。
若 $W \subset U$ 是 $R$-等变开子概形，则 $R$ 在 $W$ 上的限制就是
$R_W = s^{-1}(W) = t^{-1}(W)$。类似地，若
$Z \subset U$ 是 $R$-不变闭子概形，则 $R$ 在 $Z$ 上的限制就是
$R_Z = s^{-1}(Z) = t^{-1}(Z)$。

\begin{lemma}
\label{groupoids-lemma-constructing-invariant-opens}
设 $S$ 为概形，$(U, R, s, t, c)$ 为 $S$ 上的群胚概形。
\begin{enumerate}
\item 对任意子集 $W \subset U$，子集 $t(s^{-1}(W))$
在集合论意义下是 $R$-不变的。
\item 若 $s$ 与 $t$ 是开态射，则对每个开子集 $W \subset U$，
开集 $t(s^{-1}(W))$ 是 $R$-不变开子概形。
\item 若 $s$ 与 $t$ 开且拟紧，则 $U$ 有一个由
$R$-不变拟紧开子概形组成的开覆盖。
\end{enumerate}
\end{lemma}

\begin{proof}
第 (1) 点由引理 \ref{groupoids-lemma-pre-equivalence-equivalence-relation-points} 和
\ref{groupoids-lemma-groupoid-pre-equivalence} 得出：$t(s^{-1}(W))$
正是 $U$ 中与 $W$ 的某一点等价的点集。
接着假设 $s$ 和 $t$ 开，且 $W \subset U$ 开。
由于 $s$ 开，集合 $W' = t(s^{-1}(W))$ 是 $U$ 的开子集。
最后，假设 $s$, $t$ 都开且拟紧。若 $W \subset U$ 拟紧开，
则 $W' = t(s^{-1}(W))$ 也拟紧开，并由上面的讨论可知它不变。
让 $W$ 遍历 $U$ 的全体仿射开集，即得 (3)。
\end{proof}

\begin{lemma}
\label{groupoids-lemma-first-observation}
设 $S$ 为概形，$(U, R, s, t, c)$ 为 $S$ 上的群胚概形。
假设 $s$ 与 $t$ 拟紧且平坦，而 $U$ 拟分离。若
$W \subset U$ 拟紧开，则 $t(s^{-1}(W))$ 是
$U$ 中某个非空族拟紧开子集的交。
\end{lemma}

\begin{proof}
注意 $s^{-1}(W)$ 是 $R$ 中的拟紧开集。作为连续映射，
$t$ 把拟紧子集 $s^{-1}(W)$ 映到拟紧子集 $t(s^{-1}(W))$。
由于 $t$ 平坦且 $s^{-1}(W)$ 对泛化封闭，$t(s^{-1}(W))$
也对泛化封闭；见
（《概形态射》引理 \ref{morphisms-lemma-generalizations-lift-flat} 及
《拓扑学》引理 \ref{topology-lemma-lift-specializations-images}）。
选取包含 $t(s^{-1}(W))$ 的拟紧开集 $W' \subset U$。
由《概形的性质》引理 \ref{properties-lemma-quasi-compact-quasi-separated-spectral}，
$W'$ 是谱空间（这里用到了 $U$ 拟分离）。对
$t(s^{-1}(W)) \subset W'$ 应用《拓扑学》引理 \ref{topology-lemma-make-spectral-space}，即得结论。
\end{proof}

\begin{lemma}
\label{groupoids-lemma-second-observation}
假设与记号同引理 \ref{groupoids-lemma-first-observation}。
存在 $R$-不变开集 $V \subset U$ 以及拟紧开集 $W'$，
使 $W \subset V \subset W' \subset U$。
\end{lemma}

\begin{proof}
令 $E = t(s^{-1}(W))$。回忆 $E$ 在集合论意义下是 $R$-不变的
（引理 \ref{groupoids-lemma-constructing-invariant-opens}）。
由引理 \ref{groupoids-lemma-first-observation}，存在包含 $E$ 的拟紧开集
$W'$。令 $Z = U \setminus W'$，并考虑
$T = t(s^{-1}(Z))$。注意 $Z \subset T$，且
$E \cap T = \emptyset$，因为
$s^{-1}(E) = t^{-1}(E)$ 与 $s^{-1}(Z)$ 不交。
由于 $T$ 是闭子集 $s^{-1}(Z) \subset R$ 在拟紧态射
$t : R \to U$ 下的像，闭包 $\overline{T}$ 中任意点
$\xi$ 都是 $T$ 中某点的特化；见
《概形态射》引理 \ref{morphisms-lemma-reach-points-scheme-theoretic-image}
（并见《概形态射》引理 \ref{morphisms-lemma-quasi-compact-scheme-theoretic-image}，
它说明概形论像就是像的闭包）。
设 $\xi' \leadsto \xi$，其中 $\xi' \in T$。若 $r \in R$ 且
$s(r) = \xi$，则由 $s$ 平坦，可在 $R$ 中找到一个特化
$r' \leadsto r$，使 $s(r') = \xi'$
（《概形态射》引理 \ref{morphisms-lemma-generalizations-lift-flat}）。
于是 $t(r') \leadsto t(r)$。由于 $T$ 由
引理 \ref{groupoids-lemma-constructing-invariant-opens}
在集合论意义下不变，故 $t(r') \in T$。
因此 $\overline{T}$ 是集合论意义下 $R$-不变的闭子集，
而 $V = U \setminus \overline{T}$ 正是所求开集。
它包含于 $W'$，证明完成。
\end{proof}




\section{商层}
\label{groupoids-section-quotient-sheaves}

\noindent
设 $\tau \in \{Zariski, \etale, fppf, smooth, syntomic\}$，
$S$ 为概形，$j : R \to U \times_S U$ 为 $S$ 上的预关系。
假设 $U, R, S$ 是某个 $\tau$-位点 $\Sch_\tau$ 的对象
（见《概形上的拓扑》第 \ref{topologies-section-procedure} 节）。
于是可考虑函子
$$
h_U, h_R :
(\Sch/S)_\tau^{opp}
\longrightarrow
\textit{Sets}.
$$
它们都是层，见《下降》引理 \ref{descent-lemma-fpqc-universal-effective-epimorphisms}。
态射 $j$ 诱导映射 $j : h_R \to h_U \times h_U$。
对每个对象 $T \in \Ob((\Sch/S)_\tau)$，
可取由 $j(T) : R(T) \to U(T) \times U(T)$ 生成的等价关系
$\sim_T$，并取其商。这样得到预层
\begin{equation}
\label{groupoids-equation-quotient-presheaf}
(\Sch/S)_\tau^{opp}
\longrightarrow
\textit{Sets}, \quad
T \longmapsto U(T)/\sim_T
\end{equation}

\begin{definition}
\label{groupoids-definition-quotient-sheaf}
令 $\tau$, $S$ 及预关系 $j : R \to U \times_S U$ 如上。
在此情形下，与 $j$ 关联的{\it 商层 $U/R$}，是预层
(\ref{groupoids-equation-quotient-presheaf}) 在 $\tau$-拓扑下的层化。
若 $j : R \to U \times_S U$ 来自群概形 $G/S$
在 $U$ 上的作用，如引理 \ref{groupoids-lemma-groupoid-from-action}，
则有时把该商层记作 $U/G$。
\end{definition}

\noindent
这恰好意味着图
$$
\xymatrix{
h_R \ar@<1ex>[r] \ar@<-1ex>[r] &
h_U \ar[r] &
U/R
}
$$
是 $(\Sch/S)_\tau$ 上集合层范畴中的余等化子图。
利用 Yoneda 嵌入，可以把 $(\Sch/S)_\tau$
看成 $(\Sch/S)_\tau$ 上层范畴的全子范畴，从而把概形与可表函子相识别。
沿用这一记号滥用，通常把上图简写为
$$
\xymatrix{
R \ar@<1ex>[r]^s \ar@<-1ex>[r]_t &
U \ar[r] &
U/R
}
$$
在考虑群胚或等价关系的商层时，我们主要使用 fppf 拓扑。

\begin{definition}
\label{groupoids-definition-representable-quotient}
在定义 \ref{groupoids-definition-quotient-sheaf} 的情形下，
若层 $U/R$ 可表，则称预关系 $j$ 有{\it 可表商}。
若取 $j = (t, s)$ 时商 $U/R$ 可表，则称群胚
$(U, R, s, t, c)$ 有{\it 可表商}。
\end{definition}

\noindent
下述引理刻画表示该商的概形 $M$。例如，当
$\tau = fppf$、$U \to M$ 平坦、有限表现且满，并且
$R \cong U \times_M U$ 时，该引理适用。

\begin{lemma}
\label{groupoids-lemma-criterion-quotient-representable}
在定义 \ref{groupoids-definition-quotient-sheaf} 的情形下，
假设存在概形 $M$ 以及态射 $U \to M$，满足：
\begin{enumerate}
\item 态射 $U \to M$ 等化 $s, t$；
\item 态射 $U \to M$ 在 $\tau$-拓扑下诱导层的满射
$h_U \to h_M$；
\item 诱导映射 $(t, s) : R \to U \times_M U$
在 $\tau$-拓扑下诱导层的满射
$h_R \to h_{U \times_M U}$。
\end{enumerate}
此时 $M$ 表示商层 $U/R$。
\end{lemma}

\begin{proof}
条件 (1) 说明 $h_U \to h_M$ 通过 $U/R$ 分解；
条件 (2) 说明层映射 $U/R \to h_M$ 是满射；
条件 (3) 说明层映射 $U/R \to h_M$ 是单射。
故结论成立。
\end{proof}

\noindent
若不要求 $j$ 为预等价关系（例如仅为预关系），则下述引理不成立。

\begin{lemma}
\label{groupoids-lemma-quotient-pre-equivalence}
设 $\tau \in \{Zariski, \etale, fppf, smooth, syntomic\}$，
$S$ 为概形，$j : R \to U \times_S U$ 为 $S$ 上的预等价关系，
并假设 $U, R, S$ 是某个 $\tau$-位点 $\Sch_\tau$ 的对象。
对 $T \in \Ob((\Sch/S)_\tau)$ 及
$a, b \in U(T)$，下列条件等价：
\begin{enumerate}
\item $a$ 与 $b$ 映到 $(U/R)(T)$ 中的同一元素；
\item 存在 $T$ 的一个 $\tau$-覆盖
$\{f_i : T_i \to T\}$ 以及态射 $r_i : T_i \to R$，使
$a \circ f_i = s \circ r_i$ 且 $b \circ f_i = t \circ r_i$。
\end{enumerate}
换言之，此时 $\tau$-层的映射
$$
h_R \longrightarrow h_U \times_{U/R} h_U
$$
是满射。
\end{lemma}

\begin{proof}
证明略。提示：该论证之所以成立，是因为此时预层
(\ref{groupoids-equation-quotient-presheaf}) 实际由
$T \mapsto U(T)/j(R(T))$ 给出，而
$j(R(T)) \subset U(T) \times U(T)$ 是等价关系；见
定义 \ref{groupoids-definition-equivalence-relation}。
\end{proof}

\begin{lemma}
\label{groupoids-lemma-quotient-pre-equivalence-relation-restrict}
设 $\tau \in \{Zariski, \etale, fppf, smooth, syntomic\}$，
$S$ 为概形，$j : R \to U \times_S U$ 为 $S$ 上的预等价关系，
$g : U' \to U$ 为 $S$ 上的概形态射。令
$j' : R' \to U' \times_S U'$ 为 $j$ 在 $U'$ 上的限制。
假设 $U, U', R, S$ 是某个 $\tau$-位点 $\Sch_\tau$ 的对象。
则商层的映射
$$
U'/R' \longrightarrow U/R
$$
是单射。若 $g$ 定义层的满射 $h_{U'} \to h_U$，
其拓扑为 $\tau$（例如，若 $\{g : U' \to U\}$
是 $\tau$-覆盖），则 $U'/R' \to U/R$ 是同构。
\end{lemma}

\begin{proof}
假设截面 $\xi, \xi' \in (U'/R')(T)$ 映到 $U/R$ 的同一截面。
可以找到 $T$ 的一个 $\tau$-覆盖
$\mathcal{T} = \{T_i \to T\}$，使
$\xi|_{T_i}, \xi'|_{T_i}$ 分别由
$a_i, a_i' \in U'(T_i)$ 给出。由
引理 \ref{groupoids-lemma-quotient-pre-equivalence} 及位点公理，
加细 $\mathcal{T}$ 后可假设存在态射 $r_i : T_i \to R$，使
$g \circ a_i = s \circ r_i$、$g \circ a_i' = t \circ r_i$。
按照构造，
$R' = R \times_{U \times_S U} (U' \times_S U')$，
故 $(r_i, (a_i, a_i')) \in R'(T_i)$。这说明 $a_i$ 与
$a_i'$ 在 $T_i$ 上定义 $U'/R'$ 的同一截面。
由层条件可知 $\xi = \xi'$。

\medskip\noindent
若 $h_{U'} \to h_U$ 是层的满射，则
$U'/R' \to U/R$ 当然也是满射。若
$\{g : U' \to U\}$ 是 $\tau$-覆盖，则层映射
$h_{U'} \to h_U$ 是满射，见《位点与层》引理 \ref{sites-lemma-covering-surjective-after-sheafification}。
因此在此情形下 $U'/R' \to U/R$ 也是满射。
\end{proof}

\begin{lemma}
\label{groupoids-lemma-quotient-groupoid-restrict}
设 $\tau \in \{Zariski, \etale, fppf, smooth, syntomic\}$，
$S$ 为概形，$(U, R, s, t, c)$ 为 $S$ 上的群胚概形，
$g : U' \to U$ 为 $S$ 上的概形态射。令
$(U', R', s', t', c')$ 为 $(U, R, s, t, c)$ 在 $U'$ 上的限制。
假设 $U, U', R, S$ 是某个 $\tau$-位点 $\Sch_\tau$ 的对象。
则商层的映射
$$
U'/R' \longrightarrow U/R
$$
是单射。若复合
$$
\xymatrix{
U' \times_{g, U, t} R \ar[r]_-{\text{pr}_1} \ar@/^3ex/[rr]^h
& R \ar[r]_s & U
}
$$
在 $\tau$-拓扑下定义层的满射，则该映射是双射。
例如，下列任一条件均足够：
$\{h : U' \times_{g, U, t} R \to U\}$ 是 $\tau$-覆盖；
$U' \to U$ 在 $\tau$-拓扑下定义层的满射；或
$\{g : U' \to U\}$ 是 $\tau$-拓扑下的覆盖。
\end{lemma}

\begin{proof}
由引理 \ref{groupoids-lemma-groupoid-pre-equivalence} 和
\ref{groupoids-lemma-quotient-pre-equivalence-relation-restrict} 即得单射性。
为证明满射性（层满射的刻画见
《位点与层》第 \ref{sites-section-sheaves-injective} 节），作如下论证。
假设 $T$ 是概形，$\sigma \in U/R(T)$。存在覆盖
$\{T_i \to T\}$，使 $\sigma|_{T_i}$ 是某个元素
$f_i \in U(T_i)$ 的像。因此可假设 $\sigma$ 是
$f \in U(T)$ 的像。由 $h$ 是层满射这一假设，可以找到
$\tau$-覆盖 $\{\varphi_i : T_i \to T\}$ 以及态射
$f_i : T_i \to U' \times_{g, U, t} R$，使
$f \circ \varphi_i = h \circ f_i$。记
$f'_i = \text{pr}_0 \circ f_i : T_i \to U'$。
于是 $f'_i \in U'(T_i)$ 映到
$g \circ f'_i \in U(T_i)$，并且
$g \circ f'_i \sim_{T_i} h \circ f_i = f \circ \varphi_i$；
记号同 (\ref{groupoids-equation-quotient-presheaf})。具体地，
给出该关系的 $R(T_i)$ 中元素是 $\text{pr}_1 \circ f_i$。
这说明 $\sigma$ 在 $T_i$ 上的限制属于
$U'/R'(T_i) \to U/R(T_i)$ 的像，如所需。

\medskip\noindent
若 $\{h\}$ 是 $\tau$-覆盖，则它诱导层的满射，见
《位点与层》引理 \ref{sites-lemma-covering-surjective-after-sheafification}。
若 $U' \to U$ 满，则 $h$ 也满，因为 $s$ 有截面
（即群胚概形的单位元 $e$）。
\end{proof}

\begin{lemma}
\label{groupoids-lemma-criterion-fibre-product}
设 $S$ 为概形，$f : (U, R, j) \to (U', R', j')$
为 $S$ 上等价关系之间的态射。假设
$$
\xymatrix{
R \ar[d]_s \ar[r]_f & R' \ar[d]^{s'} \\
U \ar[r]^f & U'
}
$$
是笛卡尔的。对任意
$\tau \in  \{Zariski, \etale, fppf, smooth, syntomic\}$，图
$$
\xymatrix{
U \ar[d] \ar[r] & U/R \ar[d]^f \\
U' \ar[r] & U'/R'
}
$$
是 $\tau$-层的纤维积方块。
\end{lemma}

\begin{proof}
由引理 \ref{groupoids-lemma-quotient-pre-equivalence}，商层有一个简单描述；
下文将不再说明而直接使用它。先证明
$$
U \longrightarrow U' \times_{U'/R'} U/R
$$
是单射。假设 $a, b \in U(T)$ 映到右侧同一元素，则
$f(a) = f(b)$。把 $T$ 替换为某个 $\tau$-覆盖的各成员后，
可假设存在 $r \in R(T)$，使 $a = s(r)$ 且 $b = t(r)$。
于是 $r' = f(r)$ 是 $R'$ 的一个 $T$-值点，满足
$s'(r') = t'(r')$。因此 $r' = e'(f(a))$（其中 $e'$
是与 $j'$ 关联的群胚概形的单位，见
引理 \ref{groupoids-lemma-equivalence-groupoid}）。
由于引理中的第一个图是笛卡尔的，这蕴含 $r$ 必须等于
$e(a)$，故 $a = b$。

\medskip\noindent
最后证明所示箭头满。设 $T$ 为 $S$ 上的概形，
$(a', \overline{b})$ 为层
$U' \times_{U'/R'} U/R$ 在 $T$ 上的截面。
把 $T$ 替换为某个 $\tau$-覆盖的各成员后，可假设
$\overline{b}$ 是元素 $b \in U(T)$ 的类。
再次把 $T$ 替换为某个 $\tau$-覆盖的各成员后，可假设存在
$r' \in R'(T)$，使 $a' = t(r')$ 且 $s'(r') = f(b)$。
由于引理中的第一个图是笛卡尔的，可找到 $r \in R(T)$，
使 $s(r) = b$ 且 $f(r) = r'$。显然，
$a = t(r) \in U(T)$ 是映到 $(a', \overline{b})$ 的截面。
\end{proof}








\section{用群胚表述下降}
\label{groupoids-section-groupoids-descent}

\noindent
笛卡尔态射定义如下。

\begin{definition}
\label{groupoids-definition-cartesian-morphism}
设 $S$ 为概形，$f : (U', R', s', t', c') \to (U, R, s, t, c)$
为 $S$ 上群胚概形的态射。若图
$$
\xymatrix{
R' \ar[r]_f \ar[d]_{s'} & R \ar[d]^s \\
U' \ar[r]^f & U
}
$$
是概形范畴中的纤维积方块，则称 $f$ 是{\it 笛卡尔的}，
或称{\it $(U', R', s', t', c')$ 在 $(U, R, s, t, c)$ 上是笛卡尔的}。
所谓{\it 在 $(U, R, s, t, c)$ 上笛卡尔的群胚概形态射}，
是指与指向 $(U, R, s, t, c)$ 的结构态射相容的群胚概形态射。
\end{definition}

\noindent
笛卡尔态射与下降数据相关。先证明一个一般引理，
描述固定群胚概形上的笛卡尔群胚概形范畴。

\begin{lemma}
\label{groupoids-lemma-characterize-cartesian-schemes}
设 $S$ 为概形，$(U, R, s, t, c)$ 为 $S$ 上的群胚概形。
在 $(U, R, s, t, c)$ 上笛卡尔的群胚概形范畴等价于
如下偶对 $(V, \varphi)$ 的范畴：$V$ 是 $U$ 上的概形，而
$$
\varphi :
V \times_{U, t} R
\longrightarrow
R \times_{s, U} V
$$
是 $R$ 上的同构，满足 $e^*\varphi = \text{id}_V$，并且作为
$R \times_{s, U, t} R$ 上的概形态射有
$$
c^*\varphi = \text{pr}_1^*\varphi \circ \text{pr}_0^*\varphi
$$
\end{lemma}

\begin{proof}
引理中的拉回记号表示基变换。所示公式有意义，因为作为
$R \times_{s, U, t} R$ 上的概形有
$$
(R \times_{s, U, t} R) \times_{\text{pr}_1, R, \text{pr}_1} (V \times_{U, t} R)
=
(R \times_{s, U, t} R) \times_{\text{pr}_0, R, \text{pr}_0} (R \times_{s, U} V)
$$

\medskip\noindent
给定 $(V, \varphi)$，令 $U' = V$、$R' = V \times_{U, t} R$。
令 $t' : R' \to U'$ 为投影 $V \times_{U, t} R \to V$；
令 $s'$ 为 $\varphi$ 后接投影 $R \times_{s, U} V \to V$；
令 $c'$ 为复合
\begin{align*}
R' \times_{s', U', t'} R'
& \xrightarrow{\varphi, 1}
(R \times_{s, U} V) \times_V (V \times_{U, t} R) \\
& \xrightarrow{}
R \times_{s, U} V \times_{U, t} R \\
& \xrightarrow{\varphi^{-1}, 1}
V \times_{U, t} (R \times_{s, U, t} R) \\
& \xrightarrow{1, c}
V \times_{U, t} R = R'
\end{align*}
略去的计算表明，这给出 $(U, R, s, t, c)$ 上的群胚概形；
显然，该群胚概形在 $(U, R, s, t, c)$ 上是笛卡尔的。

\medskip\noindent
反之，给定笛卡尔态射
$f : (U', R', s', t', c') \to (U, R, s, t, c)$，
则态射
$$
U' \times_{U, t} R \xleftarrow{t', f} R' \xrightarrow{f, s'} R \times_{s, U} U'
$$
都是同构。可以令 $V = U'$，并令 $\varphi$ 为复合
$(f, s') \circ (t', f)^{-1}$。我们略去验证 $\varphi$
满足引理中的条件，也略去验证这两个构造互逆。
\end{proof}

\noindent
设 $S$ 为概形，$f : X \to Y$ 为 $S$ 上的概形态射。
于是得到 $S$ 上的群胚概形
$(X, X \times_Y X, \text{pr}_1, \text{pr}_0, c)$。
事实上，$j : X \times_Y X \to X \times_S X$ 是等价关系，
可取其关联群胚；见引理 \ref{groupoids-lemma-equivalence-groupoid}。

\begin{lemma}
\label{groupoids-lemma-cartesian-equivalent-descent-datum}
设 $S$ 为概形，$f : X \to Y$ 为 $S$ 上的概形态射。
引理 \ref{groupoids-lemma-characterize-cartesian-schemes} 的构造给出等价
$$
\begin{matrix}
\text{群胚概形范畴} \\
\text{在下列对象上 Cartesian： } (X, X \times_Y X, \ldots)
\end{matrix}
\longrightarrow
\begin{matrix}
\text{ 下降数据范畴} \\
\text{ 相对于 } X/Y
\end{matrix}
$$
\end{lemma}

\begin{proof}
这直接源于引理 \ref{groupoids-lemma-characterize-cartesian-schemes}
以及《下降》定义 \ref{descent-definition-descent-datum}
中概形上下降数据的定义。
\end{proof}







\section{分离性条件}
\label{groupoids-section-separation}

\noindent
给定 $S$ 上的群胚 $(U, R, s, t, c)$ 时，这实际上是对态射
$j : R \to U \times_S U$ 的条件。与前一节一样，先写出相应的图。

\begin{lemma}
\label{groupoids-lemma-diagram-diagonal}
设 $S$ 为概形，$(U, R, s, t, c)$ 为 $S$ 上的群胚，
$G \to U$ 为稳定子群概形。交换图
$$
\xymatrix{
R \ar[d]^{\Delta_{R/U \times_S U}} \ar[rrr]_{f \mapsto (f, s(f))} & & &
R \times_{s, U} U \ar[d] \ar[r] & U \ar[d] \\
R \times_{(U \times_S U)} R \ar[rrr]^{(f, g) \mapsto (f, f^{-1} \circ g)} & & &
R \times_{s, U} G \ar[r] & G
}
$$
中的两个左侧水平箭头都是同构，右侧方块是纤维积方块。
\end{lemma}

\begin{proof}
证明略。这是关于定义以及代数几何函子观点的练习。
\end{proof}

\begin{lemma}
\label{groupoids-lemma-diagonal}
设 $S$ 为概形，$(U, R, s, t, c)$ 为 $S$ 上的群胚，
$G \to U$ 为稳定子群概形。
\begin{enumerate}
\item 下列条件等价：
\begin{enumerate}
\item $j : R \to U \times_S U$ 分离；
\item $G \to U$ 分离；
\item $e : U \to G$ 是闭浸入。
\end{enumerate}
\item 下列条件等价：
\begin{enumerate}
\item $j : R \to U \times_S U$ 拟分离；
\item $G \to U$ 拟分离；
\item $e : U \to G$ 拟紧。
\end{enumerate}
\end{enumerate}
\end{lemma}

\begin{proof}
群概形 $G \to U$ 是 $R \to U \times_S U$
沿对角态射 $U \to U \times_S U$ 的基变换，见
引理 \ref{groupoids-lemma-groupoid-stabilizer}。因此，若
$j$ 分离（resp.\ 拟分离），则
$G \to U$ 分离（resp.\ 拟分离）
（见《概形》引理 \ref{schemes-lemma-separated-permanence}）。
故在 (1) 与 (2) 中均有 (a) $\Rightarrow$ (b)。

\medskip\noindent
若 $G \to U$ 分离（resp.\ 拟分离），则态射
$U \to G$ 作为结构态射 $G \to U$ 的截面，是闭浸入
（resp.\ 拟紧）；见《概形》引理 \ref{schemes-lemma-section-immersion}。
故在 (1) 与 (2) 中均有 (b) $\Rightarrow$ (a)。

\medskip\noindent
由引理 \ref{groupoids-lemma-diagram-diagonal} 的结果
（以及《概形》引理 \ref{schemes-lemma-base-change-immersion}
和 \ref{schemes-lemma-quasi-compact-preserved-base-change}），
若 $e$ 是闭浸入（resp.\ 拟紧），则
$\Delta_{R/U \times_S U}$ 是闭浸入（resp.\ 拟紧）。
故在 (1) 与 (2) 中均有 (c) $\Rightarrow$ (a)。
\end{proof}









\section{有限平坦群胚：仿射情形}
\label{groupoids-section-finite-flat}

\noindent
设 $S$ 为概形，$(U, R, s, t, c)$ 为 $S$ 上的群胚概形。
假设 $U = \Spec(A)$ 与 $R = \Spec(B)$ 仿射。此时得到两个环同态
$s^\sharp, t^\sharp : A \longrightarrow B$。
令 $C$ 为 $s^\sharp$ 与 $t^\sharp$ 的等化子，即
\begin{equation}
\label{groupoids-equation-invariants}
C = \{a \in A \mid t^\sharp(a) = s^\sharp(a) \}.
\end{equation}
有时称它为{\it $U$ 上 $R$-不变函数环}。
$M = \Spec(C)$ 有什么性质？首先，图
$$
\xymatrix{
R \ar[r]_s \ar[d]_t & U \ar[d] \\
U \ar[r] & M
}
$$
交换；换言之，态射 $U \to M$ 等化 $s, t$。
此外，若 $T$ 为任意仿射概形，且态射 $U \to T$ 等化 $s, t$，
则 $U \to T$ 通过 $U \to M$ 分解。也就是说，
$U \to M$ 是仿射概形范畴中的余等化子。

\medskip\noindent
我们希望找出保证态射 $U \to M$ 在概形范畴中确为“商”的条件。
这一问题将在别处详论（此处插入未来的引用）；这里仅讨论若干特殊情形，
主要关注 $s, t$ 有限局部自由的情形。

\begin{example}
\label{groupoids-example-quotient-projective-line}
设 $k$ 为域，$U = \text{GL}_{2, k}$，而
$B \subset \text{GL}_2$ 是上三角矩阵的闭子群概形。
则商层 $\text{GL}_{2, k}/B$（在 Zariski、\'etale 或 fppf 拓扑下，
见定义 \ref{groupoids-definition-quotient-sheaf}）由射影直线表示：
$\mathbf{P}^1 = \text{GL}_{2, k}/B$。（细节略。）
另一方面，此情形下的不变函数环仅为 $k$。注意，此时态射
$s, t : R = \text{GL}_{2, k} \times_k B \to \text{GL}_{2, k} = U$
是相对维数为 $3$ 的光滑态射。
\end{example}

\noindent
回忆在《习题》习题 \ref{exercises-exercise-trace-det} 和
\ref{exercises-exercise-trace-det-rings} 中，我们已经定义了
有限局部自由模及有限局部自由环扩张的行列式和范数。
若 $\varphi : A \to B$ 是有限局部自由环同态，则以
$\text{Norm}_\varphi(b) \in A$ 表示 $b \in B$ 的范数。
对于有限局部自由概形态射，范数构造见
《除子》引理 \ref{divisors-lemma-finite-locally-free-has-norm}。

\begin{lemma}
\label{groupoids-lemma-determinant-trick}
设 $S$ 为概形，$(U, R, s, t, c)$ 为 $S$ 上的群胚概形。
假设 $U = \Spec(A)$ 与 $R = \Spec(B)$ 仿射，并且
$s, t : R \to U$ 有限局部自由。令 $C$ 如
(\ref{groupoids-equation-invariants})，并取 $f \in A$。则
$\text{Norm}_{s^\sharp}(t^\sharp(f)) \in C$。
\end{lemma}

\begin{proof}
考虑引理 \ref{groupoids-lemma-diagram} 的交换图
$$
\xymatrix{
& U & \\
R \ar[d]_s \ar[ru]^t &
R \times_{s, U, t} R
\ar[l]^-{\text{pr}_0} \ar[d]^{\text{pr}_1} \ar[r]_-c &
R \ar[d]^s \ar[lu]_t \\
U & R \ar[l]_t \ar[r]^s & U
}
$$
把 $f \in \Gamma(U, \mathcal{O}_U)$ 看作函数。
图上半部分的交换性说明，作为
$\Gamma(R \times_{S, U, t} R, \mathcal{O})$ 中的元素，有
$\text{pr}_0^\sharp(t^\sharp(f)) = c^\sharp(t^\sharp(f))$。
观察右下笛卡尔方块，范数构造与基变换的相容性给出
$s^\sharp(\text{Norm}_{s^\sharp}(t^\sharp(f))) =
\text{Norm}_{\text{pr}_1^\sharp}(c^\sharp(t^\sharp(f)))$。
类似地得到
$t^\sharp(\text{Norm}_{s^\sharp}(t^\sharp(f))) =
\text{Norm}_{\text{pr}_1^\sharp}(\text{pr}_0^\sharp(t^\sharp(f)))$。
因此由本证明中的第一个等式，
$s^\sharp(\text{Norm}_{s^\sharp}(t^\sharp(f))) =
t^\sharp(\text{Norm}_{s^\sharp}(t^\sharp(f)))$，如所需。
\end{proof}

\begin{lemma}
\label{groupoids-lemma-finite-locally-free-disjoint-free}
设 $S$ 为概形，$(U, R, s, t, c)$ 为 $S$ 上的群胚概形。
假设 $s, t : R \to U$ 有限局部自由。则
$$
U = \coprod\nolimits_{r \geq 1} U_r
$$
是 $R$-不变开集的不交并，并且 $R$ 在 $U_r$ 上的限制
$R_r$ 满足 $s, t : R_r \to U_r$ 有限局部自由、秩为 $r$。
\end{lemma}

\begin{proof}
由《概形态射》引理 \ref{morphisms-lemma-finite-locally-free}，
存在分解 $U = \coprod\nolimits_{r \geq 0} U_r$，使
$s : s^{-1}(U_r) \to U_r$ 有限局部自由、秩为 $r$。
由于 $s$ 满，$U_0 = \emptyset$。注意，
$u \in U_r \Leftrightarrow$ 当且仅当概形论纤维
$s^{-1}(u)$ 在 $\kappa(u)$ 上次数为 $r$。若
$z \in R$ 满足 $s(z) = u$、$t(z) = u'$，则采用
引理 \ref{groupoids-lemma-diagram} 的记号，
$$
\text{pr}_1^{-1}(z) \to \Spec(\kappa(z))
$$
由所引引理知，它同时是
$s^{-1}(u) \to \Spec(\kappa(u))$ 与
$s^{-1}(u') \to \Spec(\kappa(u'))$ 的基变换。因此
$u \in U_r \Leftrightarrow u' \in U_r$；换言之，诸开子集
$U_r$ 都是 $R$-不变的。特别地，$R$ 在 $U_r$ 上的限制就是
$s^{-1}(U_r)$，且 $s : R_r \to U_r$ 有限局部自由、秩为 $r$。
由 $R_r$ 的逆，$t : R_r \to U_r$ 同构于 $s$，故其秩也为 $r$。
\end{proof}

\begin{lemma}
\label{groupoids-lemma-integral-over-invariants}
设 $S$ 为概形，$(U, R, s, t, c)$ 为 $S$ 上的群胚概形。
假设 $U = \Spec(A)$ 与 $R = \Spec(B)$ 仿射，并且
$s, t : R \to U$ 有限局部自由。令 $C \subset A$ 如
(\ref{groupoids-equation-invariants})。则 $A$ 在 $C$ 上整。
\end{lemma}

\begin{proof}
首先，由引理 \ref{groupoids-lemma-finite-locally-free-disjoint-free}，
$(U, R, s, t, c)$ 是群胚概形 $(U_r, R_r, s, t, c)$ 的不交并，
其中每个 $s, t : R_r \to U_r$ 的秩恒为 $r$。
由于 $U$ 拟紧，除有限多个 $r$ 外均有 $U_r = \emptyset$。
只需对每个 $(U_r, R_r, s, t, c)$ 证明该引理，故可假设
$s, t$ 有限局部自由、秩为 $r$。

\medskip\noindent
假设 $s, t$ 有限局部自由、秩为 $r$。取 $f \in A$，
考虑元素 $x - f \in A[x]$，其中把 $x$ 看作
$\mathbf{A}^1$ 上的坐标。由于
$$
(U \times \mathbf{A}^1, R \times \mathbf{A}^1,
s \times \text{id}_{\mathbf{A}^1},
t \times \text{id}_{\mathbf{A}^1},
c \times \text{id}_{\mathbf{A}^1})
$$
也是源和靶均有限的群胚概形，可对它应用
引理 \ref{groupoids-lemma-determinant-trick}，得到
$P(x) = \text{Norm}_{s^\sharp}(t^\sharp(x - f))$ 属于
$C[x]$。由于 $s^\sharp : A \to B$ 有限局部自由、秩为 $r$，
$P$ 是次数为 $r$ 的首一多项式。此外，由 Cayley--Hamilton 定理
有 $P(f) = 0$（《交换代数》引理 \ref{algebra-lemma-charpoly}）。
更准确地，Cayley--Hamilton 定理说明多项式
$s^\sharp(P) \in B[x]$ 零化元素 $t^\sharp(f)$。
由于 $P$ 的系数在 $C$ 中，得到 $t^\sharp(P(f)) = 0$ 于 $B$ 中。
又因 $t$ 忠实平坦，得到 $P(f) = 0$。
\end{proof}

\begin{lemma}
\label{groupoids-lemma-invariants-base-change}
设 $S$ 为概形，$(U, R, s, t, c)$ 为 $S$ 上的群胚概形。
假设 $U = \Spec(A)$ 与 $R = \Spec(B)$ 仿射，
$s, t : R \to U$ 有限局部自由。令 $C \subset A$ 如
(\ref{groupoids-equation-invariants})。给定环同态 $C \to C'$，并令
$U' = \Spec(A \otimes_C C')$、
$R' = \Spec(B \otimes_C C')$。则：
\begin{enumerate}
\item 映射 $s, t, c$ 诱导映射 $s', t', c'$，使
$(U', R', s', t', c')$ 成为群胚概形。令 $C^1 \subset A'$
为 $U'$ 上的 $R'$-不变函数。
\item 典范映射 $\varphi : C' \to C^1$ 满足：
\begin{enumerate}
\item 对每个 $f \in C^1$，存在 $n > 0$ 及多项式
$P \in C'[x]$，使其在 $C^1[x]$ 中的像为 $(x - f)^n$；
\item 对每个 $f \in \Ker(\varphi)$，存在 $n > 0$ 使
$f^n = 0$。
\end{enumerate}
\item 若 $C \to C'$ 平坦，则 $\varphi$ 是同构。
\end{enumerate}
\end{lemma}

\begin{proof}
第 (1) 点的证明略。记 $A' = A \otimes_C C'$、
$B' = B \otimes_C C'$。则
$$
C^1
= \{a \in A' \mid (t')^\sharp(a) = (s')^\sharp(a) \}
= \{a \in A \otimes_C C' \mid t^\sharp \otimes 1(a) = s^\sharp \otimes 1(a) \}.
$$
换言之，$C^1$ 是差映射
$(t^\sharp - s^\sharp) \otimes 1$ 的核，而后者正是
$C$-线性映射 $t^\sharp - s^\sharp : A \to B$
沿 $C \to C'$ 的基变换。因此 (3) 成立。

\medskip\noindent
证明 (2)(b)。由于 $C \to A$ 整
（引理 \ref{groupoids-lemma-integral-over-invariants}）且单，
$\Spec(A) \to \Spec(C)$ 满，见
《交换代数》引理 \ref{algebra-lemma-integral-overring-surjective}。
因此其满态射基变换 $\Spec(A') \to \Spec(C')$ 也满
（《概形态射》引理 \ref{morphisms-lemma-base-change-surjective}）。
于是 $\Spec(C^1) \to \Spec(C')$ 也满。这例如由
《交换代数》引理 \ref{algebra-lemma-image-dense-generic-points}
蕴含 (2)(b)。

\medskip\noindent
证明 (2)(a)。由引理 \ref{groupoids-lemma-finite-locally-free-disjoint-free}，
群胚概形 $(U, R, s, t, c)$ 分解为有限多个群胚概形
$(U_r, R_r, s, t, c)$ 的不交并，其中
$s, t : R_r \to U_r$ 有限局部自由、秩为 $r$。
沿 $U' = \Spec(C') \to U$ 拉回，得到 $U'$ 及
$U^1 = \Spec(C^1)$ 的类似分解。下一段将证明，对相应环系
$A_r, B_r, C_r, C'_r, C^1_r$，取 $n = r$ 时 (2)(a) 成立。
给定 $f \in C^1$，令 $P_r \in C_r[x]$ 为如下多项式：
其在 $C^1_r[x]$ 中的像是 $(x - f)^r$ 的像。
选择充分可除的整数 $n$，可得多项式 $P \in C'[x]$，
其在 $C^1[x]$ 中的像为 $(x - f)^n$；具体地，取
$P$ 为 $C'[x]$ 中唯一满足其在 $C'_r[x]$ 中的像为
$P_r^{n/r}$ 的元素。

\medskip\noindent
本段在环同态 $s^\sharp, t^\sharp : A \to B$
有限局部自由且秩恒为 $r$ 的情形下证明 (2)(a)。
取 $f \in C^1 \subset A' = A \otimes_C C'$。
选择平坦 $C$-代数 $D$ 以及满射 $D \to C'$，并选择
$f$ 的提升 $g \in A \otimes_C D$。考虑多项式
$$
P = \text{Norm}_{s^\sharp \otimes 1}(t^\sharp \otimes 1(x - g))
$$
于 $(A \otimes_C D)[x]$ 中。由
引理 \ref{groupoids-lemma-determinant-trick} 及本引理的 (3)，
$P$ 的系数属于 $D$（比较引理 \ref{groupoids-lemma-integral-over-invariants} 的证明）。
另一方面，$P$ 在 $(A \otimes_C C')[x]$ 中的像是
$(x - f)^r$，因为
$t^\sharp \otimes 1(x - f) = s^\sharp(x - f)$，且
$s^\sharp$ 有限局部自由、秩为 $r$。
取陈述 (2)(a) 中的 $P$ 为上述 $P$ 在映射
$D[x] \to C'[x]$ 下的像，即得所需。
\end{proof}

\begin{lemma}
\label{groupoids-lemma-points}
设 $S$ 为概形，$(U, R, s, t, c)$ 为 $S$ 上的群胚概形。
假设 $U = \Spec(A)$ 与 $R = \Spec(B)$ 仿射，
$s, t : R \to U$ 有限局部自由，并令 $C \subset A$ 如
(\ref{groupoids-equation-invariants})。则 $U \to M = \Spec(C)$ 具有下列性质：
\begin{enumerate}
\item 点集映射 $|U| \to |M|$ 是满射，并且
$u_0, u_1 \in |U|$ 映到同一点当且仅当存在
$r \in |R|$，满足 $t(r) = u_0$ 且 $s(r) = u_1$；即
$$
|M| = |U|/|R|
$$
\item 对任意代数闭域 $k$，有
$$
M(k) = U(k)/R(k)
$$
\end{enumerate}
\end{lemma}

\begin{proof}
由于 $C \to A$ 整（引理 \ref{groupoids-lemma-integral-over-invariants}）且单，
$\Spec(A) \to \Spec(C)$ 满，见
《交换代数》引理 \ref{algebra-lemma-integral-overring-surjective}。
因此 $|U| \to |M|$ 满。

\medskip\noindent
设 $k$ 为代数闭域，$C \to k$ 为环同态。
满态射在基变换下保持
（《概形态射》引理 \ref{morphisms-lemma-base-change-surjective}），
故 $A \otimes_C k$ 非零。此时 $k \subset A \otimes_C k$
是非零整扩张。因此 $A \otimes_C k$ 的任意剩余域都是
$k$ 的代数扩张，从而等于 $k$。于是 $U(k) \to M(k)$ 满。

\medskip\noindent
设 $a_0, a_1 : A \to k$ 为两个环同态。若存在环同态
$b : B \to k$，使 $a_0 = b \circ t^\sharp$ 且
$a_1 = b \circ s^\sharp$，则按定义
$a_0|_C = a_1|_C$。因此映射 $U(k) \to M(k)$
等化两个映射 $R(k) \to U(k)$。反之，假设
$a_0|_C = a_1|_C$，并将此代数同态记作 $c : C \to k$。
考虑图
$$
\xymatrix{
& &
B \ar@{-->}[lld] \\
k & &
A
\ar@<0.5ex>[ll]^{a_0}
\ar@<-0.5ex>[ll]_{a_1}
\ar@<1ex>[u]
\ar@<-1ex>[u] \\
& &
C \ar[u] \ar[llu]^c
}
$$
若能构造使该图交换的虚线箭头，则引理第 (2) 点的证明完成。
由于 $s : A \to B$ 有限，存在有限多个环同态
$b_1, \ldots, b_n : B \to k$，使
$b_i \circ s^\sharp = a_1$。若虚线箭头不存在，则
$a'_i = b_i \circ t^\sharp$（$i = 1, \ldots, n$）
均不等于 $a_0$。于是 $A \otimes_C k$ 的极大理想
$$
\mathfrak m'_i = \Ker(a_i' \otimes 1 : A \otimes_C k \to k)
$$
均不同于
$\mathfrak m = \Ker(a_0 \otimes 1 : A \otimes_C k \to k)$。
由《交换代数》引理 \ref{algebra-lemma-silly}，可得元素
$f \in A \otimes_C k$，满足 $f \in \mathfrak m$，但对
$i = 1, \ldots, n$ 有 $f \not \in \mathfrak m_i'$。
考虑范数
$$
g = \text{Norm}_{s^\sharp \otimes 1}(t^\sharp \otimes 1(f))
\in
A \otimes_C k
$$
由引理 \ref{groupoids-lemma-determinant-trick}，它属于经映射
$c : C \to k$ 基变换所得群胚的不变函数
$C^1 \subset A \otimes_C k$。一方面，$a_1(g) \in k^*$，
因为 $t^\sharp(f)$ 在 $a_1$ 上方所有点
（对应于 $b_1, \ldots, b_n$）处的值都可逆
（此处插入关于行列式性质的未来引用）。
另一方面，由于 $f \in \mathfrak m$，$f$ 不是单位，
故 $t^\sharp(f)$ 也不是单位（因为
$t^\sharp \otimes 1$ 忠实平坦），所以它的范数不是单位
（此处插入关于行列式性质的未来引用）。
于是 $C^1$ 含有一个既非幂零又非单位的元素。
下面说明这导致矛盾。把引理 \ref{groupoids-lemma-invariants-base-change}
应用于映射 $c : C \to C' = k$，可知 $k$ 到不变函数
$C^1$ 的映射是单射；而且对任意元素 $x \in C^1$，
存在整数 $n > 0$ 使 $x^n \in k$。因此 $C^1$ 的每个元素
要么是单位，要么幂零。

\medskip\noindent
还需完成 (1) 的证明。已知 $|U| \to |M|$ 满，且显然
$|U| \to |M|$ 是 $|R|$-不变的。最后，假设
$u_0, u_1 \in U$ 映到同一点 $m \in M$。
诱导的域扩张 $\kappa(u_0)/\kappa(m)$ 和
$\kappa(u_1)/\kappa(m)$ 都是代数扩张（如上所用，因为
$A$ 在 $C$ 上整）。若 $k$ 是 $\kappa(m)$ 的代数闭包，
则可找到 $\kappa(m)$-嵌入
$\overline{u}_0 : \kappa(u_0) \to k$ 和
$\overline{u}_1 : \kappa(u_1) \to k$。它们确定映到
$M(k)$ 中同一点的 $k$-值点
$\overline{u}_0, \overline{u}_1 \in U(k)$。
由 (2)，存在点 $\overline{r} \in R(k)$，满足
$s(\overline{r}) = \overline{u}_0$ 且
$t(\overline{r}) = \overline{u}_1$。
$\overline{r}$ 的像 $r \in R$ 满足
$s(r) = u_0$ 且 $t(r) = u_1$，如所需。
\end{proof}

\begin{lemma}
\label{groupoids-lemma-etale}
设 $S$ 为概形，$f : (U', R', s', t') \to (U, R, s, t, c)$
为 $S$ 上群胚概形的态射。假设：
\begin{enumerate}
\item $U$, $R$, $U'$, $R'$ 都仿射；
\item $s, t, s', t'$ 有限局部自由；
\item 图
$$
\xymatrix{
R' \ar[d]_{s'} \ar[r]_f & R \ar[d]^s \\
U' \ar[r]^f & U
}
\quad
\quad
\xymatrix{
R' \ar[d]_{t'} \ar[r]_f & R \ar[d]^t \\
U' \ar[r]^f & U
}
\quad
\quad
\xymatrix{
G' \ar[d] \ar[r]_f & G \ar[d] \\
U' \ar[r]^f & U
}
$$
均为笛卡尔的，其中 $G$ 与 $G'$ 为稳定子群概形；
\item $f : U' \to U$ 是 \'etale 态射。
\end{enumerate}
则从 $U$ 上 $R$-不变函数到 $U'$ 上 $R'$-不变函数的映射
$C \to C'$ 是 \'etale 的，并且
$U' = \Spec(C') \times_{\Spec(C)} U$。
\end{lemma}

\begin{proof}
令 $M = \Spec(C)$、$M' = \Spec(C')$。写
$U = \Spec(A)$、$U' = \Spec(A')$、$R = \Spec(B)$ 及
$R' = \Spec(B')$。下文将直接使用引理 \ref{groupoids-lemma-integral-over-invariants}、
\ref{groupoids-lemma-invariants-base-change} 和
\ref{groupoids-lemma-points} 的结果。

\medskip\noindent
假设 $C$ 是严格 Hensel 局部环。设 $p \in M$ 为闭点，
$p' \in M'$ 映到 $p$。断言：此时存在
$(U, R, s, t, c)$ 上的不交并分解
$(U', R', s', t', c') = (U, R, s, t, c) \amalg (U'', R'', s'', t'', c'')$，
使得对 $M$ 上相应的不交并分解
$M' = M \amalg M''$，点 $p'$ 对应于 $p \in M$。

\medskip\noindent
该断言蕴含引理。假设 $M_1 \to M$ 是仿射概形的平坦态射。
可将一切基变换到 $M_1$，而不影响假设 (1) -- (4)。
由引理 \ref{groupoids-lemma-invariants-base-change}，
$M_1$、resp.\ $M_1'$ 分别是 $U_1$ 上 $R_1$-不变函数、
resp.\ $U'_1$ 上 $R'_1$-不变函数的谱。
假设 $p' \in M'$ 映到 $p \in M$。令 $M_1$ 为
$\mathcal{O}_{M, p}$ 的严格 Hensel 化之谱，其闭点为
$p_1 \in M_1$。选择映到 $p_1$ 与 $p'$ 的点
$p'_1 \in M'_1$。由该断言得到
$$
(U'_1, R'_1, s'_1, t'_1, c'_1) =
(U_1, R_1, s_1, t_1, c_1) \amalg
(U''_1, R''_1, s''_1, t''_1, c''_1)
$$
并且相应地，作为 $M_1$ 上的概形有
$M'_1 = M_1 \amalg M''_1$。写 $M_1 = \Spec(C_1)$，
并把 $C_1 = \colim C_i$ 写成 \'etale $C$-代数的滤过余极限。
令 $M_i = \Spec(C_i)$。则 $M_1 = \lim M_i$，其他概形亦然。
由《概形的极限》引理 \ref{limits-lemma-descend-opens} 和
\ref{limits-lemma-descend-isomorphism}，可找到 $i$，使
$$
(U'_i, R'_i, s'_i, t'_i, c'_i) =
(U_i, R_i, s_i, t_i, c_i) \amalg
(U''_i, R''_i, s''_i, t''_i, c''_i)
$$
故 $M'_i = M_i \amalg M''_i$。特别地，经过 \'etale 基变换后，
$M' \to M$ 在 $p'$ 上方某点成为 \'etale 的。
这蕴含 $M' \to M$ 在 $p'$ 处为 \'etale
（例如由《概形态射》引理 \ref{morphisms-lemma-set-points-where-fibres-etale}）。
证明断言之后，我们将证明 $U' \cong M' \times_M U$。

\medskip\noindent
证明该断言。注意 $U_p$ 与 $U'_{p'}$ 都只有有限多个点。
对 $u \in U_p$，扩张 $\kappa(u)/\kappa(p)$ 是代数的，
故 $\kappa(u)$ 可分闭。由于 $U' \to U$ 为 \'etale，
态射 $U'_{p'} \to U_p$ 在剩余域扩张上诱导同构。
取映到 $u \in U_p$ 的 $u' \in U'_{p'}$。
由假设 (3)，概形论纤维的态射
$(s')^{-1}(u') \to s^{-1}(u)$、
$(t')^{-1}(u') \to t^{-1}(u)$ 及
$G'_{u'} \to G_u$ 都是同构。注意
$U_p = t(s^{-1}(u))$（集合论意义下），可知
$U'_{p'}$ 的点满射到 $U_p$ 的点。
假设 $U'_{p'}$ 中两点 $u'_1$ 与 $u'_2$ 映到
$U_p$ 的同一点 $u$。则存在点
$r' \in R'_{p'}$，满足 $s'(r') = u'_1$ 且 $t'(r') = u'_2$。
考虑两个域塔
$$
\kappa(r')/\kappa(u'_1)/\kappa(u)/\kappa(p) \quad
\kappa(r')/\kappa(u'_2)/\kappa(u)/\kappa(p)
$$
它们的两个“端点”分别与下面两个域塔的两个“端点”相同：
$$
\kappa(r')/\kappa(u'_1)/\kappa(p')/\kappa(p) \quad
\kappa(r')/\kappa(u'_2)/\kappa(p')/\kappa(p)
$$
由于 $(U'_{p'}, R'_{p'}, s', t', c')$ 是 $p'$ 上的群胚，
后两者诱导相同映射 $\kappa(p') \to \kappa(r')$。
又因 $\kappa(u)/\kappa(p)$ 纯不可分，所诱导的两个映射
$\kappa(u) \to \kappa(r')$ 相同。因此 $r'$ 映到纤维
$G_u$ 的一点。由假设 (3)，$r' \in (G')_{u'_1}$。
具体地，可把 $G$ 看成 $R$ 的闭子概形，其中后者经
$s$ 视为 $U$ 上的概形；再利用其向 $U'$ 的基变换给出
$G' \subset R'$。特别地，$u'_1 = u'_2$。
故 $U'_{p'} \to U_p$ 在点上双射，并在剩余域上诱导同构。
于是 $U'_{p'}$ 是有限闭点集
（《交换代数》引理 \ref{algebra-lemma-finite-residue-extension-closed}），
因而 $U'_{p'}$ 在 $U'$ 中闭。令 $J' \subset A'$ 为在集合论意义下
截出 $U'_{p'}$ 的根理想。

\medskip\noindent
断言证明的第二部分。令 $\mathfrak m \subset C$ 为极大理想。
由《代数进阶》引理 \ref{more-algebra-lemma-integral-over-henselian-pair}，
$(A, \mathfrak m A)$ 是 Hensel 对。令
$J = \sqrt{\mathfrak m A}$。由《代数进阶》引理 \ref{more-algebra-lemma-change-ideal-henselian-pair}，
$(A, J)$ 是 Hensel 对；根据上面的讨论，\'etale 环同态
$A \to A'$ 诱导同构 $A/J \to A'/J'$。
由《代数进阶》引理 \ref{more-algebra-lemma-characterize-henselian-pair}，
$A' = A \times A''$。考虑相应的不交并分解
$U' = U \amalg U''$。开集 $(s')^{-1}(U)$ 是 $R'$
中可特化到 $R'_{p'}$ 一点的点集；$(t')^{-1}(U)$ 亦然。
类似地，$(s')^{-1}(U'') = (t')^{-1}(U'')$，
因为这是不能特化到 $R'_{p'}$ 的点集。因此得到不交并分解
$$
(U', R', s', t', c') =
(U, R, s, t, c) \amalg
(U'', R'', s'', t'', c'')
$$
这立即给出 $M' = M \amalg M''$，断言得证。

\medskip\noindent
最后还需证明典范映射 $U' \to M' \times_M U$ 是同构。
它是 \'etale 态射
（《概形态射》引理 \ref{morphisms-lemma-etale-permanence}）。
另一方面，基变换到严格 Hensel 局部环
（如证明第三段所做），并利用断言证明过程中得到的双射
$U'_{p'} \to U_p$，可知 $U' \to M' \times_M U$
普遍双射（略去若干细节）。然而，普遍双射的 \'etale 态射
是同构（《下降》引理 \ref{descent-lemma-universally-injective-etale-open-immersion}），
证明完成。
\end{proof}

\begin{lemma}
\label{groupoids-lemma-basis}
设 $S$ 为概形，$(U, R, s, t, c)$ 为 $S$ 上的群胚概形。假设：
\begin{enumerate}
\item $U = \Spec(A)$ 与 $R = \Spec(B)$ 仿射；
\item 存在元素 $x_i \in A$（$i \in I$），使
$B = \bigoplus_{i \in I} s^\sharp(A)t^\sharp(x_i)$。
\end{enumerate}
则 $A = \bigoplus_{i\in I} Cx_i$，且 $B \cong A \otimes_C A$，
其中 $C \subset A$ 是 (\ref{groupoids-equation-invariants}) 中
$U$ 上的 $R$-不变函数。
\end{lemma}

\begin{proof}
本证明中以 $s, t : A \to B$ 代替
$s^\sharp, t^\sharp$，类似地以
$c : B \to B \otimes_{s, A, t} B$ 表示相应映射。
记 $p_0 : B \to B \otimes_{s, A, t} B$、
$b \mapsto b \otimes 1$，并记
$p_1 : B \to B \otimes_{s, A, t} B$、
$b \mapsto 1 \otimes b$。由引理 \ref{groupoids-lemma-diagram-pull}
及 $C$ 的定义，有交换图
$$
\xymatrix{
B \otimes_{s, A, t} B &
B \ar@<-1ex>[l]_-c \ar@<1ex>[l]^-{p_0} &
A \ar[l]^t \\
B \ar[u]^{p_1} &
A \ar@<-1ex>[l]_s \ar@<1ex>[l]^t \ar[u]_s &
C \ar[u] \ar[l]
}
$$
此外，左侧两个方块在环范畴中是余笛卡尔的，而顶行同构于图
$$
\xymatrix{
B \otimes_{t, A, t} B &
B \ar@<-1ex>[l]_-{p_1} \ar@<1ex>[l]^-{p_0} &
A \ar[l]^t
}
$$
由《下降》引理 \ref{descent-lemma-ff-exact}，
它是等化子图，因为条件 (2) 特别蕴含 $s$
（以及与之同构的箭头 $t$）忠实平坦。
底行按 $C$ 的定义也是等化子图。
利用 $x_i$ 可得到交换图
$$
\xymatrix{
B \otimes_{s, A, t} B &
B \ar@<-1ex>[l]_-c \ar@<1ex>[l]^-{p_0} &
A \ar[l]^t \\
\bigoplus_{i \in I} B x_i \ar[u]^{p_1} &
\bigoplus_{i \in I} A x_i \ar@<-1ex>[l]_s \ar@<1ex>[l]^t \ar[u]_s &
\bigoplus_{i \in I} C x_i \ar[u] \ar[l]
}
$$
其中，右侧竖直箭头把 $x_i$ 映到 $x_i$，
中间竖直箭头把 $x_i$ 映到 $t(x_i)$，左侧竖直箭头把
$x_i$ 映到
$c(t(x_i)) = t(x_i) \otimes 1 = p_0(t(x_i))$
（该等式来自引理 \ref{groupoids-lemma-diagram} 中图上半部分的交换性）。
于是该图交换。由假设，中间竖直箭头是同构。
由于左侧两个方块余笛卡尔，左侧竖直箭头也是同构。
另一方面，水平行正合（即均为等化子），故右侧竖直箭头也是同构。
\end{proof}

\begin{proposition}
\label{groupoids-proposition-finite-flat-equivalence}
设 $S$ 为概形，$(U, R, s, t, c)$ 为 $S$ 上的群胚概形。假设：
\begin{enumerate}
\item $U = \Spec(A)$ 与 $R = \Spec(B)$ 仿射；
\item $s, t : R \to U$ 有限局部自由；
\item $j = (t, s)$ 是等价关系。
\end{enumerate}
在此情形下，令 $C \subset A$ 如
(\ref{groupoids-equation-invariants})。则 $U \to M = \Spec(C)$
有限局部自由，且 $R = U \times_M U$。此外，$M$
表示 fppf 拓扑下的商层 $U/R$
（见定义 \ref{groupoids-definition-quotient-sheaf}）。
\end{proposition}

\begin{proof}
本证明中以 $s, t : A \to B$ 代替
$s^\sharp, t^\sharp$。由
引理 \ref{groupoids-lemma-criterion-quotient-representable}，
只需证明 $C \to A$ 有限局部自由，并且映射
$$
t \otimes s : A \otimes_C A \longrightarrow B
$$
是同构。首先，$j$ 是单态射，而且有限
（因为 $s$ 和 $t$ 已经有限）。因此由
《概形态射》引理 \ref{morphisms-lemma-finite-monomorphism-closed}，
$j$ 是闭浸入。故 $A \otimes_C A \to B$ 满。

\medskip\noindent
我们将如引理 \ref{groupoids-lemma-invariants-base-change} 那样，
沿平坦环同态 $C \to C'$ 作基变换，并使用不变函数的形成与平坦基变换交换
这一事实，见所引引理的 (3)。下面将证明：对每个素理想
$\mathfrak p \subset C$，存在局部平坦环同态
$C_{\mathfrak p} \to C_{\mathfrak p}'$，使得基变换到
$C_{\mathfrak p}'$ 后结论成立。这立即蕴含
$A \otimes_C A \to B$ 是单射（使用
《交换代数》引理 \ref{algebra-lemma-characterize-zero-local}）。
结合《交换代数》引理 \ref{algebra-lemma-local-flat-ff}、
\ref{algebra-lemma-flat-localization} 及
\ref{algebra-lemma-flatness-descends}，还可推出
$C \to A$ 平坦。又因 $U \to \Spec(C)$ 满
（引理 \ref{groupoids-lemma-points}），故 $C \to A$ 忠实平坦。
同构 $B \cong A \otimes_C A$ 蕴含 $A$ 是有限表现
$C$-模，见《交换代数》引理 \ref{algebra-lemma-descend-properties-modules}。
于是 $A$ 在 $C$ 上有限局部自由，见
《交换代数》引理 \ref{algebra-lemma-finite-projective}。

\medskip\noindent
由引理 \ref{groupoids-lemma-finite-locally-free-disjoint-free}，
$A$ 是有限多个环 $A_r$ 的积，$B$ 是有限多个环 $B_r$ 的积，
且群胚概形相应分解（见引理 \ref{groupoids-lemma-integral-over-invariants} 的证明）。
于是 $C$ 也是诸环 $C_r$ 的积，$C'$ 也相应分解。
因此可以并且确实假设环同态
$s, t : A \to B$ 有限局部自由、恒秩为 $r$。

\medskip\noindent
将使用的局部环同态 $C_{\mathfrak p} \to C_{\mathfrak p}'$
可以取为任意使 $C_{\mathfrak p}'$ 的剩余域为无限域的
局部平坦环同态。由《交换代数》引理 \ref{algebra-lemma-flat-local-given-residue-field}，这样的同态存在。

\medskip\noindent
假设 $C$ 是极大理想为 $\mathfrak m$、剩余域无限的局部环，
并假设 $s, t : A \to B$ 有限局部自由、恒秩为 $r > 0$。
由于 $C \subset A$ 整
（引理 \ref{groupoids-lemma-integral-over-invariants}），
位于 $\mathfrak m$ 上方的所有素理想都是极大理想，
且 $A$ 的所有极大理想都位于 $\mathfrak m$ 上方。
$C \subset B$ 也同样如此。选取 $A$ 中位于 $\mathfrak m$
上方的极大理想 $\mathfrak m'$（其存在性由
引理 \ref{groupoids-lemma-points}）。因为 $t : A \to B$ 有限局部自由，
$B$ 中位于 $\mathfrak m'$ 上方的极大理想至多有限多个。
再次由引理 \ref{groupoids-lemma-points}，$A$ 只有有限多个极大理想，
即 $A$ 半局部；进而 $B$ 也半局部。现在，因为
$t \otimes s : A \otimes_C A \to B$ 满，可把
《交换代数》引理 \ref{algebra-lemma-semi-local-module-basis-in-submodule}
应用于环同态 $C \to A$、$A$-模 $M = B$
（经 $t$ 视作 $A$-模）以及 $C$-子模 $s(A) \subset B$。
该引理给出 $x_1, \ldots, x_r \in A$，使 $M$ 是以
$s(x_1), \ldots, s(x_r)$ 为基的自由 $A$-模。
因此应用引理 \ref{groupoids-lemma-basis} 可知
$C \to A$ 有限自由，且 $B \cong A \otimes_C A$。
\end{proof}



\section{有限平坦群胚}
\label{groupoids-section-finite-flat-general}

\noindent
本节证明一个引理，用来说明：若每个等价类都包含于某个仿射开集，
则概形对有限平坦等价关系取商仍是概形。见
《代数空间的性质》命题 \ref{spaces-properties-proposition-finite-flat-equivalence-global}。

\begin{lemma}
\label{groupoids-lemma-find-invariant-affine}
设 $S$ 为概形，$(U, R, s, t, c)$ 为 $S$ 上的群胚概形。
假设 $s$, $t$ 有限局部自由。设 $u \in U$ 为一点，并且
$t(s^{-1}(\{u\}))$ 包含于 $U$ 的某个仿射开集。
则 $u$ 在 $U$ 中有一个 $R$-不变仿射开邻域。
\end{lemma}

\begin{proof}
由于 $s$ 有限局部自由，其纤维有限。因此
$t(s^{-1}(\{u\})) = \{u_1, \ldots, u_n\}$ 是有限集。
注意 $u \in \{u_1, \ldots, u_n\}$。
取包含 $\{u_1, \ldots, u_n\}$ 的仿射开集
$W \subset U$；特别地，$u \in W$。考虑
$Z = R \setminus s^{-1}(W) \cap t^{-1}(W)$。这是 $R$ 的闭子集。
像 $t(Z)$ 是 $U$ 的闭子集；粗略地说，它由 $U$ 中与
$U \setminus W$ 某一点 $R$-等价的点组成。因此
$W' = U \setminus t(Z)$ 是 $U$ 中包含于 $W$ 的
$R$-不变开子概形，并且
$\{u_1, \ldots, u_n\} \subset W'$。图示如下：
$$
\{u_1, \ldots, u_n\} \subset W' \subset W \subset U.
$$
取元素 $f \in \Gamma(W, \mathcal{O}_W)$，满足
$\{u_1, \ldots, u_n\} \subset D(f) \subset W'$。
这样的 $f$ 由《交换代数》引理 \ref{algebra-lemma-silly} 存在。
根据 $W'$ 的选择，有 $s^{-1}(W') \subset t^{-1}(W)$，
因而得到图
$$
\xymatrix{
s^{-1}(W') \ar[d]_s \ar[r]_-t & W \\
W'
}
$$
由假设，竖直箭头有限局部自由。令
$$
g = \text{Norm}_s(t^\sharp f) \in \Gamma(W', \mathcal{O}_{W'})
$$
按构造，$g$ 是 $W'$ 上在 $u$ 处非零的函数，因为
$t^\sharp(f)$ 在 $R$ 中位于 $u$ 上方的每一点处都非零，
这是由于 $f$ 在 $u_1, \ldots, u_n$ 处非零。
类似地，$D(g) \subset W'$ 恰由如下点 $w$ 构成：
$f$ 在与 $w$ 等价的任一点处都不为零。这说明 $D(g)$
是 $W'$ 的 $R$-不变仿射开集。最终图示为
$$
\{u_1, \ldots, u_n\} \subset D(g) \subset D(f) \subset W' \subset W \subset U
$$
故结论成立。
\end{proof}







\section{拟射影概形的下降}
\label{groupoids-section-quasi-projective}

\noindent
可用引理 \ref{groupoids-lemma-find-invariant-affine}
说明某一类下降数据是有效的。

\begin{lemma}
\label{groupoids-lemma-descend-along-finite-with-bound}
设 $X \to Y$ 为满的有限局部自由态射，$d$ 为正整数。
假设对每个几何点 $\bar{y} : \mathrm{Spec}(k) \to Y$，
纤维 $X_{\bar{y}}$ 至多有 $d$ 个点。设 $V$ 为 $X$ 上的概形，
并且对所有 $(y, v_1, \ldots, v_d)$，其中 $y \in Y$ 且
$v_1, \ldots, v_d \in V_y$，都存在仿射开集
$U \subset V$，使 $v_1, \ldots, v_d \in U$。
则 $V/X/Y$ 上的任意下降数据都是有效的。
\end{lemma}

\begin{proof}
设 $\varphi$ 为《下降》定义 \ref{descent-definition-descent-datum} 意义下的下降数据。
回忆从 $Y$ 上概形到相对于 $\{X \to Y\}$ 的下降数据的函子全忠实，
见《下降》引理 \ref{descent-lemma-refine-coverings-fully-faithful}。
因此利用《概形的构造》引理 \ref{constructions-lemma-relative-glueing}，只需在 $Y$ 仿射时证明。
略去若干细节（若 $Y$ 分离或其对角仿射，则可避免这一论证，
因为此时从仿射概形到 $X$ 的每个态射都仿射）。

\medskip\noindent
假设 $Y$ 仿射。若 $V$ 也仿射，则有效性由
《下降》引理 \ref{descent-lemma-affine} 得出。因此由
《下降》引理 \ref{descent-lemma-effective-for-fpqc-is-local-upstairs}，
只需证明 $V$ 的每个点 $v$ 都有一个
$\varphi$-不变仿射开邻域。考虑群胚
$(X, X \times_Y X, \text{pr}_1, \text{pr}_0, \text{pr}_{02})$。
由引理 \ref{groupoids-lemma-cartesian-equivalent-descent-datum}，
下降数据 $\varphi$ 确定且被如下群胚概形的笛卡尔态射确定：
$$
(V, R, s, t, c)
\longrightarrow
(X, X \times_Y X, \text{pr}_1, \text{pr}_0, \text{pr}_{02})
$$
该态射在 $\Spec(\mathbf{Z})$ 上。由于 $X \to Y$ 有限局部自由，
$\text{pr}_i : X \times_Y X \to X$，从而 $s$ 与 $t$
都有限局部自由，且其几何纤维至多有 $d$ 个点。特别地，
点 $v \in V$ 的 $R$-轨道 $t(s^{-1}(\{v\}))$
至多有 $d$ 个点。再次使用引理 \ref{groupoids-lemma-cartesian-equivalent-descent-datum} 的范畴等价，
可知 $V$ 的 $\varphi$-不变开集恰是 $V$ 的 $R$-不变开集。
由假设，存在包含轨道 $t(s^{-1}(\{v\}))$ 的 $V$ 的仿射开集，
因为该轨道的所有点都映到 $Y$ 的同一点。因此引理 \ref{groupoids-lemma-find-invariant-affine} 给出包含 $v$ 的
$R$-不变仿射开集。
\end{proof}

\begin{lemma}
\label{groupoids-lemma-descend-along-finite}
设 $X \to Y$ 为满的有限局部自由态射。 
设 $V$ 为 $X$ 上的概形，并且对所有 
$(y, v_1, \ldots, v_d)$，其中 $y \in Y$ 且
$v_1, \ldots, v_d \in V_y$，都存在仿射开集
$U \subset V$，使 $v_1, \ldots, v_d \in U$。
则 $V/X/Y$ 上的任意下降数据都是有效的。
\end{lemma}

\begin{proof}
与引理 \ref{groupoids-lemma-descend-along-finite-with-bound} 的证明一样，
可假设 $Y$ 仿射，特别是拟紧。于是存在整数 $d$，使
$X \to Y$ 的所有几何纤维至多有 $d$ 个点，见
《概形态射》第 \ref{morphisms-section-universally-bounded} 节中的讨论。
关于 $V$ 的假设显然蕴含可应用
引理 \ref{groupoids-lemma-descend-along-finite-with-bound}，证明完成。
\end{proof}

\begin{lemma}
\label{groupoids-lemma-descend-along-finite-quasi-projective}
设 $X \to Y$ 为满的有限局部自由态射，$V$ 为 $X$ 上的概形，
并假设下列条件之一成立：
\begin{enumerate}
\item $V \to X$ 射影；
\item $V \to X$ 拟射影；
\item $V$ 上存在充足可逆层；
\item $V$ 上存在 $X$-充足可逆层；
\item $V$ 上存在 $X$-很充足可逆层。
\end{enumerate}
则 $V/X/Y$ 上的任意下降数据都是有效的。
\end{lemma}

\begin{proof}
验证引理 \ref{groupoids-lemma-descend-along-finite} 中的条件。
设 $y \in Y$，且 $v_1, \ldots, v_d \in V$ 是 $y$ 上方的点。
情形 (1) 是 (2) 的特殊情形，见
《概形态射》引理 \ref{morphisms-lemma-projective-quasi-projective}。
情形 (2) 是 (4) 的特殊情形，见
《概形态射》定义 \ref{morphisms-definition-quasi-projective}。
若 $V$ 上存在充足可逆层，则由《概形的性质》引理 \ref{properties-lemma-ample-finite-set-in-affine}，
存在包含 $v_1, \ldots, v_d$ 的仿射开集。因此 (3) 成立。
在情形 (4) 与 (5) 中，可无妨把 $Y$ 替换为 $y$ 的仿射开邻域；
此时 $X$ 也仿射。在情形 (4) 中，由《概形态射》定义 \ref{morphisms-definition-relatively-ample}，$V$ 有充足可逆层，
故结论由情形 (3) 得出。在情形 (5) 中，可把 $V$
替换为包含 $v_1, \ldots, v_d$ 的拟紧开集，再由
《概形态射》引理 \ref{morphisms-lemma-ample-very-ample}
归约到情形 (4)。
\end{proof}

\begin{lemma}
\label{groupoids-lemma-descend-along-finite-radicial}
设 $X \to Y$ 为满的有限局部自由根态射。 
设 $V$ 为 $X$ 上的概形。则 $V/X/Y$ 上的任意下降数据都是有效的。
\end{lemma}

\begin{proof}
取 $d = 1$，应用引理 \ref{groupoids-lemma-descend-along-finite-with-bound}。
\end{proof}
